%=============================================================================
% UIDT Ontology v4.0 MSS-000 EXECUTOR CANDIDATE — Phase-6 Rebase:
% Relational-Algebraic Pregeometry, Open Continuum Passage, and Verification Ledger
% Author: Philipp Rietz (ORCID: 0009-0007-4307-1609)
% STATUS: MSS-000 EXECUTOR CANDIDATE under PI Decision D20; not applied; PI GO required.
% Basis manuscript: UIDT_Ontology_v3_9_9_DEWRAPPED-006 (DOI 10.5281/zenodo.20319634)
% Parent framework record: DOI 10.5281/zenodo.17835200
% License: CC BY 4.0
%
% MSS-000 EXECUTOR NOTE (non-rendering):
% This candidate is derived from PR-#15 blob 601cc57d under D20. It is not the
% applied manuscript, not a merge authorization, not a canonical freeze, not a
% proof of UIDT, and not an evidence-class upgrade. Claude Desktop audits the
% full diff before application; explicit PI GO gates the single-writer apply.
%
% FAIL-FAST ANCHORS (each occurs exactly once; patch scripts must abort if
% an anchor count differs from 1):
%   ANCHOR-P2-PART-I-CONTENT / ANCHOR-P3-PART-II-CONTENT /
%   ANCHOR-P4-PART-III-CONTENT / ANCHOR-P5-BIBLIOGRAPHY
%   (live form: double-percent delimited, one occurrence each)
%=============================================================================
%=============================================================================
% PROCESSING-COVERAGE LEDGER  (non-rendering; audit artifact)
% Updated: R3.1 tension-ledger pass, 2026-07-13 (T9 row + T1 copy-edit;
%          gated by O-1/O-1-followup/O-2/O-3, advisory, PI signature pending). Percentages are share of source lines actually
% READ (not skimmed, not summarised). "Used" = content entered the manuscript.
% This block exists so that a fresh session can see, without asking, what has
% and has not been consumed. Do not delete. Update at every phase close.
%
% --- PHASES ------------------------------------------------------------------
%  P1 skeleton .................. 100%  done
%  P2 Part I .................... 100%  done
%  P3 Part II ................... 100%  done
%  P4 Part III .................. 100%  done
%  P5 back matter + bib ......... 100%  done
%  R0 reading stage .............  60%  Modules B/C/E only partially read
%  R1 coda (philosophy+adversary) 100%  done
%  R2 Part I additions .......... 100%  done
%  R3 Part II additions .........  ~90% II-1, II-2, II-3, DIR-S-02/03,
%                                       observer taxonomy (Class 0-III, D20
%                                       Class-0 reframe: kernel-level, not the historical scalar model),
%                                       AGI underdetermination red team, and
%                                       consciousness-physics demarcation matrix
%                                       (D20 primitive-cell reframe) DRAFTED.
%                                       R3 prose complete; OPEN only: compile.
%                                       Compile + four-checker NOT run this pass
%                                       (must run in Desktop/Claude Code).
%  R4 MSS-000 architecture ....... 100%  executor candidate; audit pending
%  R5 post-apply reconciliation ..   0%  executor verification pending
%
% --- SOURCES -----------------------------------------------------------------
%  -006 canonical (3125 ln) ............... ~88% read | 23/26 sections relocated
%      OPEN: §412 empirical baseline; §1050 deeper gap/milestones;
%            §1411 psychedelics/NDE; §1748 architecture/synthesis
%            (§1321 observer taxonomy + §1463 demarcation matrix relocated R3)
%  Lean corpus @ 61bc24da (26 files) ...... 100% static audit | F1/F2/F3 open
%                                            (GitHub issue #10, advisory)
%  Module A  FOUNDATION      (2,046 ln) ...  ~5% structure map only
%  Module B  MATRIX THERMO  (29,505 ln) ...  ~4% targeted extraction only
%                                            (= raw chat UIDT-CHAT-002 x26)
%  Module C  COSMOLOGY/PBH   (2,026 ln) ... ~10% structure + PBH claim scan
%                                            (contains NO UIDT PBH model)
%  Module D  OPERATOR NORMS    (170 ln) ... 100% read | audit2 PATCH C refuted
%  Module E  MISC/LITERATURE(23,348 ln) ...  ~2% structure map only
%                                            (= raw chat UIDT-CHAT-001 x14)
%  01_KNOWLEDGE_GRAPH.json ................    0% unread
%  Ultimate_Pipeline/*.tex (1,995 ln) .....    0% unread (5 modules)
%  json_output/ (111 files) ............... ~5% (DR-010,018,024,026 only)
%  UNIVERSUM_SIM-master LEDGER/ docs/ ..... 0% unread
%
% --- BIBLIOGRAPHY ------------------------------------------------------------
%  99 bibitems | 59 cited | 0 dangling  (bibitem/cite counts grep-verified on
%    this file at R3 pass; parker2019 + rabinovici2021 now cited by II-3.
%    overfull/undefined-ref status pending pdflatex, not run this pass)
%  Verified in this session: cacic2011 (pp/venue), vansuijlekom2019 (authorship
%    corrected), cubitt2015, vanraamsdonk2010
%  NOT resolved -> may not be cited: Kocharovsky BEC papers; DR-006 (OCR-degraded)
%  Resolved this pass: cubitt2015, vanraamsdonk2010 (R2); connesrovelli1994 (R2);
%    parker2019, rabinovici2021 (post-R2, pre-R3)
%
% --- STANDING GATES ----------------------------------------------------------
%  G-DIFF   run at every phase close: map every -006 section to
%           relocated / superseded-with-reason / dropped-with-reason
%  G-SOURCE no manuscript block written from a summary of a source
%  G-CLASS  every externally supplied evidence tag re-derived from scratch
%
% --- REFUSED (do not reintroduce) --------------------------------------------
%  audit2 Gap 3.3  Markov blanket "forces" (3,2,1)      -> target hunting
%  audit2 Gap 10   container smash as [A] no-go theorem -> no theorem exists
%  audit2 PATCH C  fluctuation floor ~2a^2 tagged [A]   -> refuted by its source
%  Module B        "inner derivations lay the foundation for gauge groups"
%                                                       -> Prop. inner-derivation
%=============================================================================
\PassOptionsToPackage{hyphens}{url}
\PassOptionsToPackage{table,svgnames,dvipsnames}{xcolor}

\documentclass[12pt,a4paper,twoside]{article}

%=============================================================================
% PACKAGES (identical to -006 baseline)
%=============================================================================
\usepackage[utf8]{inputenc}
\usepackage[T1]{fontenc}
\usepackage{microtype}
\usepackage{fix-cm}
\usepackage{mathpazo}

\usepackage{geometry}
\usepackage{setspace}
\usepackage{titlesec}
\usepackage{fancyhdr}
\usepackage{booktabs}
\usepackage{longtable}
\usepackage{array}
\usepackage{multirow}
\usepackage{enumitem}
\usepackage{float}
\usepackage{caption}
\usepackage[numbers,sort&compress]{natbib}
\usepackage{etoolbox}
\usepackage{needspace}
\usepackage{placeins}

\usepackage{amsmath,amssymb,amsfonts,amsthm}
\usepackage{mathtools}
\usepackage{bm}
\usepackage{physics}

\usepackage{graphicx}
\usepackage{xcolor}

\usepackage{tikz}
\usetikzlibrary{shapes,arrows,positioning,shadows}
\usepackage[most]{tcolorbox}
\tcbuselibrary{skins,breakable}

\usepackage{listings}

\usepackage{hyperref}
\usepackage{cleveref}

%=============================================================================
% COLORS (monochrome master-class palette, unchanged from -006)
%=============================================================================
\definecolor{catA}{RGB}{0,0,0}
\definecolor{catAminus}{RGB}{0,0,0}
\definecolor{catB}{RGB}{0,0,0}
\definecolor{catC}{RGB}{0,0,0}
\definecolor{catD}{RGB}{0,0,0}
\definecolor{catE}{RGB}{0,0,0}
\definecolor{darkblue}{RGB}{0,0,0}
\definecolor{dimgray}{RGB}{105,105,105}
\definecolor{pioverride}{RGB}{0,0,0}
\definecolor{uidtBlack}{RGB}{0,0,0}
\definecolor{uidtRule}{gray}{0.22}
\definecolor{uidtHairline}{gray}{0.62}
\definecolor{uidtSoft}{gray}{0.985}
\definecolor{uidtTitleShade}{gray}{0.955}
\definecolor{uidtTitleRule}{gray}{0.12}
\providecolor{darkgray}{RGB}{64,64,64}

%=============================================================================
% PAGE SETUP (unchanged)
%=============================================================================
\geometry{a4paper, left=3cm, right=3cm, top=3cm, bottom=3cm}
\setlength{\headheight}{14.5pt}
\addtolength{\topmargin}{-2.5pt}
\clubpenalty=10000
\widowpenalty=10000
\raggedbottom
\setlength{\parskip}{0.65em plus 0.12em minus 0.04em}
\setlength{\parindent}{0pt}
\setstretch{1.32}

%=============================================================================
% SECTION HEADING STYLE (unchanged)
%=============================================================================
\titleformat{\section}
 {\normalfont\fontsize{16}{19}\selectfont\bfseries} {\thesection}{1em}{}
\titlespacing*{\section}{0pt}{3.0ex plus 1ex minus 0.2ex}{1.8ex plus 0.2ex}

\titleformat{\part}[display]
 {\normalfont\huge\bfseries\filcenter} {\partname\ \thepart}{0.6em}{\fontsize{20}{24}\selectfont}
\titlespacing*{\part}{0pt}{0pt}{2.5ex}

%=============================================================================
% HEADER/FOOTER (draft-status variant)
%=============================================================================
\fancypagestyle{uidtontology}{%
 \fancyhf{}
 \fancyhead[LE,RO]{\thepage}
 \fancyhead[RE]{\small\textit{UIDT v4.0 MSS-000: Phase-6 Rebase Candidate}}
 \fancyhead[LO]{\small\textit{P.\ Rietz}}
 \fancyfoot[C]{%
 \footnotesize\textcolor{darkgray}{%
 \copyright\ 2026 P.\ Rietz $\cdot$ CC BY 4.0 $\cdot$ MSS-000 candidate (D20; application pending) $\cdot$
 basis: \href{https://doi.org/10.5281/zenodo.20319634}{DOI 10.5281/zenodo.20319634} | \href{https://github.com/Mass-Gap/UIDT-Framework-v3.9-Canonical}{UIDT Canonical Repository}}%
 }
 \renewcommand{\headrulewidth}{0.4pt}
 \renewcommand{\footrulewidth}{0.4pt}
}

%=============================================================================
% METADATA
%=============================================================================
\hypersetup{
 pdftitle={Relational-Algebraic Pregeometry: UIDT v4.0 MSS-000 Phase-6 Rebase Candidate}, pdfauthor={Philipp Rietz}, pdfsubject={Falsification-first research programme for finite relational-algebraic pregeometry with open continuum passage and verification-ledger strata}, pdfkeywords={Relational Algebra, Noncommutative Geometry, Spectral Triple, Krajewski Diagram, Pregeometry, Verification Ledger, UIDT}, colorlinks=true, linkcolor=black, citecolor=black, urlcolor=black, pdflang={en}
}

%=============================================================================
% THEOREM ENVIRONMENTS (unchanged)
%=============================================================================
\BeforeBeginEnvironment{theorem}{\par\nopagebreak\needspace{4\baselineskip}}
\BeforeBeginEnvironment{definition}{\par\nopagebreak\needspace{4\baselineskip}}

\theoremstyle{plain}
\newtheorem{theorem}{Theorem}[section]
\newtheorem{proposition}[theorem]{Proposition}
\newtheorem{lemma}[theorem]{Lemma}
\newtheorem{corollary}[theorem]{Corollary}

\theoremstyle{definition}
\newtheorem{definition}[theorem]{Definition}
\newtheorem{axiom}{Axiom}
\newtheorem{prediction}{Testable Prediction}
\newtheorem{falsmark}{Falsification Criterion}

\theoremstyle{remark}
\newtheorem{remark}[theorem]{Remark}
\newtheorem{limitation}[theorem]{Limitation}
\newtheorem{openquestion}{Open Question}

\crefname{axiom}{axiom}{axioms}
\Crefname{axiom}{Axiom}{Axioms}
\crefname{prediction}{prediction}{predictions}
\Crefname{prediction}{Prediction}{Predictions}
\crefname{openquestion}{open question}{open questions}
\Crefname{openquestion}{Open Question}{Open Questions}
\crefname{limitation}{limitation}{limitations}
\Crefname{limitation}{Limitation}{Limitations}

% Custom Commands (carried over; Part III depends on these)
\newcommand{\UIDT}{\textsc{uidt}}
\newcommand{\GeV}{\,\mathrm{GeV}}
\newcommand{\MeV}{\,\mathrm{MeV}}
\newcommand{\catmark}[1]{\textsuperscript{\textsf{\textbf{[#1]}}}}
\newcommand{\LegacySx}{S_{\mathrm{legacy}}(x)} % museum-only historical notation; never an active primitive
\newcommand{\Sx}{\LegacySx} % compatibility alias used only by byte-preserved D19 governance records
\newcommand{\lambdaS}{5/12 \approx 0.4167} % [A] display alias for exact RG value 5\kappa^2/3.


% Phase-6 architecture guard: \LegacySx is permitted only for historical,
% negative-boundary, or museum references. It is not an active field macro.
%===================================================================
% EVIDENCE-DISCIPLINE MACROS (carried over verbatim; no new class invented)
%===================================================================
\newcommand{\Bdelta}{\hyperref[box:pi-override-delta]{\catmark{B}}}
\newcommand{\Bpending}{\catmark{B}\,\textit{(pending source audit)}}
%===================================================================
% UIDT PARAMETER LEDGER MACROS (v3.9 ledger values remain authoritative;
% RACE CONDITION LOCK: do NOT move mp.dps = 80 to config. Display aliases
% only; authoritative values live in LEDGER/CONSTANTS.md.)
%===================================================================
\newcommand{\DeltaGap}{1.710 \pm 0.015\GeV} % [B] PI Decision D18
\newcommand{\kappaVal}{0.500} % [A] dimensionless EFT-normalised coefficient
\newcommand{\GammaGeom}{16.339} % [A-] calibrated, NEVER derived
\newcommand{\GammaMC}{16.374 \pm 1.005} % [A-] Monte Carlo mean; tension entry vs 16.339
\newcommand{\vVEV}{47.7\MeV} % [A]
\newcommand{\ETorsion}{2.44\MeV} % [E pending] observer-context only
\newcommand{\LambdaGrid}{0.66\,\mathrm{nm}} % [D]
\newcommand{\KappaRG}{\tfrac{1}{2}\ \text{(exact)}} % [A] exact framework input
%===================================================================

% Lean 4 listing style for Part I (typographic only; no execution semantics)
\lstdefinelanguage{Lean4}{
  morekeywords={structure,class,instance,def,lemma,theorem,where,import,namespace,end,variable,Prop,Type,fun,let,by,intros,simp,rfl,open},
  sensitive=true,
  morecomment=[l]{--},
  morecomment=[s]{/-}{-/},
  morestring=[b]",
}
\lstset{
  language=Lean4,
  basicstyle=\small\ttfamily,
  commentstyle=\itshape\color{dimgray},
  keywordstyle=\bfseries,
  columns=fullflexible,
  keepspaces=true,
  breaklines=true,
  frame=none,
  xleftmargin=1.2em,
  aboveskip=0.9em,
  belowskip=0.9em,
}

%=============================================================================
% TCOLORBOX ENVIRONMENTS (unchanged from -006)
%=============================================================================
\tcbset{%
 uidtboxbase/.style={%
  enhanced, breakable, sharp corners, colback=uidtSoft, colframe=uidtRule, coltitle=uidtBlack, boxrule=0pt, arc=0pt, left=11pt, right=11pt, top=9pt, bottom=9pt, boxsep=0pt, before skip=1.15em plus 0.25em minus 0.15em, after skip=1.15em plus 0.25em minus 0.15em, fontupper=\small\setstretch{1.18}, parbox=false, borderline north={0.45pt}{0pt}{uidtHairline}, borderline south={0.45pt}{0pt}{uidtHairline}, },
 uidttitle/.style={%
  colbacktitle=white,
  boxed title style={%
   sharp corners, colback=white, colframe=white, boxrule=0pt, left=0pt, right=0pt, top=0pt, bottom=0pt, }, attach boxed title to top left={xshift=11pt,yshift=-1.8mm}, },
}

\newtcolorbox{evidencebox}[2][]{%
 uidtboxbase, uidttitle, title={#1},
 fonttitle=\scriptsize\bfseries\sffamily\MakeUppercase\lsstyle,
 top=17pt, borderline west={1.1pt}{0pt}{uidtTitleRule},
 overlay unbroken and first={%
 \node[anchor=north east, font=\tiny\sffamily\bfseries\lsstyle, text=uidtHairline] at ([xshift=-7pt,yshift=14pt]frame.north east) {EVIDENCE}; },
}

\newtcolorbox{killswitch}{%
 uidtboxbase, colframe=uidtBlack, colback=white,
 fonttitle=\scriptsize\bfseries\sffamily\MakeUppercase\lsstyle,
 title={Falsification Criterion}, boxrule=0pt, borderline north={0.65pt}{0pt}{uidtBlack}, borderline south={0.65pt}{0pt}{uidtBlack}, borderline west={1.8pt}{0pt}{uidtBlack},
}

\newtcolorbox{axiombox}[1][]{%
 uidtboxbase, uidttitle, colframe=uidtRule,
 fonttitle=\scriptsize\bfseries\sffamily\MakeUppercase\lsstyle,
 title={#1}, borderline west={0.9pt}{0pt}{uidtRule},
}

\newtcolorbox{probbox}[2][]{%
 uidtboxbase, uidttitle, colframe=uidtBlack, colback=white,
 fonttitle=\scriptsize\bfseries\sffamily\MakeUppercase\lsstyle,
 title={#1 \hfill #2}, boxrule=0pt, borderline north={0.55pt}{0pt}{uidtBlack}, borderline south={0.55pt}{0pt}{uidtBlack}, borderline west={1.4pt}{0pt}{uidtBlack},
}

\newtcolorbox{pioverridebox}[1][]{%
 uidtboxbase, uidttitle, colframe=uidtBlack, colback=white,
 fonttitle=\scriptsize\bfseries\sffamily\MakeUppercase\lsstyle,
 title={#1}, boxrule=0pt, borderline north={0.75pt}{0pt}{uidtBlack}, borderline south={0.75pt}{0pt}{uidtBlack}, borderline west={2.2pt}{0pt}{uidtBlack},
}

\newtcolorbox{governancebox}[1][]{%
 uidtboxbase, uidttitle, colframe=uidtHairline,
 fonttitle=\scriptsize\bfseries\sffamily\MakeUppercase\lsstyle,
 title={#1}, boxrule=0pt, borderline west={1.1pt}{0pt}{uidtRule},
}

%=============================================================================
\begin{document}
\pagestyle{uidtontology}

%=============================================================================
% TITLE
%=============================================================================
\begin{titlepage}
\vspace*{1.6cm}
\begin{center}
{\LARGE\bfseries Relational-Algebraic Pregeometry:\\[0.25em]
A Phase-6 Rebase of the UIDT\\[0.25em]
Research Programme}\\[0.9em]
{\small\bfseries A Falsification-First Programme\catmark{D}\ for a Finite Algebraic Kernel,\\[0.25em]
an Explicitly Open Passage to Apparent Continuity,\\[0.25em]
and a Non-Ontological Verification and Parameter Ledger}\\[1.5em]
{\large\itshape Unified Information-Density Theory (UIDT) v4.0\\[0.25em]
MSS-000 Executor Candidate under PI Decision D20}\\[2.6em]
{\small Philipp Rietz}\\[0.5em]
{\small ORCID: \href{https://orcid.org/0009-0007-4307-1609}{0009-0007-4307-1609}}\\[0.2em]
{\small Independent Researcher}\\[3em]
{\normalsize Successor draft to UIDT Ontology v3.9.9 (basis: -006)}\\[0.5em]
{\small Basis manuscript (Ontology v3.9):
\href{https://doi.org/10.5281/zenodo.20319634}{DOI 10.5281/zenodo.20319634}}\\[0.25em]
{\small UIDT Framework (parent record):
\href{https://doi.org/10.5281/zenodo.17835200}{DOI 10.5281/zenodo.17835200}}\\[0.5em]
{\small 19~July~2026 --- EXECUTOR CANDIDATE, not applied, not archived, not frozen}\\[2.2em]
\rule{0.6\textwidth}{0.4pt}\\[1em]
{\small\itshape ``Reality is not a collection of things,\\
but a process of structural individuation\\
that we are licensed to describe only as far as we can falsify it.''}\\[1.4em]
\end{center}
\end{titlepage}

%=============================================================================
% GOVERNANCE AND STATUS PAGE (renders before the table of contents)
%=============================================================================
\thispagestyle{uidtontology}

\begin{governancebox}[Governance Status of This Executor Candidate]
This document is the MSS-000 executor candidate derived from PR-\#15 blob
\texttt{601cc57d082840424bfa2aedc6542eb1b6d754e4}. It is not the applied
manuscript, not a canonical freeze, not a merge authorization, and not an
evidence-class upgrade. PI Decision D20 authorizes the architectural rebase
under the binding guardrails stated below; it does not authorize application.
Claude Desktop must audit the complete diff against the baseline blob and may
recommend GO or HOLD. Only explicit PI GO permits Claude, as the single
writer of the active \texttt{.tex}, to apply the approved patch. Build,
repository reconciliation, and technical post-application verification then
return to the executor lane. Final integration remains PI-only.
\end{governancebox}

\begin{pioverridebox}[PI Decision D19 --- Ontological Re-Architecture (DRAFT; signature pending)]
\label{box:pi-decision-d19}
\textbf{Decision under review.} The continuous real scalar field $\Sx$,
which UIDT v3.9.x treated as the ontological primitive (Variant~B), is in
this draft \emph{demoted} to a macroscopic effective-field-theory limit
(Part~III). The candidate primitive investigated in its place is a finite
algebraic structure --- a real spectral-triple kernel with Krajewski
combinatorics --- presented throughout Part~I strictly as a
\emph{research programme at evidence class \catmark{D}}, never as an
established foundation. Three binding conditions attach to this decision:
(i)~Part~I is declared a formal research programme\catmark{D}, not a
theorem corpus; (ii)~no Standard-Model-like block partition is hard-coded
into any core formal structure --- partitions enter only as inserted
candidate instances under the Gap-Localization-before-Construction (GLBC)
principle; (iii)~all v3.9 ledger results retain their existing evidence
classes unchanged inside Part~III. \textbf{This box is a governance
artifact, not a physics result. It takes effect only upon PI signature in
LEDGER/CLAIMS.json; until then, Variant~B remains canonical.}
\end{pioverridebox}

\begin{pioverridebox}[PI Decision D20 --- Part III Architectural Reassignment]
\label{box:pi-decision-d20}
\textbf{Authorized architectural change.} The quantitative content of the
former effective-scalar Part~III --- including $v$, $\lambda_S$, $\kappa$,
the EFT machinery, benchmark batteries, and cosmological calibration ledger
--- has no established derivation from the pregeometric kernel and therefore
does not function as active ontology. Under D20 it is reorganized as a
verification, benchmark, and parameter-ledger stratum. No quantitative result
is deleted or reclassified. The finite relational-algebraic kernel of Part~I
remains a falsifiable research programme at evidence class D; the passage of
Part~II remains open and supplies no continuum, metric, signature, Born-rule,
or Standard-Model gauge-group derivation.

\textbf{Excluded decisions.} D20 does not decide PI-DELTA-P6,
PI-ET-P6, PI-MUSEUM, or P2. Consequently $\Delta=\DeltaGap$ remains
\catmark{B} under D18 with its uncertainty and all locators intact; the
observer-context class of $E_T$ is unchanged while the subatomic-to-observer
firewall remains binding; and scalar-history material is retained reversibly
in a non-imported historical appendix rather than deleted.

\textbf{Application status.} Authorized for executor authorship and audit,
but not yet applied. Claude Desktop performs the pre-application audit; PI GO
or HOLD gates the single-writer application.
\end{pioverridebox}

\begin{evidencebox}[Reading Order and Epistemic Weight]{}
Part~I (\emph{The Pregeometric Kernel}) remains a \catmark{D}-class research
programme. It localizes obstructions, states conditional lower bounds, and
records finite-algebraic and Lean~4 structures without promoting an inserted
candidate to an established foundation. Part~II (\emph{The Open Passage to
Apparent Continuity}) is capped at \catmark{E} wherever it introduces
coarse-graining, observer functors, thermal-time interpretation, or apparent
continuity; it is an obligation register, not a supplied mechanism. Part~III
(\emph{Verification, Benchmarks, and Parameter Ledger}) retains every
quantitative result at its incoming class: $\kappa=\KappaRG$\,\catmark{A},
$\lambda_S=\lambdaS$\,\catmark{A}, $v=\vVEV$\,\catmark{A},
$\gamma=\GammaGeom$\,\catmark{A-} (calibrated, never derived), and
$\Delta=\DeltaGap$\,\catmark{B} under D18 and Limitation~L12. Cosmological
statements remain capped at \catmark{C}. The historical scalar/EFT model is
preserved in Appendix~\ref{app:scalar-museum} as a non-imported museum record;
its equations are not premises of Parts~I--III.
\end{evidencebox}

\newpage
\tableofcontents
\newpage

%=============================================================================
% ABSTRACT (deflationary, audit-reconciled)
%=============================================================================
\begin{abstract}
\begingroup \normalfont\normalsize
\setlength{\parskip}{0.65em plus 0.12em minus 0.04em}%
\setlength{\parindent}{0pt}%
\noindent This MSS-000 candidate rebases the UIDT v4.0 working manuscript
onto a Phase-6 relational-algebraic architecture. The active foundational
programme contains no classical scalar field, background spacetime, metric,
or smooth-manifold primitive. Its candidate kernel is expressed through
finite combinatorial and algebraic data --- in particular the separation
\texttt{BlockPartition $\to$ FiniteAlgebraSignature $\to$ SpectralTriple} ---
and remains a falsifiable research programme at evidence class \catmark{D}.
The known $d^2=0$ scalar-gradient obstruction, the unresolved origin of
$G_{\mathrm{SM}}$, the absence of a Born-rule derivation, and the absence of
solutions to the $H_0$ and $S_8$ tensions remain explicit boundaries.

Part~II records the open passage from the finite kernel to apparent
continuity. It supplies no dimensional-emergence theorem, Lorentzian
signature, background metric, or non-circular observer derivation and is
capped at \catmark{E} wherever it becomes interpretive. Under PI Decision
D20, Part~III no longer functions as active ontology. It is a verification,
benchmark, and parameter-ledger stratum that preserves the incoming
quantitative content and evidence classes without deletion or upgrade:
$\kappa=\KappaRG$\,\catmark{A}, $\lambda_S=\lambdaS$\,\catmark{A},
$v=\vVEV$\,\catmark{A}, $\gamma=\GammaGeom$\,\catmark{A-} as a permanent
calibration, and $\Delta=\DeltaGap$\,\catmark{B} under D18 and Limitation~L12.
Cosmological entries remain calibrations capped at \catmark{C}.

The former scalar/EFT architecture is retained intact enough for audit in a
non-imported historical appendix. It is not silently deleted, and its
Lagrangian, coordinate derivatives, background metric ansatz, and scalar
correlator programme are not premises of the active architecture. This
candidate is not applied and carries no merge authority: Claude Desktop must
audit the complete diff against blob \texttt{601cc57d}; explicit PI GO or
HOLD gates application.
\endgroup
\end{abstract}

\newpage

%=============================================================================
% CHAPTER 0 — EVIDENCE APPARATUS
%=============================================================================
\section{Evidence Grammar, Strata, and Audit Discipline}
\label{sec:evidence-apparatus}

\noindent The order of this manuscript is deliberate. Before the
pregeometric kernel is defined, before a single formal structure is
introduced, this chapter fixes the discipline by which every later
assertion will be judged. A re-architecture is precisely the moment at
which a framework is most tempted to convert enthusiasm into evidence; the
apparatus below exists to make that conversion impossible.

\subsection{Evidence Classes}
\label{subsec:evidence-classes}

Every nontrivial claim in this draft carries exactly one of six evidence
classes, rendered as a typographic tag. \catmark{A} marks mathematical
proof closure, an exact identity, or a verified deterministic theorem
under stated assumptions; internal residual closure below $10^{-14}$ at
80-digit \texttt{mpmath} precision qualifies only for deterministic
arithmetic and is never external physical agreement. \catmark{A-} marks a
calibrated internal closure: calibrated, not derived. \catmark{B} marks
lattice-compatible or externally benchmarked physics with verified sources
and explicit uncertainty handling. \catmark{C} is the cosmology cap; no
cosmological statement in this draft exceeds it, and no statement at
\catmark{C} claims to resolve any observational tension. \catmark{D} marks
a prediction or research programme: falsifiable but unconfirmed.
\catmark{E} marks speculative, interpretive, withdrawn, or museum content.
No other class exists; composite or inflected tags (such as A+ or B$-$)
are not part of the grammar and their appearance anywhere in this document
would itself constitute an audit failure.

\subsection{Strata}
\label{subsec:strata}

Claims are additionally stratified by provenance. Stratum~I comprises
empirical measurements and externally sourced values with uncertainties.
Stratum~II comprises consensus theory: standard quantum field theory,
lattice methodology, the Standard Model, $\Lambda$CDM formalism, and ---
central to Part~I --- the established mathematics of finite real spectral
triples and Krajewski classification. Stratum~III comprises \UIDT\
interpretation: the pregeometric-kernel hypothesis, the rejection of the historical scalar ansatz as active ontology,
every mapping between formal structure and physical ontology, and all
observer-theoretic material. The cardinal rule is that Stratum~III is
never written as Stratum~I, and that the relevance of a Stratum~II
mechanism is never presented as validation of a Stratum~III
interpretation.

\subsection{Falsification-First Rule}
\label{subsec:falsification-first}

Every nontrivial scientific, mathematical, or ontological claim in this
draft carries at least one of: an empirical falsification condition, a
mathematical proof or verification condition, a numerical reproduction
condition, or an explicit downgrade condition. Where no such condition can
be stated, the claim is downgraded to \catmark{D} or \catmark{E} on the
spot. A claim that cannot fail is not treated as a physics result.

\subsection{The Gap-Localization-before-Construction Principle}
\label{subsec:glbc}

Part~I is governed by a methodological rule inherited from the v3.9 audit
cycle and here promoted to a structural constraint: \emph{Gap Localization
before Construction} (GLBC)\catmark{E}\ (methodology-class content; the
tag records that GLBC is a working discipline, not a theorem). Formal
types are defined empty before any physical scale, any ledger constant, or
any Standard-Model-like pattern is applied to them. In particular, the
calibrated invariant $\gamma = \GammaGeom$\,\catmark{A-} and the spectral
gap $\Delta = \DeltaGap$\,\catmark{B} never appear inside any core formal
structure of Part~I --- not as targets, not as losses, not as initial
conditions, not as stopping rules --- and the Standard-Model-like block
partition enters exclusively as an inserted candidate instance whose
endogenous selection is recorded as open. Any violation of this rule
constitutes target leakage and voids the affected construction.

\subsection{Numerical Discipline}
\label{subsec:numerical-discipline}

Where proof-critical arithmetic occurs, it is executed in Python~3 with
\texttt{mpmath} at \texttt{mp.dps = 80}, set locally inside the
proof-critical function rather than globally at module scope; binary
floating-point heuristics, \texttt{float()} coercion, and precision-losing
conversions are barred from proof-critical paths. High-precision output is
rendered via \texttt{mp.nstr(value, 80)}. No new 80-digit values are
invented or estimated in this draft; numerical statements either reproduce
ledger values at their recorded classes or are absent.

\subsection{Provenance of This Draft}
\label{subsec:provenance}

This draft consolidates: the canonical v3.9.9 manuscript
(\texttt{-006}, four-checker verified); the corrected G1--G4
admissibility lattice, with the block-repetition misreading of the
multiplicity condition explicitly quarantined as a documented category
error; the Phase-7 Lean~4 axiom-initialization architecture
(\texttt{RealStructure}, star-compatible representations, and an abstract
first-order-condition schema); the honest null result of the matrix-%
condensation metastability protocol; and the division-algebra programme
recorded as a \catmark{D} candidate for the intermediate algebra. Sources
that failed primary-source verification during the audit cycle ---
including an external audit document that rejected the
Standard-Model-like partition on a mathematically false multiplicity
argument --- are excluded from evidentiary use and cited, where cited at
all, only as documented failure history.

%=============================================================================
% PART I — THE PREGEOMETRIC KERNEL
%=============================================================================
\newpage
\part{The Pregeometric Kernel\texorpdfstring{\ \catmark{D}}{[D]}}
\label{part:pregeometric-kernel}

\noindent Part~I develops the candidate primitive of the v4.0
re-architecture: a finite algebraic structure of real spectral-triple type,
formalized in a typed proof assistant under the GLBC principle. Everything
in this part is a research programme at evidence class \catmark{D} unless
a narrower tag is stated. The part establishes, in order: the primitive
data and their Lean~4 typing; the Krajewski combinatorial frame; the
G1--G4 obstruction set with the corrected admissibility lattice; the lower
bounds any admissible intermediate algebra must satisfy; and the open
selection problems --- Wedderburn moduli, real structure, KO-dimension,
and chirality --- that define the programme's falsifiable milestones.
There is no space in this part, and no time: both are claims Part~II must
earn, not assumptions Part~I is permitted.

%%ANCHOR-P2-PART-I-CONTENT%%

%============================ P2 PART I CONTENT ============================
% Inserted at ANCHOR-P2-PART-I-CONTENT by fail-fast patch script.
% All Lean listings are ASCII transliterations of the verbatim repository
% sources at HEAD 61bc24da (Unicode tokens transliterated for typesetting);
% the authoritative text is the repository file.
%===========================================================================

\section{The Demotion of the Continuous Scalar and the Kernel Primitive}
\label{sec:demotion}

\subsection{The Negative Results That Motivate the Re-Architecture}
\label{subsec:negative-results}

The v3.9 audit cycle produced three independent negative results whose
joint force motivates the present re-architecture. None of them is new to
this draft; each is carried over at its established evidence class, and the
re-architecture adds no evidential weight to any of them.

First, the scalar-gradient route to gauge curvature fails
exactly\catmark{A}. If a gauge potential is built as $A = \dd S$ from a
smooth scalar on a manifold, then $F = \dd A = \dd^2 S = 0$ identically in
ordinary exterior calculus. This is a clean negative result for that
construction --- and only for that construction; it is not a universal
no-go theorem for enriched algebraic routes.

Second, the $G_{\mathrm{SM}}$-origin gap is localized, not closed. Within
established theorems, no gapless ab initio derivation
a structureless real scalar $\to G_{\mathrm{SM}}$ exists; the obstruction is concentrated in four
jointly unsatisfied conditions (G1--G4) on fiber multiplicity, Wedderburn
invariants, real structure, and chirality
(\cref{sec:g1g4}). On every known route, the Standard-Model gauge group is
selected, not forced.

Third, the dynamical-selection programme returned an honest null. The
blinded, preregistered metastability pre-protocol for matrix condensation
(\texttt{PREREG-PR-B1-PRE-001}, PI-signed 2026-06-17; detector, coupling
window $\tilde\alpha \in \{0.40, 0.55, 0.625, 0.75, 0.90\}$, and matrix
sizes $N \in \{16, 24, 32\}$ frozen before any result was seen) found
that a planted Standard-Model-like block configuration decays out of its
class --- in the detector's class-membership sense --- within
approximately ten Monte-Carlo steps of thermalization at every tested
coupling, confirming the preregistered literature null (single-sphere
vacuum) with an independently validated instrument.\catmark{D} The
protocol, its frozen outcome table, and the recorded outcome
(\texttt{PRE-O1}) are archived at \texttt{verification/prereg/PR-B1/}
with the corresponding ledger entry in \texttt{LEDGER/CLAIMS.json}. No
verified mechanism forcing a stable asymmetric block vacuum is known;
that search remains an open \catmark{D}/\catmark{E} programme.

A primitive that cannot generate the gauge sector, cannot curve anything
through its own gradient, and is not dynamically selected by the one
condensation mechanism tested, retains its primitive status only by
stipulation. The re-architecture withdraws the stipulation.

\begin{pioverridebox}[What the Demotion Does and Does Not Claim]
The removal of the historical scalar ansatz from active ontology under PI Decision D20 (authorized architecture; application pending;
\hyperref[box:pi-decision-d20]{governance page}) is an ontological reassignment, not a physics
result. It does \emph{not} assert that a discrete pregeometric kernel is
established; it asserts that the continuous scalar no longer carries the
burden of primitivity, and that the candidate now under investigation is
the finite algebraic structure of Part~I, at evidence class \catmark{D}
throughout. Replacing one unproven Stratum-III primitive by another
unproven Stratum-III primitive yields no upgrade anywhere; what it yields
is a candidate whose obstructions are sharper, whose lower bounds are
theorem-shaped, and whose formalization is machine-checkable in principle.
\end{pioverridebox}

\subsection{Design Directive DIR-S-01: The Primitive as a Pregeometric Operator}
\label{subsec:dir-s-01}

The formal seat of the demotion is Design Directive DIR-S-01, encoded in
the formalization layer as a foundation module. The primitive of the v4.0
programme is a pregeometric operator $\mathbf S$, not a classical field
$S \colon M \to \mathbb R$ on a pre-given four-dimensional spacetime.
Coordinates and metric are required to emerge at a later ontological
level; historical routes through $A = \dd S$ on a smooth manifold are excluded at
the type level, in direct response to the $d^2 = 0$ obstruction.

\begin{lstlisting}
/-- [D/E] Pre-geometric UIDT Primitive Operator S.
    DIR-S-01: S is an operator in a suitable C*-algebra /
    on a Hilbert space, not a classical field
    S : spacetime -> R. The concrete realization (Matrix,
    NCG, Tensor-Network) is an additional instance. -/
class PrimitiveOperator (S : Type*) where
  isPrimitive : Prop := True
\end{lstlisting}

Three features of this definition deserve emphasis, because each is a
deliberate act of restraint. There is no spacetime dependence: the module
does not define $S \colon \mathrm{Spacetime} \to \mathbb R$, and no metric
tensor is available at this level. The class is intentionally minimal:
algebraic axioms (self-adjointness, positivity, spectral properties) are
deferred until the formalization matures, rather than asserted
decoratively. And the $d^2=0$ obstruction is recorded in the module as a
negative design commitment --- the smooth-field route is rejected
entirely, not formalized and then argued around.

\begin{killswitch}
The kernel-primitive hypothesis is downgraded from \catmark{D} to
\catmark{E} if any of the following occurs: (i)~a theorem shows that no
finite real spectral triple satisfying the Part-I lower bounds
(\cref{subsec:lower-bounds}) can couple consistently to the Part-III
effective sector; (ii)~the conditional selection structure of
\cref{sec:selection} is shown to be empty or inconsistent once H1 is
derived from intersection-form axioms; (iii)~the formalization programme
is abandoned with its axiom layer unclosed. Conversely, no success of
Part~III arithmetic upgrades the kernel hypothesis.
\end{killswitch}

\section{The Lean~4 Formalization Architecture}
\label{sec:lean-architecture}

\subsection{Scope, Source of Truth, and Audit Boundary}
\label{subsec:lean-scope}

{\sloppy The formalization lives in the working repository
\path{UIDT-Framework-V3.9-UNIVERSUM_SIM} (module namespace
\texttt{Antigravit2}), audited for this draft at commit
\texttt{61bc24da} (2026-07-08, merge of the phase-10b axiom pull
request). The corpus at that head comprises 26 Lean files across five
layers: \texttt{Foundation} (the primitive operator),\par}
\texttt{MatrixThermo} (block partitions), \texttt{Filters} (enumeration
and admissibility predicates), \texttt{NCG} (real structure, first-order
condition, Krajewski combinatorics, spectral-triple assembly), and
\texttt{Meta} (evidence tags and the anti-target-leakage policy).

Two boundaries are binding for every claim in this section. First, the
audit underlying this draft is a \emph{static structural audit} of the
repository snapshot at the stated head; no \texttt{lake build} was
executed in the audit session, and no build result is asserted here.
Second, a successful build would in any case certify only type-checking;
it would not discharge the open proof obligations catalogued in
\cref{subsec:formalization-status}, and it would never convert a formal
result into a physics claim. Lean-proven statements are \catmark{A}
strictly \emph{within the formal system}; their physical import is bounded
by the epistemic status of the hypotheses they consume.

\subsection{Gap Localization before Construction, Operationalized}
\label{subsec:glbc-operational}

The GLBC principle of \cref{subsec:glbc} is not a stylistic preference in
this corpus; it is implemented as type-system structure. The gap parameter
module defines the spectral gap purely formally and blocks its premature
physical identification:

\begin{lstlisting}
/-- [D] Formal spectral gap parameter.
    No physical value (like 1.710) is assigned here. -/
structure SpectralGapParam where
  delta : Real
  delta_pos : 0 < delta

/-- [D] Explicit blocker for premature physical
    identification. The glueball mass conjecture is stated
    as False in the formal core, preventing any theorem
    from accidentally relying on Delta = M(0++). -/
def DeltaGlueballConjecture (Delta : Real) : Prop := False
\end{lstlisting}

The second definition deserves a sentence of appreciation and a sentence
of precision. It renders Limitation~L12 as a type-level fact: any formal
development that attempted to consume the glueball identification would
have to prove \texttt{False}, which is to say it could not proceed. The
precision: this blocks the identification \emph{inside the formal corpus};
it neither strengthens nor weakens the manuscript-level status of L12,
which remains a Stratum-II/III discipline boundary. Likewise, the ledger
constants $\gamma = \GammaGeom$\,\catmark{A-} and
$\Delta = \DeltaGap$\,\catmark{B} appear in the corpus only inside warning
comments; a literal scan of the head confirms no occurrence in any
definition, hypothesis, loss, or initial condition.

\subsection{The Real-Structure Layer (Phase 7)}
\label{subsec:real-structure}

The Phase-7 axiom layer follows the reviewed specification exactly:
anti-linearity is carried by a dedicated structure, conjugation is
expressed through the field's star ring endomorphism, and the operator $J$
is deliberately \emph{not} equipped with invertibility or unitarity ---
$J^{-1}$ is excluded from the phase in every form, including symbolically.

\begin{lstlisting}
/-- [D] Anti-linear map over field C.
    J(a . x) = conj(a) . J(x). -/
structure AntiLinearMap (H : Type _)
    [AddCommGroup H] [Module Complex H] where
  toFun    : H -> H
  map_add  : forall x y,
    toFun (x + y) = toFun x + toFun y
  map_smul : forall (a : Complex) x,
    toFun (a . x) = (starRingEnd Complex a) . toFun x

/-- [D] RealStructure: J and the KO sign data.
    J is anti-linear but NOT required to be invertible
    or unitary here. J-inverse is excluded from Phase 7. -/
class RealStructure (H : Type _)
    [AddCommGroup H] [Module Complex H] where
  koDimension : ZMod 8
  J    : AntiLinearMap H
  eps  : Int
  epsD : Int
  epsG : Int
  J_involutive : forall x, J (J x) = (eps : Complex) . x
\end{lstlisting}

A trivial instance on $\mathbb C$ --- $J$ realized as complex conjugation
with all signs $+1$ --- closes the regression layer without \texttt{sorry}
and certifies that the architecture is self-consistent on the smallest
non-degenerate carrier. One deviation from the drafting-session
specification is recorded for traceability: the head keeps only the
involutivity relation $J^2 = \varepsilon\,\mathrm{id}$ inside the class,
deferring the $JD$ and $J\gamma$ commutation relations, while adding a
\texttt{koDimension} field of type $\mathbb Z/8$. This is a conservative
narrowing --- fewer axioms asserted, none invented --- and is consistent
with the phase philosophy.

\subsection{The First-Order Condition as a Commutator Prop.\ (Phase 10b)}
\label{subsec:foc}

The first-order condition has graduated from the Phase-7 envelope to a
genuine structural proposition. The head defines a minimal spectral-triple
core --- Dirac operator, real-structure carrier, grading, and separated
left and right representations, the latter through an opposite-algebra
wrapper --- and states the condition in its standard commutator form:

\begin{lstlisting}
/-- Commutator abstraction for endomorphisms. -/
def Commute {H : Type _} (F G : H -> H) : Prop :=
  forall x, F (G x) = G (F x)

/-- First-Order Condition (FOC):
    [[D, a], b^o] = 0 for all a in A, b^o in A^o. -/
def firstOrderCondition {A H : Type _} [AddGroup H]
    (ST : SpectralTripleCore A H) : Prop :=
  forall (a : A) (b_op : Opposite A),
    Commute
      (fun x => ST.D (ST.leftRep a x)
                - ST.leftRep a (ST.D x))
      (ST.rightRep b_op)
\end{lstlisting}

This is the honest formal shape of the NCG first-order axiom at the
current stage: a proposition that can be stated and instantiated, with the
analytic content (boundedness, adjointability, the bimodule structure of
the right action) deferred to later phases. Nothing here asserts that any
physically interesting triple satisfies the condition.

\subsection{The Krajewski Combinatorial Kernel and the Inserted Candidate}
\label{subsec:krajewski}

The combinatorial layer is deliberately stripped of physical labels. An
edge of a Krajewski diagram is a bare record of two block indices and a
multiplicity; chirality-capability appears only as the derived structural
property that an edge connects distinct blocks; and diagram admissibility
at this level demands nothing beyond the absence of self-loops.

\begin{lstlisting}
structure Bimodule where
  leftBlock    : Nat
  rightBlock   : Nat
  multiplicity : Nat

def Bimodule.asymmetric (b : Bimodule) : Prop :=
  b.leftBlock /= b.rightBlock

structure KrajewskiDiagram where
  edges : List Bimodule
\end{lstlisting}

Against this neutral kernel, the Standard-Model-like configuration enters
exactly once, and exactly as governance condition~(ii) of D20 requires: as
an \emph{inserted candidate instance}, in a dedicated instance module,
decoupled from the filter hypotheses and from every physical constant. The
instance fixes the block list $[3,2,1]$ and the three off-diagonal edges,
and proves --- unconditionally, because the property is purely structural
--- that the resulting diagram has no self-loops. That theorem is
\catmark{A} within the formal system and physically weightless: it
certifies combinatorial well-formedness of an inserted diagram, nothing
more.

\begin{remark}[{Notation: $(1,2,3)$, $(2,3)$, and $[3,2,1]$}]
\label{rem:partition-notation}
The finite algebra $\mathbb C \oplus \mathbb H \oplus M_3(\mathbb C)$ has
matrix dimensions $(1,2,3)$ with total carrier dimension $N = 6$. The
v3.9 appendix literature labels this pattern by its non-abelian pair as
``the $(2,3)$ partition''; the Lean corpus writes the same object as the
descending list $[3,2,1]$. All three notations name one candidate
configuration, and this draft uses the Lean form wherever the formal layer
is in view.
\end{remark}

\subsection{Evidence Tags and Anti-Target-Leakage in Code}
\label{subsec:code-evidence}

The corpus formalizes its own audit discipline. A meta-module defines the
six evidence grades as an inductive type for structured documentation, and
states the anti-target-leakage policy as a first-class object: no proof
may contain its own target as input; the candidate partition must not
appear in any definition, axiom, or hypothesis, and may appear only in
theorem conclusions and test assertions, enforceable by review and by
grep-based continuous-integration checks. The policy also fixes the
correct reading of every Lean tag in this Part: all definitions are
\catmark{D}/\catmark{E} unless a proof is supplied, in which case the
proven statement is \catmark{A} within the formal system and is not
thereby a physics claim.

\subsection{Formalization Status at Head \texttt{61bc24da} and Audit Findings}
\label{subsec:formalization-status}

The table below records the verified status of the corpus at the audited
head. ``Closed'' means no \texttt{sorry} in the file at head; counts are
from a literal whole-word scan of the snapshot.

\begin{table}[H]
\centering\small
\begin{tabular}{@{}llp{5.2cm}@{}}
\toprule
Module & Proof state at head & Character \\
\midrule
\texttt{Foundation/PrimitiveOperator} & closed & type/class layer, DIR-S-01 \\
\texttt{MatrixThermo/BlockPartition} & closed & partitions; 14 decidable regressions \\
\texttt{NCG/RealStructure} & closed & Phase-7 axiom layer, trivial instance \\
\texttt{NCG/FirstOrderCondition} & closed & FOC as commutator proposition \\
\texttt{NCG/KrajewskiCore} & closed & combinatorial kernel \\
\texttt{NCG/Krajewski321} & closed & inserted candidate instance \\
\texttt{NCG/DeltaScale} & closed & gap parameter, L12 type blocker \\
\texttt{NCG/SpectralTriple} & 5 open & canonical-form regressions (registered) \\
\texttt{NCG/AxiomAudit} & 1 open & unregistered \\
\texttt{Filters/Enumeration} & 11 open & normal-form obligations (registered class) \\
\texttt{Filters/EliminationN6} & 34 open & includes the flagship proposition \\
\bottomrule
\end{tabular}
\caption{Lean corpus proof state at commit \texttt{61bc24da}; 51
substantive open obligations in total. Static audit; no build executed.}
\label{tab:lean-status}
\end{table}

\begin{probbox}[Audit Finding F1]{Register--Code Contradiction}
The corpus's own proof-obligation register asserts that the flagship
selection proposition (\cref{sec:selection}) is closed by native
decision-procedure evaluation. At the audited head this is false: the
theorem body is \texttt{sorry}, as are its thirty-three per-partition
verdict examples. Until the discrepancy is resolved on the remote, the
proposition's ``formally verified'' tag inside the corpus must be read as
aspirational, and this manuscript cites it strictly as \emph{stated, not
machine-checked at head \texttt{61bc24da}}. The finding is a
documentation-integrity defect, not a mathematical one; the intended
decidable proof route is plausible and small.
\end{probbox}

\begin{probbox}[Audit Finding F2]{Target-Shaped Admissibility Definition}
A wrapper in the candidate-instance layer defines a carrier to be
admissible if and only if its KO-dimension is $6$ \emph{and its block list
equals} $[3,2,1]$, and then proves the inserted candidate admissible by
reflexivity. This is definitional target-matching: the theorem is vacuous,
and the definition violates the corpus's own anti-target-leakage policy,
which bars the candidate partition from definitions. Two Prop-valued
signature fields are additionally discharged with the trivial proposition
rather than with proofs of the intended properties. The repair is
mechanical --- define admissibility through the generic filter predicates
and derive the candidate's status --- and is queued in the fix handout of
Phase~P5.
\end{probbox}

\begin{probbox}[Audit Finding F3]{Under-Registered Proof Debt}
Of the 51 open obligations at head, the register covers only the five
canonical-form regressions and the enumeration class; the 34 obligations
of the elimination module and one audit-module obligation are
unregistered. The corpus's own rule --- no \texttt{sorry} survives
silently --- is therefore violated at head. Register synchronization is
queued alongside F1/F2.
\end{probbox}

\section{The G1--G4 Obstruction Set and the Corrected Admissibility Lattice}
\label{sec:g1g4}

\subsection{The Four Conditions}
\label{subsec:g1g4-statement}

For the strict-scalar primitive with symmetries implemented by canonical
linear transformations, the non-derivability of $G_{\mathrm{SM}}$ is
localized to a single step --- the enlargement of the abelian single-mode
local algebra to a noncommutative algebra of local multiplicity at least
two~\cite{krajewski1998,paschkesitarz1998} --- and quantified by four conditions that no known mechanism
satisfies jointly. \textbf{(G1)} Fiber-multiplicity generation
$1 \to n \geq 2$ without inserted variables, blocked homologically (loop
phases factor through $H_1(\Gamma) = \pi_1(\Gamma)^{\mathrm{ab}}$, bracket
identically zero) and algebraically (a scalar-generated algebra is
commutative, with trivial inner automorphisms). \textbf{(G2)}
Theorem-determined Wedderburn invariants: at present the matrix dimension,
block partition, and real forms are free moduli~\cite{cacic2011}, and gauge structure is
the automorphism shadow of an inserted algebra (Skolem--Noether).
\textbf{(G3)} An endogenous real structure $J$ of KO-dimension~6: at
present an external axiom choice. \textbf{(G4)} Chirality-capable matter:
blocked by Nielsen--Ninomiya~\cite{nielsen1981} for local translation-invariant lattices,
with no established non-abelian chiral construction carrying
Standard-Model content. One partial escape deserves explicit statement, because it is the only
established endogenous enhancement of a single scalar to a non-abelian
current algebra: the Frenkel--Kac construction, which realises
$\widehat{\mathfrak{su}}(2)_1$ from a compactified free
boson~\cite{frenkelkac1980,goddardolive1986}. It satisfies a weakened
(G1) \emph{only} in $1{+}1$ dimensions, under target compactification
and radius tuning, and is bounded by a rank-budget argument: a level-one
simply-laced affine algebra so obtained supplies $\mathrm{rank} = 1$
against the requirement $\mathrm{rank}(\mathfrak g_{\mathrm{SM}}) = 4$.
This is the origin of the rank bound of
\cref{subsec:lower-bounds}.\catmark{A}

The consequence stands as in v3.9: on every known
route, $G_{\mathrm{SM}}$ is selected, not forced; this is a gap
localization, not an impossibility theorem for all conceivable enriched
routes.

\subsection{The Multiplicity Precision Remark and a Quarantined Category Error}
\label{subsec:multiplicity-precision}

\begin{remark}[Binding precision on ``multiplicity $\geq 2$'']
\label{rem:multiplicity}
The G1 condition concerns the \emph{non-abelian fiber}: at least one
algebra summand of matrix dimension $n_i \geq 2$, so that the inner
automorphism group is non-abelian. It is \emph{not} the requirement that
block dimensions repeat, and not a pairing constraint of $J$ on blocks.
The real structure constrains the \emph{bimodule multiplicity matrix}
$V_{ij}$~\cite{krajewski1998,paschkesitarz1998,cacic2011} --- each representation must be accompanied by its conjugate ---
and Wigner's anti-unitary argument guarantees that such a $J$ exists
already at $\dim V_{ij} = 1$. The Standard-Model algebra
$\mathbb C \oplus \mathbb H \oplus M_3(\mathbb C)$, each summand appearing
once, carries a KO-dimension-6 real structure~\cite{ccm2007,vansuijlekom2019}. The $[3,2,1]$ pattern
therefore \emph{satisfies} G1.
\end{remark}

An external audit document circulated during the v3.9 cycle rejected the
Standard-Model-like partition under G1 on the ground that ``each block
size occurs only once, and $J$ requires paired basis vectors.'' That
argument conflates block repetition with bimodule multiplicity and is
mathematically false by the result just cited; the same document
consequently admitted the all-equal partition $[2,2,2]$ and rejected
$[3,2,1]$ --- inverting the corrected lattice on both entries. The
document is quarantined at \catmark{E} as documented failure history and
is not used evidentially anywhere in this draft.

\subsection{The Corrected Per-Partition Lattice}
\label{subsec:corrected-lattice}

\begin{table}[H]
\centering\footnotesize
\begin{tabular}{@{}lllp{5.1cm}@{}}
\toprule
Partition & G1 & Verdict & Operative mechanism \\
\midrule
$[3,2,1]$ ($\mathbb C\oplus\mathbb H\oplus M_3(\mathbb C)$) & passes &
admissible\,\catmark{B} & real structure exists at $\dim V_{ij}=1$;
anomaly freedom via the standard hypercharge assignment \\
Pati--Salam-type & passes & mechanism correct\,\catmark{B} & admissible
without the order-one condition, per the established NCG route~\cite{ccvs2013,ccm2015} \\
$[6,3]$-type (unequal pair) & passes & rejected\,\catmark{B} &
intersection-form/Poincar\'e-duality bound $k=\ell\pm1$~\cite{iss2004}; the earlier
``order-one plus anomaly'' rationale was incorrect \\
$[2,2,2]$ & passes & rejected\,\catmark{D} & mass-degeneracy argument;
plausible but not verified against a primary-source theorem --- the
weakest verdict in the lattice \\
\bottomrule
\end{tabular}
\caption{Corrected G1--G4 admissibility verdicts. Verdict tags grade the
\emph{justification}, not the desirability, of each entry; the
$[2,2,2]$ rejection is the open verification item.}
\label{tab:lattice}
\end{table}

\subsection{Lower Bounds on Any Admissible Intermediate Algebra}
\label{subsec:lower-bounds}

Any candidate intermediate algebra $\mathcal A$ in a route
$\mathbf S \to \mathcal A \to G_{\mathrm{SM}}$ must carry: local
multiplicity at least two (a non-abelian fiber); rank resources at least
four, against $\mathrm{rank}(\mathfrak g_{\mathrm{SM}}) = 4$; a real
structure with the $J^2$-sign pattern of KO-dimension~6; and a chirality
mechanism evading the Nielsen--Ninomiya and finite-2-group obstructions.
Any candidate failing one bound is excluded before dynamical analysis.
These bounds are the constructive yield of the gap analysis: they tell a
model-builder what $\mathcal A$ must minimally contain, and they are the
standing acceptance test for every candidate discussed in
\cref{sec:candidates}.

\section{Filters H1/H2 and the Conditional $N=6$ Selection Proposition}
\label{sec:selection}

\subsection{The Two Hypotheses and Their Honest Epistemic Grades}
\label{subsec:h1h2}

The formal selection layer rests on two Boolean predicates over sorted
block lists, and the corpus itself grades them with unusual candour.
\textbf{H1 (design-level)}, the intersection-form filter: consecutive
block sizes in a sorted partition differ by at most one. The motivation is
the NCG intersection-form constraint on coupled blocks, of which the
$k = \ell \pm 1$ bound of \cref{tab:lattice} is the two-block instance;
the corpus states explicitly that no general dimension-jump theorem of
exactly this form has been derived, and marks the derivation from
spectral-triple axioms as the upgrade path. \textbf{H2 (heuristic)}, the
mass-non-degeneracy filter: at least two summands, all pairwise distinct.
The motivation is the experimentally non-degenerate fermion spectrum; the
corpus states explicitly that ``dynamical mass non-degeneracy'' is a
physical claim about the Yukawa sector for which the combinatorial filter
is a proxy. H2 is what eliminates $[2,2,2]$ --- which is precisely the
\catmark{D}-graded verdict of \cref{tab:lattice}, so the formal layer and
the manuscript lattice are consistent about where the weakest link sits.

\subsection{The Elimination Protocol and the Proposition}
\label{subsec:elimination}

Among the eleven sorted integer partitions of $N=6$, the two filters
eliminate ten:

\begin{table}[H]
\centering\small
\begin{tabular}{@{}lccl@{}}
\toprule
Partition & H1 & H2 & Eliminating filter \\
\midrule
$[6]$ & pass & fail & H2 (simple algebra) \\
$[5,1]$ & fail & pass & H1 (jump 4) \\
$[4,2]$ & fail & pass & H1 (jump 2) \\
$[4,1,1]$ & fail & fail & H1 and H2 \\
$[3,3]$ & pass & fail & H2 (repetition) \\
$[3,2,1]$ & pass & pass & --- \\
$[3,1,1,1]$ & fail & fail & H1 and H2 \\
$[2,2,2]$ & pass & fail & H2 (repetition) \\
$[2,2,1,1]$ & pass & fail & H2 (repetition) \\
$[2,1,1,1,1]$ & pass & fail & H2 (repetition) \\
$[1,1,1,1,1,1]$ & pass & fail & H2 (all equal) \\
\bottomrule
\end{tabular}
\caption{Elimination protocol for the $N=6$ partition space under
H1 and H2, as documented in the formal corpus.}
\label{tab:elimination}
\end{table}

\begin{proposition}[Conditional $N=6$ selection; \catmark{D}; checker
status open]
\label{prop:unique321}
Under H1 and H2, $[3,2,1]$ is the unique admissible partition of $N=6$.
\end{proposition}

The grading of this proposition is layered, and the layers must not be
collapsed. The \emph{logical} content --- exhaustive filtering of a
finite, machine-enumerable list --- is decidable and, once the pending
proof obligations are discharged, will be \catmark{A} within the formal
system. At the audited head it is stated, not machine-checked
(Finding~F1). The \emph{physical} content is bounded above by its
weakest hypothesis: design-level H1 and heuristic H2 place the
proposition at \catmark{D}, and no completion of the Lean proof will
change that; only deriving H1 from intersection-form axioms and replacing
H2 by a stability theorem for the free-energy landscape could. The
falsification structure is explicit and welcome: weakening H1 admits
$[5,1]$ and $[4,2]$; weakening H2 admits $[2,2,1,1]$ and the all-equal
partitions. A selection principle that survives only at one exact filter
strength is a hypothesis about that strength, and is recorded as such.

\begin{openquestion}[Staircase conjecture; \catmark{D}]
\label{oq:staircase}
For general $N$, the Phase-9-admissible partitions consist exclusively of
staircase partitions $[k, k{-}1, \ldots, 1]$ with $k(k{+}1)/2 = N$, so
that admissible total dimensions are triangular numbers. The $N=6$ case
is the first non-trivial instance; $N=10$ and $N=15$ are the next test
cases. Falsification path: exhibit an $N$ and a non-staircase partition
satisfying both H1 and H2.
\end{openquestion}

\subsection{Target-Leakage Audit of the Selection Layer}
\label{subsec:selection-leakage}

The generic layer passes the audit: the filter definitions quantify over
arbitrary partitions, the candidate list appears in the enumeration only
as one member of the complete $p(6) = 11$ inventory, and the ledger
constants appear nowhere outside warning comments. The single violation
is the wrapper of Finding~F2, which is confined to the instance module
and touches neither the filters nor the enumeration. The structural
conclusion of the audit is therefore: the selection result, once
machine-checked, will owe nothing to hard-coded bias --- but it will owe
everything to H1 and H2, which is exactly what its \catmark{D} grade
says.

\section{Candidate Intermediate Algebras}
\label{sec:candidates}

\subsection{The Division-Algebra Programme as a \catmark{D} Candidate}
\label{subsec:division-algebra}

A substantial peer-reviewed programme constructs Standard-Model structure
from the normed division algebras $\mathbb R, \mathbb C, \mathbb H,
\mathbb O$ and their Clifford multiplication algebras. It is the
strongest existing source of a \emph{principled} reason for the specific
factors of $\mathbb C \oplus \mathbb H \oplus M_3(\mathbb C)$, and it is
recorded here as a candidate for $\mathcal A$ at evidence class
\catmark{D}, bounded exactly as its own primary sources bound it. What
the programme establishes, per sources verified in the audit cycle: one
generation of Standard-Model Weyl representations arises as a single copy
of $\mathbb R \otimes \mathbb C \otimes \mathbb H \otimes \mathbb O$
acting on itself; colour and electromagnetism arise from minimal left
ideals of $\mathbb C\ell(6)$, the weak sector from $\mathbb C\ell(4)$~\cite{furey2018};
and gauge structure appears as stabilizer or automorphism structure of a
rigid algebra, with $SU(3)$ the stabilizer of an imaginary octonionic
unit inside $G_2 = \mathrm{Aut}(\mathbb O)$, and the exceptional Jordan
algebra $J_3(\mathbb O)$ providing a parallel route.

What the programme does \emph{not} establish, stated by its own authors:
the gauge group is obtained by restricting to symmetries preserving a
\emph{chosen} subalgebra --- an inserted structure, the same move as
G2/G3; an algebraic foundation for three generations remains elusive,
with generations introduced through an inserted discrete symmetry and
anomaly cancellation flagged as future work~\cite{gourlay2024,furey2025}; and the starting datum is a
rich algebra, not a scalar --- no construction in this programme begins
from a single real scalar field, so the programme presupposes the very
multiplicity that G1 identifies as the unavoidable insertion. It
satisfies the lower bounds of \cref{subsec:lower-bounds} by
\emph{exhibiting} the inserted structure the gap analysis predicts as
necessary. Any UIDT use of it must carry these three bounds explicitly;
it must not be cited as a verified mechanism for unique
$G_{\mathrm{SM}}$ selection.

\subsection{The Relocation Necessity Conjecture}
\label{subsec:rnc}

\begin{openquestion}[Relocation Necessity Conjecture; \catmark{D};
priority research path]
\label{oq:rnc}
In every moduli space of finite real spectral triples in which a
continuous family of Dirac operators varies, the Wedderburn invariants
and the KO-dimension are necessarily discrete (locally constant).
\end{openquestion}

If proved, the conjecture converts the lower bounds of
\cref{subsec:lower-bounds} from ``currently unmet by any mechanism'' to
``provably not obtainable by continuous deformation'' --- the moduli-space
formulation would then not close the multiplicity gap but demonstrate
that it can only be \emph{relocated}, never dissolved, by deformation
arguments. It is an independent, publishable structural question
regardless of the fate of the kernel hypothesis, and the staircase
question (\cref{oq:staircase}) is its natural combinatorial companion.

%============================== R2 PART I ADDITIONS =========================
% Inserted before the Part I claims register.
% Sources: -006 Appendix B (relocated, references adapted); R0 extraction
% (Inner-Derivation Relocation, verified by exact symbolic computation);
% Module D (Casimir algebra, exact). Evidence classes unchanged from -006.
%===========================================================================

\section{Named Constraints: Localisation, Lower Bounds, Moduli}
\label{sec:named-constraints}

\noindent The obstruction analysis of \cref{sec:g1g4} is used repeatedly
in what follows and in the appendix of \cref{sec:connes-attractor}. It is
therefore stated once, in citable form. The three statements below add no
content to \cref{sec:g1g4}; they name it.

\begin{theorem}[Gap Localisation]
\label{thm:gap-localisation}
Let the primitive be a single real scalar degree of freedom whose
symmetries are implemented by canonical linear transformations. Then the
non-derivability of $G_{\mathrm{SM}}$ within the present theorem corpus
is localised to a single step: the enlargement of the abelian
single-mode local algebra to a noncommutative algebra of local
multiplicity $n \geq 2$. Conditions G1--G4 are jointly unsatisfied by
every mechanism surveyed in \cref{sec:candidates}. This is a
localisation, not an impossibility theorem for all conceivable enriched
routes.\catmark{A}\ \emph{re: corpus}
\end{theorem}

\begin{corollary}[Lower Bounds on any Admissible Intermediate Algebra]
\label{cor:algebra-lower-bounds}
Any $\mathcal A$ in a route $\mathbf S \to \mathcal A \to
G_{\mathrm{SM}}$ must carry: (i)~local multiplicity $\geq 2$, that is a
non-abelian fiber; (ii)~rank resources $\geq 4$, against
$\mathrm{rank}(\mathfrak g_{\mathrm{SM}}) = 4$; (iii)~a real structure
with the $J^2$-sign pattern of KO-dimension~6; and (iv)~a chirality
mechanism evading the Nielsen--Ninomiya and finite-2-group obstructions.
A candidate failing one bound is excluded before dynamical analysis.
\end{corollary}

\begin{proposition}[Wedderburn Moduli]
\label{prop:wedderburn-moduli}
For finite spectral triples with complex matrix summands: (i)~the
condensed factors are complex unitary algebras $M_{n_i}(\mathbb C)$;
(ii)~the partition $(n_i)$ is a function of couplings, chemical
potentials, and entropy, not of a selection theorem; (iii)~the
quaternionic summand $\mathbb H$ requires an antilinear structure that a
complex ensemble does not supply. The Wedderburn invariants are, at
present, free moduli~\cite{cacic2011}.\catmark{A}\ \emph{re: corpus}
\end{proposition}

\section{The Inner-Derivation Route: Relocation, Not Closure}
\label{sec:inner-derivation}

\noindent The $d^2 = 0$ obstruction (\cref{prop:d2-obstruction}) is
specific to the exterior derivative on a smooth manifold. The natural
algebraic response, and the one the v4.0 developmental record adopts, is
to replace the exterior derivative by an \emph{inner derivation}
$\delta(\cdot) = [\mathbf S, \cdot]$ on a noncommutative matrix algebra.
Since $\delta \circ \delta$ is not identically zero, the nilpotency that
killed the gradient route no longer applies, and the record concludes
that the foundation for gauge groups has thereby been laid. That
conclusion is incorrect, and the correct statement is a sharper negative
result.

\begin{proposition}[Inner-Derivation Relocation]
\label{prop:inner-derivation}
Let $\mathbf S = \mathbf S^{\ast} \in M_N(\mathbb C)$ and let
$\mathcal A_{\mathbf S} = C^{\ast}(\mathbf S)$ denote the unital
$C^{\ast}$-algebra it generates. Then:
\begin{enumerate}[label=(\roman*),leftmargin=2.2em,itemsep=0.15em]
\item $\mathcal A_{\mathbf S} \cong C(\sigma(\mathbf S))$ is commutative
by the spectral theorem, hence $\mathrm{ad}_{\mathbf
S}\!\restriction_{\mathcal A_{\mathbf S}} \equiv 0$: the inner derivation
vanishes identically on the algebra generated by the primitive, and no
curvature is produced within it.
\item $\mathrm{ad}_{\mathbf S} \neq 0$ on $M_N(\mathbb C)$ whenever
$\mathbf S \notin \mathbb C\!\cdot\!\mathbb 1$, and
$\mathrm{ad}_{\mathbf S}^{2} \neq 0$ there; but $M_N(\mathbb C)$ is not
generated by $\mathbf S$.
\end{enumerate}
Consequently the passage from exterior derivative to inner derivation
evades the identity $d^2=0$ while presupposing the ambient noncommutative
algebra. The obstruction is relocated to the insertion of that algebra,
which is condition~(G1) of \cref{thm:gap-localisation}.\catmark{A}
\end{proposition}

\begin{proof}[Proof sketch]
(i) is the spectral theorem: a self-adjoint operator generates a
commutative $C^{\ast}$-algebra, isomorphic to the continuous functions on
its spectrum; commutativity gives $[\mathbf S, p(\mathbf S)] = 0$ for
every polynomial, and by continuity for every element. (ii) is immediate
from the fact that the centre of $M_N(\mathbb C)$ is $\mathbb
C\!\cdot\!\mathbb 1$. The content of the proposition is not the algebra
but the reading: (i) says the primitive alone generates nothing, and (ii)
says that what does generate something was supplied from outside.
\end{proof}

\begin{remark}[Verification]
\label{rem:inner-derivation-check}
The statement was verified by exact symbolic computation over
$\mathbb Q[i]$ at $N=4$, with no floating-point arithmetic anywhere in
the path: for a generic self-adjoint $\mathbf S$ and polynomials $p, q$,
the commutators $[p(\mathbf S), q(\mathbf S)]$ and
$\mathrm{ad}_{\mathbf S}(p(\mathbf S))$ evaluate to the zero matrix,
while $\mathrm{ad}_{\mathbf S}(X)$ and $\mathrm{ad}_{\mathbf S}^{2}(X)$
are non-zero for an $X \notin C^{\ast}(\mathbf S)$; membership is decided
by the rank test $\mathrm{rank}\{\mathbb 1, \mathbf S, \mathbf S^2,
\mathbf S^3\} = 4$ against $\mathrm{rank}$ of the augmented system $= 5$.
This is internal closure of an identity, not external agreement
(\cref{subsec:closure-vs-agreement}).
\end{remark}

\begin{corollary}[The Same Obstruction, Twice: No Endogenous Modular Time]
\label{cor:modular-relocation}
Let $\mathcal A_{\mathbf S} = C^{\ast}(\mathbf S)$ as in
\cref{prop:inner-derivation}, and let $\omega$ be any faithful normal
state on its weak closure. Since $\mathcal A_{\mathbf S}$ is abelian,
$\omega$ is tracial, and by the Tomita--Takesaki theorem the modular
automorphism group of a trace is the identity: $\sigma^{\omega}_t =
\mathrm{id}$ for all $t$. Hence the primitive alone generates no modular
flow. Any thermal-time construction in the sense of Connes and
Rovelli~\cite{connesrovelli1994}, in which physical time is derived as a
one-parameter automorphism group determined by a thermodynamical state,
therefore requires the same inserted ambient noncommutative algebra that
\cref{prop:inner-derivation} shows to be required for gauge
curvature.\catmark{A}
\end{corollary}

\begin{remark}[Why this matters]
\label{rem:one-obstruction}
The two most load-bearing hopes of the v4.0 architecture --- that
noncommutativity would supply gauge structure, and that modular flow
would supply time --- fail for one and the same reason, and fail on the
same object. The obstruction is not two obstructions. It is the single
insertion of multiplicity identified as (G1) in
\cref{thm:gap-localisation}, seen once through the curvature and once
through the modular group. This unification is a result of the present
draft and is, so far as the audited corpus shows, the sharpest available
statement of what the pregeometric kernel does \emph{not} yet
do.\catmark{A}
\end{remark}

\begin{governancebox}[Withdrawn Reading]
Any statement to the effect that replacing $\partial_\mu$ by a
commutator ``solves'' the $d^2 = 0$ obstruction, or that it ``lays the
foundation for gauge groups,'' is withdrawn. Prestige wording of this
kind appears in the developmental record and is quarantined at
\catmark{E}. The mathematically correct claim is
\cref{prop:inner-derivation}: an evasion of the identity, purchased by
the insertion of the very multiplicity that G1 identifies as the
unavoidable input. What the passage buys is a sharper obstruction, not a
foundation.\catmark{A}
\end{governancebox}

\section{Two Exact Auxiliary Results}
\label{sec:auxiliary-exact}

\subsection{The Casimir Level Structure}
\label{subsec:casimir-levels}

For a block of size $n$ in the fuzzy-sphere sector of a condensed matrix
ensemble, the quadratic Casimir eigenvalue is
\begin{equation}
 Q_n \;=\; \frac{\alpha^2\,(n^2-1)}{4},
 \label{eq:casimir}
\end{equation}
in agreement with the angular-momentum form $j(j+1)$ for $n = 1,\ldots,8$
under $j = (n-1)/2$ and the stated normalisation. In particular
$Q_2 = 0.75\,\alpha^2$.\catmark{A} This is exact algebra and is recorded
because the $n=2$ block is the fragile one: it is the block whose
detection determines whether a Standard-Model-like configuration can be
resolved at all. Whether it survives a given fluctuation level is
\emph{not} an algebraic question, and \cref{lim:L12x} states why.

\begin{limitation}[L12$'$: The Fluctuation Floor Is a Thermodynamic Postulate]
\label{lim:L12x}
Whether the fluctuation floor of a condensed phase is set by the lower
level gap $\Delta_{\min}$ or by a higher operator-norm scale is a
thermodynamic postulate, not a mathematical choice: it fixes the
competition between potential and entropy and therefore determines,
rather than derives, the phase structure. Fixing it a priori in order to
make a desired block configuration resolvable is target leakage
(\cref{prop:target-leakage}). Two facts sharpen the limitation.
\emph{First}, the spectral-gap problem is algorithmically undecidable for
general local translation-invariant Hamiltonians, and there exist models
for which the presence or absence of a gap is independent of the axioms
of mathematics~\cite{cubitt2015}; that theorem is about the class of such
questions, and asserts nothing about the Yang--Mills gap or about
$\Delta$.\catmark{A} \emph{Second}, and consequently, any infrared
operator matching of the kind Limitation~L12 demands presupposes a fixed
floor. L12 is therefore not merely the declarative statement that no
operator matching has been performed; it is the structural statement that
the matching cannot begin without a postulate that is undecidable in
general and leakage-prone in practice.\catmark{A}/\catmark{D}
\end{limitation}

\subsection{The Division-Algebra Programme Is Exogenous}
\label{subsec:division-exogeneity}

\begin{remark}[Sharpened critique of \cref{subsec:division-algebra}]
\label{rem:division-exogeneity}
It is not sufficient to observe that the division-algebra programme
``presupposes multiplicity.'' The sharper and more damaging point is
that its starting datum is a \emph{continuous, a-priori-given} algebraic
structure --- $\mathbb R, \mathbb C, \mathbb H, \mathbb O$ and their
Clifford algebras --- whereas the v4.0 programme begins from a discrete
relational zero in which no continuum, no field, and no algebra of
functions is available. The programme therefore does not evade G1; it
does not meet G1 either. It occupies a different starting point
altogether, and its successes are successes of \emph{that} starting
point. Any \UIDT\ use of the programme must state this exogeneity, and no
result of the programme transfers to the kernel hypothesis without a
construction of the relevant algebra from the kernel --- which does not
exist.\catmark{D}
\end{remark}

\section{Dynamical Attractors toward the Connes Vacuum}
\label{sec:connes-attractor}

\begin{governancebox}[Status of This Section]
This section is an advisory audit analysis and a \emph{programme
charter}, not a result. It does not establish, and must not be cited as
establishing, a thermodynamic emergence of the algebra
$\mathbb{C}\oplus\mathbb{H}\oplus M_3(\mathbb{C})$, a dynamical necessity
of KO-dimension~6, or an anomaly-theoretic uniqueness of the $(1,2,3)$
partition. Established external inputs are carried at \catmark{B};
conditional selections at \catmark{B}\ under stated hypotheses; the
attractor content itself as programme requirements at \catmark{D}\ or
conjectures at \catmark{E}. It operates under
\cref{thm:gap-localisation}, \cref{cor:algebra-lower-bounds}, and
\cref{prop:target-leakage}. An AI-generated research memo in the
developmental record proposes the mechanisms surveyed here; per the
corpus discipline of \cref{subsec:provenance}, that memo is contextual
\catmark{E}\ material and is not cited as evidence. The research form
$\mathbf S \to \mathcal A \to G_{\mathrm{SM}}$ remains \catmark{D}.
\end{governancebox}

\subsection{From External Postulates to Dynamical Selection}
\label{subsec:attractor-programme}

The noncommutative-geometric classification reaches
$\mathbb{C}\oplus\mathbb{H}\oplus M_3(\mathbb{C})$ only through an
externally imposed axiom set. The attractor programme asks whether each
postulate can be replaced by a dynamical or thermodynamic selection
principle acting on a larger configuration space, so that the Connes
vacuum would arise as a stable attractor rather than as an axiom. The
programme is stated here in falsification-disciplined form; none of its
attractor claims is established.

\subsection{The Matrix Arena and the Exact Classical Comparison}
\label{subsec:matrix-arena}

\paragraph{External model.}
A concrete arena in which a block partition is a representation-theoretic
object is the three-matrix Yang--Mills--Chern--Simons model
\cite{azuma2004}, with Hermitian $N\times N$ matrices and action
\begin{equation}
\label{eq:ymcs-action}
S[X]=N\,\operatorname{Tr}\!\left(
-\frac14[X_a,X_b][X_a,X_b]
+\frac{2i\alpha}{3}\,\varepsilon_{abc}X_aX_bX_c
\right).
\end{equation}
This action is an external Stratum-II input.\catmark{B}

\paragraph{Stationary branches, not an exhaustive classification.}
Direct variation of \cref{eq:ymcs-action} gives
\begin{equation}
\label{eq:ymcs-eom}
[X_b,[X_b,X_a]]+i\alpha\,\varepsilon_{abc}[X_b,X_c]=0.
\end{equation}
Consequently, $X_a=\alpha L_a$ is a stationary branch whenever
$[L_a,L_b]=i\varepsilon_{abc}L_c$; mutually commuting triples form
another stationary branch. These statements identify named branches.
They do not classify every stationary point of the model.\catmark{A}

\paragraph{Exact action comparison.}
For the representation branch, substitution and the quadratic Casimir
identity give
\begin{equation}
\label{eq:ymcs-classical-value}
S[\alpha L]= -\frac{N\alpha^4}{6}\operatorname{Tr}(L_aL_a).
\end{equation}
For one irreducible representation of dimension $N$,
\(
\operatorname{Tr}(L_aL_a)=N(N^2-1)/4
\).
For $k$ equal irreducible blocks of dimension $n=N/k$,
\(
\operatorname{Tr}(L_aL_a)=N(n^2-1)/4
\).
The exact finite-$N$ magnitude ratio is therefore
\begin{equation}
\label{eq:ymcs-equal-block-ratio}
\frac{|S_k|}{|S_1|}
=\frac{n^2-1}{N^2-1}
=\frac{N^2-k^2}{k^2(N^2-1)},
\end{equation}
which approaches $1/k^2$ only at fixed $k$ as $N\to\infty$. Mutually
commuting triples have $S=0$. The equal-block representation branch is
therefore ordered at tree level by the exact finite-$N$ expression; this
is not an all-orders no-go theorem, a global classification of extrema,
or a derivation of a Standard-Model-like partition.\catmark{A}

\paragraph{UIDT boundary.}
The algebra above is exact under the displayed assumptions. External
phase-structure claims require primary numerical sources, while the use
of the comparison as a boundary on the UIDT partition-selection
programme remains an internal \catmark{D} interpretation. No UIDT target
value enters the action, initialization, detector, or conclusion.

\subsection{Matrix Condensation and the Wedderburn Moduli Problem}
\label{subsec:attractor-condensation}

\paragraph{What is established.}
Zero-dimensional matrix ensembles of IKKT type~\cite{ikkt1997} and their
bosonic deformations exhibit condensation transitions in which
block-diagonal backgrounds form dynamically, breaking
$\mathrm{U}(N)\to\prod_i\mathrm{U}(n_i)$; fuzzy-sphere backgrounds
nucleate and dissolve in first-order transitions governed by the
competition between the classical potential and the entropy of
fluctuations, and off-diagonal (bifundamental) modes acquire masses
growing with the separation of the condensed
blocks~\cite{azuma2004,grosse2005,steinacker2010}.\catmark{B} In
free-energy language, $F[(n_i)] = U[(n_i)] - T\,S[(n_i)]$ on the
partition moduli: entropy-dominated regimes favour either a single
dominant block or near-symmetric partitions, while potential-dominated
regimes track the discrete minima of the engineered potential. In no
established regime is the asymmetric partition $(1,2,3)$ singled out by a
structural principle.\catmark{B}

\begin{remark}[The moduli problem restated]
\label{rem:moduli-restated}
\cref{prop:wedderburn-moduli} remains binding at every temperature.
Entropy maximisation is structurally blind to the $(1,2,3)$ hierarchy,
and any claim that the Connes algebra is the unique global thermodynamic
minimum of an unconstrained ensemble is rejected.\catmark{A}
\end{remark}

\begin{governancebox}[Programme Requirement PR-B1: Partition Selection \catmark{D}]
Exhibit a concrete matrix action and a coupling/temperature window in
which the global free-energy minimiser over Wedderburn partitions is
$(1,2,3)$, with the minimiser stable as $N\to\infty$, under the
anti-target-leakage protocol of \cref{prop:target-leakage}: all
parameters fixed by observables independent of the partition; the
partition read out post hoc by a sealed evaluation; null controls run on
scrambled reference partitions (for example $(1,2,4)$ and $(2,2,2)$),
which must \emph{not} be selected by the same procedure. Downgrade
condition: if the minimiser tracks engineered parameters rather than a
structural principle, the mechanism is selection by hand and is recorded
at \catmark{E}.
\end{governancebox}

\subsection{Anomaly Freedom, Unimodularity, and the Quiver Flow Conjecture}
\label{subsec:attractor-anomaly}

\paragraph{What is established.}
Within the noncommutative-geometric Standard Model the physical gauge
group is the unimodular unitary group of $\mathcal{A}_F$, and the
unimodularity condition is tied to the cancellation of the chiral gauge
anomalies of the spectral fermion content~\cite{agbm1995,cc2008}: anomaly
freedom prunes the unitary group of an \emph{already chosen} algebra to
$G_{\mathrm{SM}}$ up to a finite kernel; it does not select the
algebra.\catmark{B} The configuration space on which any selection
dynamics would act is likewise established: finite real spectral triples
are classified combinatorially by Krajewski
diagrams~\cite{krajewski1998,paschkesitarz1998}, and the order-one
condition acts there as a combinatorial admissibility constraint on the
diagram.\catmark{A}

\paragraph{Verdict and mechanism must be distinguished.}
Two cautions follow from a primary-source audit of candidate partitions,
both illustrating that a correct verdict can rest on an incorrect stated
mechanism. First, $\mathbb{C}\oplus\mathbb{H}\oplus M_3(\mathbb{C})$ is
anomaly-admissible: the unimodularity reading of anomaly freedom fixes
the single anomaly-free hypercharge assignment~\cite{agbm1995}. This is a
property of the chosen representation and hypercharge, \emph{not} of the
block partition alone; anomaly freedom is not a partition-intrinsic
test.\catmark{B} Second, the partition $(3,6)$ fails for a
\emph{specific} reason: for two simple summands of matrix sizes $k$ and
$\ell$, a non-degenerate intersection form (Poincar\'e duality) requires
$k=\ell$ or $k=\ell\pm 1$~\cite{iss2004}; since $|3-6|=3$ this bound is
violated, so $(3,6)$ is excluded by the intersection-form constraint, not
by an order-one or anomaly argument.\catmark{B} The symmetric partition
$(2,2,2)$ is not excluded by anomaly freedom at all; the established
obstruction to such symmetric configurations in the irreducible
classification is the \emph{degeneracy of the fermion mass
spectrum}~\cite{iss2004}, a dynamical-non-degeneracy criterion external
to G1--G4. Its admissibility under G1--G4 proper therefore remains
open.\catmark{D}\ This is the primary-source basis of the corrected
lattice of \cref{tab:lattice} and of tension entry~T8.

\paragraph{What is not established.}
No renormalisation-group flow on the space of Krajewski diagrams is
currently defined. By the scaffold discipline of
\cref{subsec:gferg-scaffold}, a claimed flow must specify, in advance,
the flow equation, the regulator analogue on the discrete moduli, the
truncation, the boundary conditions, and the failure criteria; absent
these, renormalisation-group vocabulary on quiver space is not a
calculation. Moreover, uniqueness manifestly fails at the level of the
known classification: relaxation of the order-one condition enlarges the
solution set to Pati--Salam~\cite{ccvs2013,farnsworth2015}.\catmark{A}

\begin{remark}[Conjecture C-B1: anomaly-constrained quiver attractor, status \catmark{E}]
\label{rem:quiver-conjecture}
Within the anomaly-free, order-one-admissible stratum of Krajewski
diagrams there exists a well-defined coarse-graining flow whose unique
stable fixed point at three summands is the $(1,2,3)$ diagram. No such
flow is defined; no uniqueness argument exists; the conjecture inherits
the full burden of \cref{prop:target-leakage}, since a flow whose fixed
point is verified only against the known Standard-Model diagram predicts
nothing.\catmark{E}
\end{remark}

\subsection{Real Structure and KO-Dimension 6: A Conditional Selection}
\label{subsec:attractor-ko6}

\paragraph{What is established, conditionally.}
The KO-dimension of a real spectral triple is fixed modulo~8 by the sign
triple $(\varepsilon,\varepsilon',\varepsilon'')$ in
$J^2=\varepsilon$, $JD=\varepsilon'DJ$, $J\gamma=\varepsilon''\gamma J$;
KO-dimension~6 carries $(1,1,-1)$. Within the almost-commutative spectral
Standard Model, imposing simultaneously (i)~the elimination of the
fermion quadrupling of the Euclidean formulation and (ii)~the
admissibility of Majorana mass terms for right-handed neutrinos selects
KO-dimension~6 for the finite
triple~\cite{barrett2007,connes2006,ccm2007}.\catmark{B} With
$\varepsilon''=-1$ the finite Dirac operator admits the block form
\[
D_F=\begin{pmatrix} S & T^{*}\\[2pt] T & \bar S\end{pmatrix},
\qquad T=T^{T},
\]
where $S$ contains the Dirac masses and the symmetric block $T$ couples
right-handed neutrinos to their conjugates, enabling the seesaw
mechanism~\cite{connes2006,ccm2007}.\catmark{B} The selection is
\emph{conditional}: requirements (i) and (ii) are physical inputs, not
consequences of any dynamics. Necessity of KO-dimension~6 is therefore
necessity relative to (i) and (ii), and nothing stronger may be asserted;
formulations such as ``absolute mathematical necessity'' or ``rigorously
proven not to be an insertion'' are rejected wherever they appear in the
developmental record.\catmark{A}

\paragraph{What is open.}
A dynamical origin of the real structure $J$ itself is not established.
The heuristic that unoriented configurations are entropically or
energetically disfavoured on a stochastic substrate is exactly that, a
heuristic: no defined ensemble, measure, or flow currently demonstrates
it, and \cref{prop:wedderburn-moduli}(iii) records that the complex
substrate does not supply $J$.\catmark{A}

\begin{governancebox}[Programme Requirement PR-B2: Endogenous Real Structure \catmark{D}]
Exhibit a substrate dynamics whose stable configurations are invariant
under an antilinear involution $J$ with the KO-dimension-6 sign triple,
where $J$-invariance emerges as a stability or extremality property of
the ensemble rather than as a restriction imposed on the configuration
space at the outset, and show that configurations without such $J$ are
dynamically disfavoured. Downgrade condition: if $J$-invariance is
imposed on the measure by hand (an orientifold projection inserted at
step zero), the mechanism is an insertion and PR-B2 is not met.
\end{governancebox}

\subsection{The Residual Freezing Gap and a Spectral Falsification Protocol}
\label{subsec:attractor-freezing}

\paragraph{Layer-C declaration.}
The deepest unresolved step of the programme is the dynamical origin of
the bipartite order-one condition \emph{prior to} any flow: the condition
is doing genuine selection work, since its relaxation enlarges the
algebra to Pati--Salam~\cite{ccvs2013,farnsworth2015}, yet no mechanism
explains why the commutant structure of an emergent geometry should
freeze into the first-order form at all. This is recorded as an open
question, not as a difficulty in the course of being
resolved.\catmark{A}

\begin{openquestion}[The freezing gap]
\label{oq:freezing-gap}
What dynamical or thermodynamic principle, acting on the configuration
space of finite real spectral triples, enforces the order-one condition
$[[D_F,a],Jb^{*}J^{-1}]=0$ as a stability or extremality property, prior
to and independently of any imposed axiom? No candidate principle is
currently established; PR-B1 and PR-B2 are necessary inputs to any
answer but are not sufficient for it.\catmark{D}
\end{openquestion}

\paragraph{Falsification handle.}
If a concrete attractor realisation is ever exhibited, it acquires a
discriminating observable through the spectral statistics of the
condensed finite Dirac operator. Random Dirac ensembles over finite
spectral triples are an established, simulable
arena~\cite{barrettglaser2016}, and with
$Z(\beta+it)=\operatorname{Tr}e^{-(\beta+it)D_F^{2}}$ the spectral form
factor $g(t;\beta)=\lvert Z(\beta+it)\rvert^{2}/Z(\beta)^{2}$
distinguishes random-matrix spectral rigidity (the dip--ramp--plateau
profile, with the symmetry class fixed by the sign triple) from
Poissonian level statistics~\cite{cotler2017}. The protocol is a
\catmark{D}\ proposal; no simulation or measurement of it exists for any
\UIDT-relevant ensemble.

\begin{killswitch}
If a concrete attractor realisation (PR-B1 together with PR-B2) is
exhibited and the spectral form factor of its condensed phase is
incompatible with the rigidity class fixed by that realisation's own sign
triple, the realisation is falsified. If, instead, stable condensed
phases are exhibited whose algebras lack $G_{\mathrm{SM}}$ entirely, any
uniqueness reading of the Connes vacuum as attractor is invalidated and
Conjecture C-B1 is to be withdrawn. Absent any realisation, the research
form remains \catmark{D}\ and no evidence class of the main text is
altered by this section.\catmark{D}
\end{killswitch}

\subsection{Attractor Claims Register and Reproduction Note}
\label{subsec:attractor-claims}

\begin{table}[htbp]
\centering\footnotesize
\begin{tabular}{@{}lp{5.3cm}lp{4.3cm}@{}}
\toprule
\# & Claim & Class & Falsification / downgrade \\
\midrule
APP-B-01 & Block condensation with entropy/potential competition and massive bifundamentals is established in matrix models & \catmark{B} & source audit of the cited literature \\
APP-B-02 & Entropy maximisation is blind to $(1,2,3)$; no functional selects the partition & \catmark{A} re corpus & a selection theorem satisfying PR-B1 under \cref{prop:target-leakage} \\
APP-B-03 & Anomaly freedom prunes $\mathcal U(\mathcal A_F)$ to $G_{\mathrm{SM}}$; it does not select $\mathcal A_F$ & \catmark{B} & a theorem deriving the algebra from anomaly data alone \\
APP-B-04 & The order-one condition is a combinatorial admissibility constraint; the anomaly-consistent stratum is non-unique & \catmark{A} & exhibition of uniqueness at three summands \\
APP-B-05 & Quiver-attractor uniqueness is Conjecture C-B1 & \catmark{E} & a defined flow with regulator, truncation, boundary conditions, and a leak-free uniqueness proof \\
APP-B-06 & KO-6 is selected conditionally by fermion-doubling resolution plus Majorana admissibility & \catmark{B} cond. & a dynamical derivation of inputs (i)/(ii) or of $J$ \\
APP-B-07 & The dynamical origin of the order-one condition is open; the freezing gap stands & \catmark{A}/\catmark{D} & a principle answering \cref{oq:freezing-gap} under PR-B1/PR-B2 \\
APP-B-08 & Spectral-form-factor protocol as discriminator for any future realisation & \catmark{D} & executable only upon a concrete realisation \\
\bottomrule
\end{tabular}
\caption{Attractor claims register. No entry establishes an emergence.}
\label{tab:attractor-claims}
\end{table}

\paragraph{Reproduction note.}
The established input of \cref{subsec:attractor-anomaly} has a
machine-checkable core: the four anomaly sums of one Standard-Model
generation vanish identically in exact rational arithmetic for the
unimodular hypercharge assignment
$(Y_Q,Y_{u^c},Y_{d^c},Y_L,Y_{e^c})=(\tfrac{1}{6},-\tfrac{2}{3},\tfrac{1}{3},-\tfrac{1}{2},1)$:
\begin{verbatim}
python3 -c "from fractions import Fraction as F;
Y={'Q':(F(1,6),6),'uc':(F(-2,3),3),'dc':(F(1,3),3),'L':(F(-1,2),2),
'ec':(F(1),1)};
print(sum(m*y**3 for y,m in Y.values()),
 sum(m*y for y,m in Y.values()),
 2*Y['Q'][0]+Y['uc'][0]+Y['dc'][0],
 3*Y['Q'][0]+Y['L'][0])"
# -> 0 0 0 0  ([U(1)]^3, grav^2-U(1), [SU(3)]^2-U(1), [SU(2)]^2-U(1);
#    exact arithmetic; certifies the assignment, not any UIDT claim)
\end{verbatim}
Executed and confirmed in the present audit pass. It certifies the
anomaly arithmetic of the established input only; no attractor statement
of this section is machine-checkable, because no attractor realisation
exists to check.\catmark{A}

\begin{governancebox}[Forbidden Inferences]
Nothing in this section licenses the statements that the Connes vacuum
emerges thermodynamically, that KO-dimension~6 is dynamically necessary,
that anomaly cancellation selects the $(1,2,3)$ algebra, or that \UIDT\
derives $G_{\mathrm{SM}}$. The established inputs are external literature
results on matrix condensation, unimodularity, and conditional
KO-dimension selection; the attractor content consists of two programme
requirements (PR-B1, PR-B2), one conjecture (C-B1), and one open question
(\cref{oq:freezing-gap}). \Cref{thm:gap-localisation} and
\cref{cor:algebra-lower-bounds} remain the binding constraints, and
emergence wording is rejected wherever it appears in the developmental
record.
\end{governancebox}

%============================ END R2 PART I ADDITIONS =======================

\section{Part I Claims Register}
\label{sec:part1-claims}

\begin{table}[htbp]
\centering\footnotesize
\begin{tabular}{@{}lp{5.3cm}lp{4.3cm}@{}}
\toprule
\# & Claim & Class & Falsification / downgrade \\
\midrule
K-01 & Kernel primitive (DIR-S-01): pregeometric operator replaces
Historical scalar ansatz as candidate primitive & \catmark{D} & conditions of the
kill-switch in \cref{subsec:dir-s-01} \\
K-02 & $d^2=0$: scalar-gradient route yields $F \equiv 0$ &
\catmark{A} & counterexample in ordinary exterior calculus (none
expected) \\
K-03 & Inner-derivation relocation: $\mathrm{ad}_{\mathbf S}$ vanishes on
$C^{\ast}(\mathbf S)$; non-zero curvature requires the inserted ambient
algebra (\cref{prop:inner-derivation}) & \catmark{A} & a construction
generating a noncommutative algebra from $\mathbf S$ alone \\
K-04 & No endogenous modular time: $\sigma^\omega_t=\mathrm{id}$ on
$C^{\ast}(\mathbf S)$; thermal time needs the same inserted algebra
(\cref{cor:modular-relocation}) & \catmark{A} & a non-tracial faithful normal
state on an abelian von Neumann algebra \\
G-01 & No gapless ab initio a structureless real scalar $\to G_{\mathrm{SM}}$ in the current
theorem corpus & \catmark{A} re: corpus & a theorem generating a
noncommutative $\mathcal A$ endogenously from single-boson data \\
G-02 & Obstruction localized to G1--G4 & \catmark{A} & a mechanism
satisfying G1--G4 jointly \\
G-03 & G1 = non-abelian fiber; $J$ exists at $\dim V_{ij}=1$;
$[3,2,1]$ satisfies G1 & \catmark{A} (textual) & a primary source
requiring block repetition for $J$ \\
G-04 & Lower bounds on $\mathcal A$ (multiplicity, rank, KO-6 $J$,
chirality) & \catmark{A}/\catmark{D} & a theorem fixing the invariants
endogenously \\
G-05 & $[2,2,2]$ rejection rests on an unverified mass-degeneracy
argument & \catmark{D} & primary-source theorem settling the case
either way \\
L-01 & Lean layers Foundation/RealStructure/FOC/Krajewski/DeltaScale
closed at head \texttt{61bc24da} & \catmark{A} (formal) & scan or build
exhibiting an open obligation in these files \\
L-02 & 51 open obligations at head; register under-reports them
(F1/F3); one target-shaped definition (F2) & \catmark{A} (audit fact) &
a head at which the counts or findings no longer hold \\
S-01 & Conditional selection: H1 $\wedge$ H2 $\Rightarrow$ unique
$[3,2,1]$ at $N=6$ & \catmark{D}; checker open & discharge of F1
upgrades only the formal layer; weakening either filter voids
uniqueness \\
S-02 & Staircase conjecture for general $N$ & \catmark{D} &
non-staircase partition passing H1 $\wedge$ H2 \\
C-01 & Division-algebra programme is a \catmark{D} candidate for
$\mathcal A$; selection by inserted stabilizer; three generations open
& \catmark{D} & verified unique derivation plus spontaneous three
generations from the algebra alone \\
C-02 & Relocation Necessity Conjecture & \catmark{D} & continuous
family changing Wedderburn data or KO-dimension \\
\bottomrule
\end{tabular}
\caption{Part I claims register. No entry authorizes any ledger change;
all upgrades pass through PI governance.}
\label{tab:part1-claims}
\end{table}

%========================== END P2 PART I CONTENT ==========================


%=============================================================================
% PART II — THE CONTINUOUS LIMIT
%=============================================================================
\newpage
\part{The Open Passage to Apparent Continuity\texorpdfstring{\ \catmark{E}}{[E]}-Capped}
\label{part:continuous-limit}

\noindent Part~II examines how apparent continuity, locality, and temporal
flow could arise from a discrete kernel --- and why every currently
available account of that passage remains conjectural. The thermodynamic
limit, coarse-graining functors, modular-flow readings of time, and the
Markov-blanket-bounded observer are treated as bounded interpretive
material: capped at \catmark{E} wherever they touch phenomenology, barred
from validating any primitive, and accompanied by the documented no-go
points --- circularity risks in thermal-time readings, the absence of any
theorem forcing $(3{+}1)$ Minkowski emergence, and UV/IR-mixing
obstructions to emergent locality. The Born rule remains Stratum~II
standard quantum mechanics; no coarse-graining map in this part
constitutes a derivation of it.

%%ANCHOR-P3-PART-II-CONTENT%%

%============================ P3 PART II CONTENT ===========================
% Inserted at ANCHOR-P3-PART-II-CONTENT by fail-fast patch script.
% Relocates and adapts UIDT Ontology v3.9.9 (-006) Chapters 10-14 under the
% D20 reframe: the microscopic level of the compression picture is now the
% pregeometric kernel candidate of Part I; scalar history is appendix-only,
% observer-accessible level treated in Part III. All v3.9 observer guards
% are carried over unchanged in force. Author-name comparisons are cited
% in prose; \cite bindings attach in Phase P5 with the verified bibliography.
%===========================================================================

\section{The Unresolved Thermodynamic Passage: From Kernel to Apparent Continuity}
\label{sec:thermo-passage}

\subsection{The Programme Statement}
\label{subsec:thermo-programme}

Part~I ends with a finite algebraic candidate and a list of things it has
not earned; Part~II begins with the largest of them. If the primitive is
discrete and pregeometric, then continuity, locality, dimensionality, and
temporal flow are not available for free: each must be produced by some
limit, and the limit must be shown to exist, to be unique in the relevant
sense, and to yield the specific structure --- a $(3{+}1)$-dimensional
Lorentzian effective arena --- in which the Part-III field theory lives.
The framework's proposal is thermodynamic in character: apparent
continuity is to arise as a large-$N$, coarse-grained limit of kernel
degrees of freedom, in the way that hydrodynamics arises from molecular
dynamics, with any historical scalar ansatz confined to Appendix~\ref{app:scalar-museum} as a
order-parameter-level description of that limit.\catmark{D}

This is a programme statement, not a result, and the distinction is
enforced by the same falsification-first rule that governs Part~I. No
construction in this draft produces a continuum from the kernel; what
this Part provides is the honest inventory of what such a construction
would have to overcome, together with the observer-theoretic apparatus
--- compression, coarse-graining, self-models --- that any such limit
presupposes when it speaks of ``appearance.''

\subsection{Documented No-Go Points}
\label{subsec:no-go-points}

Three obstructions are on record from the v4.0 research cycle, and each
is carried here as a boundary condition on the programme rather than
buried in optimism.

\emph{No dimensional-emergence theorem.} There is no general theorem, in
any of the surveyed discrete-to-continuum programmes, forcing the
emergence of a $(3{+}1)$-dimensional Minkowski-signature arena from a
finite algebraic or combinatorial primitive. Candidate mechanisms exist
--- spectral-dimension flows, causal-set dynamics, tensor-network
geometry --- but dimensionality and signature are, in every case examined,
either inserted or selected among alternatives.\catmark{D} The kernel
programme inherits this gap in full.

\emph{Ultraviolet--infrared mixing.} Field theories on noncommutative or
otherwise discretized underlying structures generically exhibit
UV/IR mixing: short-distance structure contaminates long-distance
behaviour instead of decoupling from it. Any claim that the kernel's
discreteness is invisible at effective scales must confront this
mechanism explicitly; effective-theory decoupling cannot simply be
assumed where the standard locality hypotheses that underwrite it are
themselves in question.\catmark{D}

\emph{Thermal-time circularity risk.} The most natural candidate for
emergent temporality in an algebraic setting --- modular flow in the
sense of Tomita--Takesaki, read as physical time via the thermal-time
hypothesis --- carries a documented circularity risk: the selection of
the state whose modular flow is to count as time threatens to presuppose
the temporal or thermodynamic structure it is meant to produce. Until a
state-selection principle independent of temporal input is exhibited,
emergent-time claims remain at \catmark{E}, discussable as proposed
structure only.\catmark{E}

\subsection{Modular Time and the Relocation of the Thermal-Time Hypothesis}
\label{subsec:modular-time}

The thermal-time circularity flagged in \cref{subsec:no-go-points} can now
be stated with the theorem that makes it precise and the Part-I result that
sharpens it. The apparatus is standard and belongs to Stratum~II; the
reading placed on it, and the obstruction it inherits from the kernel, are
Stratum~III.

The Tomita--Takesaki theorem is the relevant piece of established
operator-algebra. For a von Neumann algebra $\mathcal M$ with a faithful
normal state $\omega$ admitting a cyclic and separating vector, there is a
canonical one-parameter \emph{modular automorphism group}
$\sigma^{\omega}_t(A) = \Delta_\omega^{\,it}\, A\, \Delta_\omega^{-it}$,
where $\Delta_\omega$ is the modular operator of the pair
$(\mathcal M, \omega)$. The modular operator $\Delta_\omega$ is a property
of the state on the algebra and must not be conflated with the UIDT
spectral gap $\Delta$ of Part~III; the two are unrelated objects that
happen to share a letter. The theorem is a mathematical fact about von
Neumann algebras and carries no interpretive content by
itself.\catmark{A}

The thermal-time hypothesis of Connes and
Rovelli~\cite{connesrovelli1994} places an interpretation on it: in a
generally covariant theory that supplies no preferred external time, one
may identify physical time with the modular flow $\sigma^{\omega}_t$
determined by the thermodynamical state, so that time becomes a
state-dependent rather than a background structure. This is a proposal
about what the modular flow \emph{means}, not a theorem that it \emph{is}
physical time; it is carried here at \catmark{E}, as interpretive structure
only.

Two inputs stand between this reading and any emergent-time claim in the
present framework, and \cref{cor:modular-relocation} fixes the first of
them exactly. If the primitive is the self-adjoint kernel operator
$\mathbf S$ alone, then $C^{\ast}(\mathbf S)$ is abelian, every faithful
normal state on its weak closure is tracial, and the modular group of a
trace is the identity: $\sigma^{\omega}_t = \mathrm{id}$ for all $t$. The
primitive alone therefore generates no modular flow, and hence no thermal
time. A nontrivial modular group requires precisely the inserted ambient
noncommutative algebra that condition~(G1) of \cref{thm:gap-localisation}
requires for gauge curvature; by \cref{rem:one-obstruction} it is the
\emph{same} insertion, seen through the modular group rather than through
the curvature. Emergent time is thus not available at the kernel level; it
becomes discussable only after the open insertion that Part~I records as
unsolved.

The second input remains open even after that insertion is granted, and it
is the circularity proper. To read $\sigma^{\omega}_t$ as time one must
first select the state $\omega$ whose flow is to count as time. In every
construction surveyed, that selection is fixed by a condition that is
itself temporal or thermodynamic --- a KMS condition at some inverse
temperature, a distinguished Hamiltonian, an equilibrium prescription ---
that is, by importing the very structure the flow was meant to produce. The
relocation of \cref{cor:modular-relocation} does not remove this second
input; it isolates it, by showing that the algebra and the state are two
separate unmet requirements rather than one.

\begin{openquestion}[State Selection Without Temporal Input]
\label{oq:state-selection}
Can a faithful normal state on the inserted ambient noncommutative algebra
be selected from pregeometric combinatorial data alone --- from the
kernel's Krajewski and multiplicity structure --- without importing a KMS,
equilibrium, or Hamiltonian condition that already presupposes temporal or
thermodynamic structure? Until such a state-selection principle is
exhibited and shown to be non-circular, modular flow as a generator of
physical time remains at \catmark{E}, and the thermal-time route to
emergent temporality is a proposed structure, not a
mechanism.\catmark{E}
\end{openquestion}

\subsection{Krylov Complexity and the Conditional Collapse of State-Space Differentiability}
\label{subsec:krylov-collapse}

A second candidate obstruction to the thermodynamic passage concerns not
time but smoothness: whether the space of states over which a continuum
limit would have to be differentiable remains differentiable at all in the
regime the kernel dynamics would occupy. The established material is the
operator-growth literature, and the discipline required is to cite it for
what it shows and not for the inference the framework draws from it.

Operator growth under Heisenberg evolution can be quantified by the Lanczos
coefficients $b_n$ generated when the Liouvillian is tridiagonalised on the
Krylov basis of a seed operator; the associated Krylov complexity measures
how far the evolving operator has spread along that basis. Parker
et~al.~\cite{parker2019} advance the Universal Operator Growth Hypothesis:
in generic non-integrable systems, in the thermodynamic limit, the Lanczos
coefficients grow asymptotically at the maximal linear rate
$b_n \sim \alpha n$, which in turn bounds the quantum Lyapunov exponent by
$\lambda_L \le 2\alpha$. This is a hypothesis, verified in specific models
(spin chains, SYK, certain classical systems), not a universal theorem, and
it is stated for the thermodynamic-limit setting. Rabinovici
et~al.~\cite{rabinovici2021} treat the complementary case relevant to a
finite $N\times N$ kernel: in a finite-entropy, finite-dimensional Krylov
space, complexity growth is bounded and terminates, the operator
wavefunction reaching the edge of the space; in chaotic models
($\mathrm{SYK}_4$) the complexity saturates at a random-matrix-governed
plateau, while in integrable ones ($\mathrm{SYK}_2$) it stays exponentially
below.

The inference the framework draws is conditional and belongs to it, not to
either source. \emph{If} an admissible kernel dynamics is chaotic in the
sense that its Lanczos sequence reaches the finite-dimensional saturation
plateau identified for $\mathrm{SYK}_4$-type systems, \emph{then} the
coarse-graining map from kernel to effective continuum cannot rely, in that
regime, on a differentiable metric structure over the space of states: the
smooth state-space geometry a continuum limit presupposes --- a
Fubini--Study-type metric on the projective state space --- is not the
object that survives once operator complexity saturates and the local
Lanczos data is replaced by random-matrix statistics. The claim is a
boundary on where the thermodynamic passage of \cref{sec:thermo-passage}
could produce a differentiable arena, not a mechanism that produces one,
and still less an explanation of any specific short-distance
breakdown.\catmark{D}

Two disciplinary statements bound this subsection. Neither
\cite{parker2019} nor \cite{rabinovici2021} asserts any collapse of
state-space differentiability; both concern operator-complexity growth and
its bounds, and the bibliography records this scope explicitly. The
differentiability consequence is a \UIDT\ inference at \catmark{D},
conditional on an antecedent --- plateau saturation for an admissible
kernel dynamics --- that is itself unestablished, since no admissible
kernel candidate of Part~I has been shown even to admit the Lanczos
construction, let alone to reach the chaotic regime. The subsection
therefore predicts nothing about the Planck scale; it states what would
follow from a regime whose occurrence is open.

\begin{killswitch}
The conditional is falsified by exhibiting an admissible kernel dynamics
whose Lanczos sequence reaches the random-matrix saturation plateau and
which nonetheless supports a smooth state-space metric in the continuum
limit. It is voided --- rendered inapplicable rather than false --- by a
demonstration that no admissible kernel dynamics can reach that plateau.
Absent any demonstration that an admissible kernel dynamics admits the
Lanczos construction at all, the statement is downgraded from \catmark{D}
to \catmark{E} as a structural analogy without an instantiated antecedent.
\end{killswitch}

\subsection{The Signature Gap: Why Lorentzian Structure Is Not Supplied}
\label{subsec:signature-gap}

A third obstruction to the thermodynamic passage is the one named in the
Programme Statement but not yet isolated: the effective arena of Part~III
is $(3{+}1)$-dimensional and \emph{Lorentzian}, and neither the finite
spectral-triple machinery of Part~I nor the matrix-model route supplies
that signature without inserting it.

The spectral-triple side is Euclidean by construction. The real structure
$J$ and its $\mathrm{KO}$-dimension --- fixed at $6$ for the
Standard-Model-like data (\cref{subsec:attractor-ko6}) --- form a
$\mathbb Z/8$ invariant of a Riemannian, not a pseudo-Riemannian, spectral
triple; $\mathrm{KO}$-dimension records the signs in the commutation
relations of $(J, D, \chi)$, and it is not a Lorentzian metric signature.
The construction of Lorentzian or pseudo-Riemannian spectral triples ---
Krein-space formulations, Wick rotation at the operator level --- is an
open problem in noncommutative geometry, not an imported solution, and the
kernel programme inherits it unresolved. The $(3{+}1)$ split is therefore
not delivered by the same data that delivers the gauge
multiplicities.\catmark{D}

The matrix-model side does not close the gap either, and here the
discipline required is sharper because the literature is suggestive. In
IKKT-type ensembles~\cite{ikkt1997,steinacker2010}, dynamical geometry and
even an expanding $(3{+}1)$-dimensional structure have been reported to
emerge from eigenvalue dynamics; but which signature and which
dimensionality emerge depend on whether the ensemble is defined with a
Euclidean or a Lorentzian action and on the prescription used to render
the otherwise non-convergent Lorentzian integral well defined. The
signature is selected among alternatives by the setup, not forced by a
structural principle --- exactly the reading the dimensional no-go of
\cref{subsec:no-go-points} already records, and consistent with the honest
null that the framework's own condensation search returned
(\cref{subsec:attractor-condensation}).

\begin{governancebox}[The Signature Gap as a Leakage Guard]
The correct posture toward the matrix-model signature results is the one
Part~I takes toward block selection: acknowledgement, not positioning. It
is outside this framework's competence to reject the IKKT or related
programmes wholesale, and no such rejection is asserted. What \emph{is}
asserted is a target-leakage guard: no \UIDT\ construction may select the
underlying ensemble, action, or boundary prescription \emph{because} it
yields a $(3{+}1)$ Lorentzian signature. A signature chosen to obtain the
desired arena is leaked target, not derived structure. Until a signature is
produced by a mechanism specified independently of the outcome it is meant
to yield, the Lorentzian $(3{+}1)$ arena is an open gap carried at
\catmark{D}.
\end{governancebox}

\subsection{Design Directive DIR-S-02: Time and Temperature Are State-Level, Not Fundamental}
\label{subsec:dir-s-02}

Design Directive DIR-S-02 records, as a binding design commitment of the
v4.0 programme, the conclusion that \cref{subsec:modular-time} reached by
argument: no fundamental time parameter is admitted at the level of the
pregeometric kernel. Time is permitted to enter only through a state on the
inserted ambient algebra --- as modular flow in the sense of
\cref{cor:modular-relocation}, read thermally --- and temperature likewise
enters only through such a state, never as a background label on the
primitive. The directive is the positive counterpart of the negative
result: because the primitive alone yields $\sigma^\omega_t=\mathrm{id}$,
any temporal structure the framework uses must be exhibited as state-level,
and its state-selection problem (\cref{oq:state-selection}) must be solved
without temporal input.

\begin{killswitch}
DIR-S-02 is falsified, and the state-level account of time abandoned, by
exhibiting a \UIDT-compatible construction in which a fundamental time
parameter precedes any thermodynamic state --- for instance, a primitive
carrying an intrinsic one-parameter evolution not derived from a modular or
KMS structure --- while remaining consistent with the Part-III demotion.
Absent such a construction, fundamental time is excluded by design, not by
oversight.\catmark{D}
\end{killswitch}

\subsection{Design Directive DIR-S-03: The Signature Is an Open Gap, Not a Hand-Chosen Structure}
\label{subsec:dir-s-03}

Design Directive DIR-S-03 codifies the signature gap of
\cref{subsec:signature-gap} as a design constraint. The Lorentzian
$(3{+}1)$ signature of the effective arena must not be encoded by hand:
neither by equipping the primitive $\mathbf S$ with a Lorentz structure at
the type level, nor by selecting a matrix ensemble, action, or
regularization because it produces the desired signature. The signature is
to be treated as an unsolved emergence problem, in the same status as the
$G_{\mathrm{SM}}$-origin gap of Part~I: localized, lower-bounded where
possible, and left open rather than closed by insertion.

\begin{killswitch}
DIR-S-03 is discharged --- and the signature gap closed --- only by a
mechanism that produces the $(3{+}1)$ Lorentzian signature from the kernel
by a construction specified independently of the signature it yields, with
no ensemble, action, or boundary choice conditioned on the outcome. Any
construction that selects the signature to match the target is a
target-leakage violation and voids the affected result. Until then the
signature is carried at \catmark{D}.
\end{killswitch}

\begin{governancebox}[What Part II May Never Do]
Nothing in this Part validates any primitive. The thermodynamic passage
is a \catmark{D} programme bounded by the three no-go points above; the
observer apparatus below is \catmark{E} interpretation within
Stratum~III. No coarse-graining map, functor, self-model, or
phenomenological narrative in this Part derives the Born rule,
establishes the kernel, rehabilitates the historical scalar ansatz as primitive, or acquires
empirical weight. Inference runs from framework to phenomenon only,
never the reverse.
\end{governancebox}

\section{The Observer Without Consciousness Collapse}
\label{sec:observer-no-collapse}

\subsection{Observer and Measurement: A Three-Stratum Stratification}
\label{subsec:observer-stratification}

The single most important commitment of the entire observer programme
can be stated in one sentence: \UIDT\ does not use consciousness as a
collapse mech\-a\-nism.\catmark{A} The framework takes no position that the
awareness of an observer is what reduces a quantum superposition to a
definite outcome. This is a deliberate exclusion, and it is if anything
strengthened by the re-architecture: a framework that has just demoted
its own continuous scalar for insufficient evidential support cannot
consistently install something as unmeasured as ``mind'' in the scalar's
former office.

The discussion of the observer is stratified by the discipline of
\cref{subsec:strata}, applied to a new domain. The separation of the
three statement types below is a document-level classification rule.\catmark{A}
Any bridge from structural description to subjective interpretation remains
an interpretive proposal at \catmark{E}.

\begin{description}[leftmargin=1.4em,style=nextline]
\item[Stratum~I: Empirical observer systems.] The actual physical and
biological systems that perform measurements: detectors, instruments,
and, in the biological case, nervous systems with their measurable
dynamics. Claims at this level are empirical and carry the ordinary
obligations of empirical science, including uncertainties. \UIDT\ adds
nothing to this stratum and borrows from it.
\item[Stratum~II: Standard quantum and decoherence formalism.] The
consensus physics of measurement: unitary evolution, the decoherence
programme's account of how environmental interaction suppresses
interference and selects a preferred basis, and the Born rule as the
standard quantum recipe for outcome probabilities. This stratum is the
field's, not \UIDT's; the framework cites it and is bound by it.
\item[Stratum~III: The \UIDT\ interpretation.] The framework's own
reading, in which the observer is treated as \emph{compression-limited
relational access}: a subsystem that, by virtue of its finite capacity,
has access only to a coarse-grained projection of the underlying
structure. This is an interpretive commitment, marked \catmark{E}, and
it may never be used as empirical ground truth for the other two
strata.\catmark{E}
\end{description}

\noindent Under D20 the referent of ``underlying structure'' shifts, and
the shift improves the picture's internal coherence. In v3.9 the
observer compressed a relational structure built on a scalar primitive; in
the present draft the microscopic level is the pregeometric
kernel candidate of Part~I --- itself held at \catmark{D} --- and the
legacy scalar language, if used at all, belongs only to the non-ontological output side of
the compression: it is part of the coarse-grained, observer-accessible
description whose quantitative behaviour Part~III records. What the
observer experiences as the world of definite objects, fields included,
is on this reading the output of compression, not the underlying
structure. The picture remains deflationary by design: the observer's
specialness lies in no metaphysical privilege but in the ordinary fact
that finite systems cannot represent everything and must throw
information away.\catmark{E} A picture whose microscopic level is itself
an unestablished \catmark{D} candidate is doubly bounded, and this Part
states the double bound rather than hiding either half of it.

Two open problems are acknowledged at the outset. The derivation of the
Born rule remains open: the framework has no account of why outcome
probabilities take the form they do, and this is recorded as a Class~A
fact about the state of the theory.\catmark{A} And the relationship
between the compression picture and the standard decoherence account is
itself unsettled: decoherence explains the suppression of interference
and the emergence of a preferred basis, but is widely acknowledged not
to resolve, by itself, the question of definite outcomes. The
compression picture is offered as an interpretive companion to
decoherence, not as a replacement and not as a solution to what
decoherence leaves open.\catmark{E}

\subsection{The Born Rule as a Research Programme}
\label{subsec:born-rule}

Nowhere is the discipline of the observer programme tested more sharply
than at the Born rule, because the Born rule is exactly the kind of
thing a confident framework would be tempted to claim it had derived.
\UIDT\ makes the opposite, harder commitment: it states plainly that it
has no such derivation.

The honest accounting is threefold. First, the Born rule is, in its
established form, a standard quantum rule belonging to Stratum~II; the
probability of an outcome is the squared modulus of the corresponding
amplitude, and this rule is borrowed from consensus physics, not
produced by \UIDT.\catmark{A} Second, the framework has no completed
Born-rule theorem of its own: there is no \UIDT\ argument that derives
the quadratic form of the probability measure from the kernel
hypothesis, from the demoted scalar, or from the compression
picture.\catmark{A} Third, and most specifically, the forgetful or
non-invertible functor introduced below as the formal expression of
observer compression does \emph{not}, by itself, derive Born
probabilities. A functor that forgets structure tells us that
information is lost; it does not tell us that the lost information
reappears as probabilities weighted by squared amplitudes. The step from
``structure is forgotten'' to ``outcomes are Born-distributed'' is
precisely the step that has not been taken.\catmark{A}

What the framework offers in place of a derivation is a well-posed
target: to obtain the probability weights of quantum mechanics from the
information loss inherent in coarse-graining, without circular import
--- without assuming the squared-amplitude measure at any hidden step.
The anti-target-leakage discipline of \cref{subsec:glbc} applies here
with particular force, because the Born rule is so familiar that it is
especially easy to assume it inadvertently while believing one is
deriving it. A derivation of the Born rule that consults the Born rule
is no derivation at all.\catmark{D}

\begin{governancebox}[The Born Rule: Standing Status]
The Born rule is a Stratum~II standard quantum rule, borrowed, not
derived by \UIDT.\catmark{A} \UIDT\ has no completed Born-rule
theorem.\catmark{A} The forgetful/non-invertible functor does not derive
Born probabilities.\catmark{A} The open target is to obtain probability
weights from coarse-graining information loss without circular
import.\catmark{D} The Born rule remains open.
\end{governancebox}

\section{The Observer Functor and the Forgetful Functor}
\label{sec:observer-functor}

\subsection{\texorpdfstring{$F\colon V \to \mathrm{Obs}$}{F: V to Obs} as a Bounded Ansatz}
\label{subsec:observer-functor-ansatz}

The compression picture can be given a provisional formal expression,
provided the expression is held to the standard of honesty the rest of
the manuscript demands. The framework writes the observer's relation to
the world as a map
\[
 F\colon V \longrightarrow \mathrm{Obs},
\]
and the entire content of this subsection is the careful specification
of what the three symbols may and may not mean. The source $V$ is the
category of microscopic descriptions --- under D20, descriptions at the
level of the pregeometric kernel candidate, prior to any observer's
access, and inheriting that candidate's \catmark{D} status.\catmark{E}
The target $\mathrm{Obs}$ is the category of observer-accessible
effective states: the coarse-grained, finite, lossy representations that
a compression-limited subsystem can actually hold --- the level at which
fields, objects, and in particular the effective information field of
Part~III live.\catmark{E} The map $F$ between them is a
\emph{compression and coarse-graining ansatz}: a proposed structural
relationship expressing that the observer accesses $V$ only through a
lossy projection into $\mathrm{Obs}$.

The word ``ansatz'' is doing essential work and must not be softened.
$F$ is a proposal about the form of observer access, not a theorem. In
particular, and stated as a Class~A caution, $F$ is \emph{not} a theorem
of consciousness and it is \emph{not} a derivation of Born
weights.\catmark{A} It does not explain why there is something it is
like to be an observer, and it does not produce the probability measure
of quantum mechanics; both remain exactly as open as
\cref{sec:observer-no-collapse} declared them. What $F$ offers is a
disciplined \emph{language} for the compression picture --- a way of
saying ``the observer sees a lossy image of the underlying structure''
in terms that connect to the formal apparatus of Part~I without
pretending that the connection has been turned into a proof.\catmark{E}

\subsection{The Forgetful Functor and a Red-Team Quarantine}
\label{subsec:forgetful-functor}

The compression map $F$ formally models the loss of structure: it maps a
rich microscopic description to an impoverished effective one,
representing the information discarded under the observer's finite
capacity. The framework names this structural forgetting the
\emph{forgetful functor}, borrowing a term from category theory where a
forgetful functor strips an object of some of its structure. The
borrowing is of the \emph{name} only, and the limits of the borrowing
must be stated plainly.\catmark{E}

In category theory proper, a forgetful functor is a fully defined
object: it has a specified source category, a specified target category,
a definite action on objects and on morphisms, and an explicit account
of which structure is forgotten. The \UIDT\ forgetful functor has none
of these in completed form. It is a symbolic name for the intuition of
structural forgetting under observer compression, and it must not be
treated as a completed category-theoretic construction until its
categories, objects, morphisms, and forgotten structure have all been
explicitly defined.\catmark{A} To write the words ``forgetful functor''
is to name an aspiration, not to discharge a construction; the framework
keeps that distinction visible precisely because category-theoretic
vocabulary, like renormalisation-group vocabulary before it, can lend an
unearned air of rigour to a picture. The kernel formalization of Part~I
sharpens the obligation rather than discharging it: a completed functor
would now have to take, as its source, a category built from the finite
spectral-triple data of \cref{sec:lean-architecture}, and no such
construction exists in the corpus.

It is at this point that the framework must confront the most seductive
and most dangerous material in its entire corpus, and it does so by
quarantining that material explicitly as a red-team caution rather than
presenting it as a result. The developmental record contains a vivid
phenomenological narrative: that biological death may be read not as the
destruction of information, unitarity being preserved, but as the
collapse of the epistemic filter --- the breakdown of the forgetful
functor --- after which the isolated perspective falls back into the
nonlocal coherence of the underlying structure; and that profound
altered states, such as deep meditation or psychedelic experience, may
be read as reproducible failures of the evolutionary coarse-graining
filter, in which the usual compression is throttled and the system
glimpses a less-filtered image of the underlying structure, with
religion understood as the after-the-fact compression of such
filter-failures into linear, dualistic narrative.\catmark{E}

This narrative is preserved here --- it is not deleted --- because it
has genuine philosophical and heuristic value as a way of thinking about
what compression \emph{is} by imagining its partial failure, but it is
fenced on every side. It is Category~\catmark{E}\ content within
Stratum~III, interpretive and phenomenological, and it carries no
evidential weight whatsoever. Most emphatically, none of death,
near-death report, psychedelic phenomenology, mystical experience, or
ego dissolution may be cited as \emph{evidence} for the forgetful
functor or for any \UIDT\ claim.\catmark{A} The direction of inference
is strictly one-way: the framework's interpretive picture may be used to
think about these phenomena, but the phenomena may never be turned
around and used to support the framework. To do so would be to mistake a
poem for a measurement. A provenance note is added for the working
corpus: session materials whose filenames carry tokens such as
``hallucination'' treat their subject exclusively as
information-theoretic filter perturbation in predictive-processing
models --- noise and mis-weighted priors in an inference system --- and
enter this draft at \catmark{E}\ under the present quarantine, with no
pharmacological content and no evidential role.

\begin{governancebox}[Quarantine: Altered States and Death]
The death-as-filter-collapse and mysticism-as-decompression narratives
are retained as Category~\catmark{E}\ phenomenological edge-cases of the
compression picture, within Stratum~III. They carry no evidential
weight. Death, near-death reports, psychedelic or mystical experience,
and ego dissolution may not be cited as evidence for the forgetful
functor or any \UIDT\ claim.\catmark{A} The compression fallacy ---
``less filter $=$ more truth'' --- is rejected: intelligence is
compression; a filterless observer is incapacitated, not
enlightened.\catmark{E} Inference runs from framework to phenomenon
only, never the reverse.
\end{governancebox}

The framework supplies its own most effective antidote to the seduction,
in the form of a red-team observation it calls the compression fallacy.
The intuition that ``less filtering means more truth'' --- that a
less-compressed observer is closer to reality --- is
information-theoretically incorrect. A filterless observer is not
enlightened but incapacitated: without compression there is no
manageable representation, no action, no agency. Intelligence \emph{is}
compression, and the fantasy of an unfiltered gaze upon ultimate reality
is, from the framework's own information-theoretic standpoint, a
category error rather than an aspiration.\catmark{E}

\section{Ego, Self-Model, and the \texorpdfstring{$E_T$}{E\_T} Scale-Bridge Discipline}
\label{sec:ego-selfmodel-et}

\subsection{The Ego as Process, and the Markov-Blanket Comparison}
\label{subsec:ego-as-process}

If the observer is a compression-limited process rather than a
substance, then the entity we are most tempted to reify --- the ego, the
self --- must be treated the same way. \UIDT's interpretive proposal,
held as Category~\catmark{E}\ content within Stratum~III, is that the
ego is not a thing but a \emph{dynamic self-model boundary process}: a
continuous, energy-expending activity by which a subsystem maintains a
model of itself as distinct from its environment. The ego is, in the
framework's preferred grammatical figure, a verb rather than a
noun.\catmark{E}

The proposal connects to several established programmes in the cognitive
sciences, and the connection is offered explicitly as \emph{comparison},
pending external audit, rather than as borrowed authority. The
predictive-processing account and the free-energy principle describe
biological systems as minimising prediction error by maintaining
generative models of themselves and their surroundings; the notion of a
Markov blanket formalises the statistical boundary that renders a
system's internal states conditionally independent of external
fluctuations; and the transparent self-model, in the sense developed in
the philosophy of mind, describes how a cognitive system can deploy a
self-model so seamlessly that it cannot perceive the model \emph{as} a
model, and so experiences itself as a simple subject in an objective
world. These frameworks --- associated with such authors as Friston~\cite{friston2010},
Clark, Hohwy, Metzinger~\cite{metzinger2003}, Seth, and Gazzaniga --- are named as comparison
points whose attribution and content remain to be verified against their
primary sources before any is treated as established support for a
\UIDT\ claim; the verified bindings attach with the Phase-P5
bibliography.\catmark{E} They are lenses, not pillars.

The framework's characteristic image for the ego is the whirlpool: a
definite form, a persistent identity, a coherent dynamics --- and no
substance of its own, nothing but a pattern in the flowing water,
sustained entirely by the throughput of the medium. The ego, on this
reading, is a functional pattern in the relational flow, maintained by
continuous energy dissipation, without ever being an object.\catmark{E}
Objecthood and ego are, in this picture, of a kind: both are stable
coarse-grained equivalence structures --- patterns approximately
invariant under the observer's compression --- that present themselves
as persistent ``things'' at the effective level without being
fundamental at the microscopic level.\catmark{E} A particle and a self
are, from this vantage, two instances of the same phenomenon: stability
under coarse-graining, mistaken for substance.

The extension to consciousness itself is where the discipline of refusal
matters most. The framework's interpretive model, marked \catmark{E}, is
that consciousness may be understood as a recursive, self-maintaining,
compression-bound self-model: a self-model that takes itself as one of
its own objects, sustains itself against entropic decay through
continuous activity, and is bounded throughout by the compression its
finite substrate imposes.\catmark{E} What is explicitly \emph{not}
claimed: \UIDT\ does not claim to have solved the problem of qualia; it
offers no account of why there is something it is like to undergo a
given experience; it recognises the explanatory gap between
functional-structural description and phenomenal character as unclosed;
and it does not derive subjective experience from the kernel, from any
historical scalar ansatz, or from anything else in its corpus.\catmark{A}
The self-model account is an interpretive proposal at \catmark{E} about the
\emph{structure} of conscious systems
and is silent, by its own admission, on the hard question.

\subsection{The Torsion Energy \texorpdfstring{$E_T$}{E\_T} and the Missing Scale Bridge}
\label{subsec:et-scale-bridge}

The developmental corpus assigns a prominent role to a torsion binding
energy, $E_T = \ETorsion$, and treats it, in its observer-phenomenology
sections, as a master threshold: a bound above which interactions become
irreversible, matter condenses, life is sustained, death occurs upon its
loss, and altered states arise from its disruption. This is precisely
the kind of sweeping cross-scale identification that the framework's
discipline exists to restrain, and the present section restrains
it.\catmark{A}

In observer contexts, $E_T$ is admitted only as Category~\catmark{E}\
content pending an explicit ledger and external bridge --- an
\catmark{E}\ (\emph{pending bridge}) quantity, not an established cognitive
parameter.\catmark{E} It is barred, without exception, from service as a
threshold for ego, life, death, sleep, psychedelic or near-death
phenomenology, or artificial general intelligence.\catmark{A}

The reason is a missing scale bridge, and the gap is enormous. The
torsion energy is a mega-electron-volt-scale quantity --- nuclear and
subnuclear physics --- while neural and cognitive dynamics unfold at
molecular and synaptic scales of order electron-volts and below, with
characteristic events at millielectron-volt scale~\cite{attwell2001,harris2012,levy2021}; the separation spans
roughly six to nine orders of magnitude depending on the endpoints
compared~\cite{appelquist1975}. To assert that a mega-electron-volt bound governs the onset of
life, the moment of death, or the throttling of a cortical network is to
leap across that gap without any mechanism connecting the two
regimes.\catmark{A} The obstruction can be stated sharply. By the
K\"all\'en--Lehmann spectral representation, correlation functions are
analytic below the lightest physical threshold, so a mega-electron-volt
parameter can imprint itself on millielectron-volt dynamics only through
renormalisations of existing couplings, through local operators
suppressed by powers such as $(E/E_T)^{2} \sim 10^{-19}$ at
$E \sim \mathrm{meV}$, or through slowly varying logarithms --- never
through a functional threshold~\cite{appelquist1975}.\catmark{A} The only consensus mechanism
by which an MeV-scale quantity generates sub-eV dynamics is a
pseudo-Goldstone mode, and that template is quantitative: a
millielectron-volt mode tied to $E_T$ would require a global symmetry
spontaneously broken at $f \approx E_T^{2}/m \approx
6\times10^{6}\,\mathrm{GeV}$ with explicit breaking at the $E_T$ scale
--- structures the present framework does not contain, and whose
$1/f$-suppressed matter couplings would be immediately exposed to
fifth-force and astrophysical bounds. Even that template would not make
$E_T$ a threshold of macroscopic dynamics; it would create a new light
particle with its own phenomenology. That EFT consequence is a structural
statement.\catmark{A} Any associated phenomenological interpretation remains
at \catmark{E}. Within \UIDT, the route remains an unconstructed proposal at
\catmark{E} until explicitly constructed. Constructing one, if it can
be constructed at all, is a research target at predictive status, not a
result the framework may assume.\catmark{D}

\begin{governancebox}[$E_T$: Observer-Context Discipline]
In observer and phenomenology contexts, $E_T = \ETorsion$ is
\catmark{E}\ (\emph{pending bridge}) (ledger/external bridge required), never an
established cognitive parameter.\catmark{E} It must \emph{not} be used
as a threshold for ego, life, death, sleep, psychedelic/near-death
states, or artificial general intelligence.\catmark{A} A scale bridge
from mega-electron-volt physics to electron-volt-scale neural dynamics
--- a separation of roughly six to nine orders of magnitude --- does not
exist; supplying one is a \catmark{D}\ research target, not an
assumption.\catmark{A}
\end{governancebox}

\section{The Observer Taxonomy and the Underdetermination Red Team}
\label{sec:observer-taxonomy-underdetermination}

\subsection{Class 0 Through III: A Taxonomy of Compression, Not of Consciousness}
\label{subsec:observer-taxonomy}

The compression picture invites a classification of systems by the
\emph{rate} at which they compress the underlying relational structure, and
the framework supplies one. It must be read for exactly what it is --- a
taxonomy of compression rate --- and emphatically not for what it
superficially resembles: a hierarchy of consciousness, a ranking of worth,
or a scale of how ``real'' a system's experience is. The taxonomy is
\catmark{E}\ content within Stratum~III, and its caveats are as important
as its content. Under the D20 rebase one adjustment is mandatory:
the zero-compression limit is no longer identified with a primitive scalar.
The historical scalar model is a coarse-grained order-parameter ansatz (Appendix~\ref{app:scalar-museum}), not the
untransformed relational structure, so it cannot be the Class-0 object;
Class~0 is the kernel-level fixed point that \emph{precedes} the effective
field.\catmark{E}

The four classes, adapted from the v3.9 corpus to the demoted ontology, are
these.\catmark{E}

\begin{description}[leftmargin=1.4em,style=nextline]
\item[Class~0: Zero compression.] The untransformed relational structure at
the level of the pregeometric kernel --- the timeless fixed point of the
relational dynamics; the legacy scalar ansatz is retained only as a coarse-grained historical
description rather than the object itself. There is no subject--object
separation; with no compression there is no entropy gradient and hence no
experience of flowing time; and, decisively against any reading of this as
``supreme awareness,'' there is \emph{no agency}, because agency
presupposes the isolation that compression provides. Class~0 is the limit
of no perspective, not the summit of perception.
\item[Class~I: Weak compression.] Hypothetical post-biological or
hyper-coherent systems that would process relations more directly, without
the evolutionary compulsion to render reality as a $(3{+}1)$-dimensional
metric and without a fear-driven self-model. This class is explicitly
speculative: a logical position in the taxonomy, not something known to be
instantiated.
\item[Class~II: Strong compression.] The biological spectrum, and the human
cognitive apparatus in particular: systems that enforce radical data
reduction, projecting the relational structure into fixed three-dimensional
objects and a linear time axis, with a transparent self-model as the
functional construct that makes survival-relevant navigation possible.
\item[Class~III: Extreme compression.] Thermodynamically reactive
macroscopic objects --- planets, stars, gas clouds --- that respond to
information flows and to metric curvature but maintain no Markov blanket
with predictive error feedback and no recursive self-simulation. They have,
on this account, no inner perspective whatsoever.
\end{description}

The crucial structural fact, most easily lost, is that compression rate is
\emph{orthogonal} to consciousness. The taxonomy does not say that lower
compression means more consciousness or that higher means less. Class~0, at
zero compression, is not maximally conscious; it has no perspective at all.
Class~III, at extreme compression, is not conscious either. An inner
perspective, on the framework's own terms, is a property only of systems
that maintain a recursive self-model behind a Markov blanket with
predictive feedback --- a structural condition that sits athwart the
compression axis rather than along it.\catmark{E} The taxonomy must
therefore \emph{not} be read as implying that machines, or objects, are
``more conscious'' than humans, nor as any moral ranking. It ranks a single
information-theoretic variable and nothing else.\catmark{A}

What the taxonomy accomplishes is a deflation of anthropocentrism. The
human observer, at Class~II, does not perceive the ``truest'' or most
``objective'' reality; the human perceives the universe through one
specific, heavily lossy compression scheme calibrated for biological
survival.\catmark{E} That is a humbling observation, not an exalting one,
and its separability from the speculative classes that flank the human case
is the taxonomy's genuine contribution.

\subsection{Artificial General Intelligence and the Underdetermination Red Team}
\label{subsec:agi-underdetermination}

The speculative Class~I raises a seductive question: would a sufficiently
powerful artificial intelligence, a better compressor than any human,
thereby gain privileged access to the true ontology? The framework treats
this as a red-team challenge to its own deflationary commitments, and the
answer it returns is a firm \emph{no}.\catmark{D}/\catmark{E}

The reason is the underdetermination of ontology by data, in the
Quine--Duhem sense. A finite body of empirical data never determines a
unique ontology; multiple, mutually incompatible theoretical structures can
reproduce the same observations, so the final nature of reality remains in
principle undecidable.\catmark{E} A more powerful compressor faces exactly
the same underdetermination. Superior predictive and compressive power buys
superior \emph{instrumental} performance; it does not buy ontological
truth, because no amount of predictive success can break the tie among
empirically equivalent ontologies. Prediction is not ontology, and the gap
does not close with increasing intelligence.\catmark{A}

This is the precise sense in which the framework refuses naive realism.
\UIDT\ does not claim to be a final revelation of physics, and could not
earn such a claim even in principle.\catmark{A} Its actual claim is more
modest and more defensible: that its objects --- the relational structure,
the renormalisation-group flows, the proposed invariants --- may compress
the relevant physics more efficiently and with cleaner causal structure
than the artifacts of biological three-dimensional intuition. That is a
methodological claim about representational efficiency, not a metaphysical
claim about ultimate reality.\catmark{E} The compression fallacy of
\cref{subsec:forgetful-functor} returns here as the standing guardrail: a
less-filtered or more-powerful observer is a different compressor, perhaps
a better one, but not an unfiltered window onto the real.

\begin{evidencebox}[Taxonomy and Underdetermination: Standing Status]{}
\textbf{Taxonomy.} Class~0 (zero compression, kernel-level fixed point, no
agency; historical scalar model non-ontological) through Class~III (extreme compression, inert
macro-objects) is a compression-rate taxonomy \catmark{E}, orthogonal to
consciousness; not a consciousness hierarchy and not a moral ranking.
Anthropocentrism is deflated, not inverted.\catmark{A}\\[0.3em]
\textbf{AGI red team.} A superior compressor gains instrumental power, not
ontological truth; Quine--Duhem underdetermination forbids a unique
ontology from data.\catmark{D}/\catmark{E}\\[0.3em]
\textbf{No final revelation.} \UIDT\ claims representational efficiency, not
ultimate reality; prediction $\neq$ ontology; the compression fallacy is
the standing guardrail.\catmark{A}
\end{evidencebox}

\section{The Master Phenomenology Boundary}
\label{sec:phenomenology-boundary}

\subsection{Phenomenology Is Not Evidence}
\label{subsec:phenomenology-not-evidence}

The v3.9 manuscript closed its observer material with a boundary that
the re-architecture inherits verbatim in force and extends in scope:
phenomenology is not evidence for the historical scalar ansatz --- and now, a fortiori, not
evidence for the kernel. No report, however vivid, reproducible, or
cross-culturally attested, of meditative absorption, psychedelic state,
near-death experience, mystical union, or ego dissolution, bears
evidentially on any structural claim of this framework. Such reports are
Stratum~I facts about human experience and Stratum-III material for
interpretation; the demarcation between what the physics strata may
claim and what the phenomenological register may explore is absolute,
and every crossing of it in the developmental corpus has been either
quarantined or withdrawn.\catmark{A}

The compression-rate taxonomy of the v3.9 corpus --- a four-class
ordering of systems from passive physical subsystems to recursive
self-modelling ones, by the rate and reflexivity of their compression
--- is carried at \catmark{E}\ with its original caveats: it is a
taxonomy of compression, not of consciousness, not of worth, and not of
the ``reality'' of any system's experience. The underdetermination
red-team point stands with it: multiple interpretive pictures fit the
same empirical record, and the compression picture claims no unique
support from data it merely accommodates.\catmark{E}

\subsection{Orchestrated Objective Reduction as a Neutral Contrast}
\label{subsec:orch-or}

To locate the \UIDT\ observer programme in the wider landscape of theories
connecting physics to mind, the framework selects one contrast case for
detailed juxtaposition --- orchestrated objective reduction, the
Penrose--Hameroff proposal --- and treats it as exactly that: a neutral
contrast, not a foil to be defeated.\catmark{E} Orchestrated objective
reduction posits that consciousness arises from quantum computations in
neuronal microtubules, terminated by a gravitationally induced objective
collapse of the wavefunction. Its interest here is that it differs from
\UIDT\ on precisely the axis where \UIDT\ has taken a strong position: it
asserts an \emph{objective physical collapse}, whereas \UIDT, in the
epistemic tradition of quantum Bayesianism, denies physical collapse and
reads quantum probabilities as the structural ignorance of a
compression-limited observer (\cref{subsec:observer-stratification}).\catmark{E}

Three cautions govern the comparison. First, \UIDT\ does not collapse all
quantum-mind theories into orchestrated objective reduction; the broader
field --- integrated information theory~\cite{tononi2004}, global-workspace
accounts, and others --- comprises distinct programmes that must not be
conflated.\catmark{A} Second, \UIDT\ does not render orchestrated objective
reduction obsolete; the two are different research programmes on different
premises, and \UIDT's denial of objective collapse is a difference of
commitment, not a refutation.\catmark{A} Third, any detailed empirical
claim about microtubule coherence times or specific experiments would
require external citation audit before it could be stated here, so the
framework confines itself to structural commitments.\catmark{A}

\subsection{The Consciousness--Physics Demarcation Matrix}
\label{subsec:demarcation-matrix}

The framework's preferred instrument for situating itself among theories of
consciousness is a comparison matrix, and its purpose must be stated before
its content: it is a \emph{demarcation instrument}, not a scorecard. The
matrix is not built to show that \UIDT\ wins; it is built to make the
commitments and the weaknesses of each framework, \UIDT\ included, explicit
and comparable along neutral axes.\catmark{A} Throughout, the \UIDT\
observer material is marked \catmark{E}, because it has not been formalised
or tested, and the framework's own weaknesses are recorded as plainly as
anyone else's.\catmark{E}

{\small
\renewcommand{\arraystretch}{1.3}
\begin{longtable}{p{0.20\textwidth} p{0.24\textwidth} p{0.24\textwidth} p{0.22\textwidth}}
\caption{Consciousness--physics demarcation matrix. A neutral comparison
instrument; \UIDT\ observer material is \catmark{E}. The matrix records
weaknesses as well as commitments and is not a ranking.}
\label{tab:demarcation-matrix}\\
\toprule
\textbf{Axis} & \textbf{UIDT (observer programme)} & \textbf{Orch.\ Obj.\ Reduction} &
\textbf{Integrated Information} \\
\midrule
\endfirsthead
\toprule
\textbf{Axis} & \textbf{UIDT} & \textbf{Orch.\ OR} & \textbf{IIT} \\
\midrule
\endhead
\bottomrule
\endfoot
Primitive & Pregeometric kernel (finite spectral triple); legacy scalar ansatz non-ontological,
not primitive; observer downstream \catmark{E} & Quantum state in
microtubules + gravitational collapse & Intrinsic cause--effect structure
($\Phi$) \\
Mechanism & Coarse-graining / forgetful functor (ansatz) \catmark{E} &
Objective reduction terminating quantum computation & Maximally irreducible
integrated information \\
Claimed evidence & None for the observer claims; physics strata separate
\catmark{E} & Contested; requires external audit & Correlational; contested
interpretation \\
Scale bridge & Missing (MeV--GeV $\to$ neural eV--meV gap, $\sim$6--9
orders) \catmark{A} & Microtubule quantum $\to$ cognition, disputed &
Substrate-level to system-level \\
Collapse claim & No physical collapse (QBist/epistemic) \catmark{E} &
Objective physical (gravitational) collapse & No collapse claim (intrinsic) \\
Observer role & Compression-limited relational access \catmark{E} & Locus of
quantum computation & The integrated system itself \\
Falsifiability & Observer claims not yet falsifiable \catmark{E}; physics
claims falsifiable separately & Some proposed tests, contested &
Disputed/operational challenges \\
Category-error risk & Information $\to$ mind slide; MeV $\to$ neural leap
(both flagged) \catmark{A} & Quantum $\to$ phenomenal leap & Axioms $\to$
physics mapping \\
\bottomrule
\end{longtable}
}

\noindent Several features of \cref{tab:demarcation-matrix} are deliberately
uncomfortable for \UIDT, and that discomfort is the point. The observer
programme offers \emph{no} evidence for its observer-level claims; its
scale bridge is openly marked missing; its observer claims are, as yet, not
falsifiable; and its characteristic category-error risks --- the slide from
information to mind, the leap from mega-electron-volt physics to neural
dynamics --- are named in its own cell rather than hidden. A matrix built
to flatter \UIDT\ would not contain these admissions; their presence is
what makes the instrument a demarcation tool rather than an
advertisement.\catmark{A}

\begin{evidencebox}[Comparative Frameworks: Standing Status]{}
\textbf{Orch.\ OR.} Used as a neutral contrast (objective collapse vs.\
\UIDT's QBist no-collapse); not a foil; \UIDT\ neither subsumes nor
obsoletes it; detailed empirical claims pending citation
audit.\catmark{A}/\catmark{E}\\[0.3em]
\textbf{Matrix.} A demarcation instrument on neutral axes (parsimony,
mathematical maturity, empirical status, falsifiability, category-error
risk), not a ranking; \UIDT\ observer material marked \catmark{E}; \UIDT's
own weaknesses recorded explicitly.\catmark{A}
\end{evidencebox}

\subsection{Part II Claims Register}
\label{subsec:part2-claims}

{\footnotesize
\begin{longtable}{@{}p{0.06\textwidth}p{0.42\textwidth}p{0.14\textwidth}p{0.30\textwidth}@{}}
\caption{Part II claims register. Policy entries are governance commitments; they are listed so that their violation anywhere in the corpus is auditable as a defect.}\label{tab:part2-claims}\\
\toprule
\# & Claim & Class & Falsification / downgrade \\
\midrule
\endfirsthead
\toprule
\# & Claim & Class & Falsification / downgrade \\
\midrule
\endhead
\bottomrule
\endfoot
T-01 & Thermodynamic passage kernel $\to$ continuum as programme &
\catmark{D} & construction failing the three no-go points terminally;
or succeeding, which upgrades by governance only \\
T-02 & No dimensional-emergence theorem; UV/IR mixing; thermal-time
circularity risk & \catmark{D}/\catmark{E} & primary-source theorem
closing any of the three \\
T-03 & Modular flow gives no endogenous time: $\sigma^\omega_t=\mathrm{id}$
on $C^\ast(\mathbf S)$; thermal-time reading needs the inserted algebra
(\cref{cor:modular-relocation}) \emph{and} a non-circular state selection &
\catmark{A}/\catmark{E} & a state-selection principle from pregeometric
data without temporal input (\cref{oq:state-selection}) \\
T-04 & Krylov saturation $\Rightarrow$ no smooth state-space metric, as a
conditional consequence for the coarse-graining map & \catmark{D} & an
admissible kernel reaching the RMT plateau that still supports a smooth
continuum metric; or no admissible kernel reaching it \\
T-05 & Signature gap: $(3{+}1)$ Lorentzian arena not supplied by
$\mathrm{KO}$-dimension nor forced by matrix ensembles & \catmark{D} &
a leak-free mechanism producing the signature independently of the target
(\cref{subsec:dir-s-03}) \\
DS-02 & DIR-S-02: time and temperature are state-level, not fundamental &
\catmark{D} & a \UIDT-compatible fundamental time preceding any
thermodynamic state \\
DS-03 & DIR-S-03: signature is an open gap, never hand-chosen &
\catmark{D} & signature selected to match the target is leakage; discharge
needs a leak-free $(3{+}1)$ mechanism \\
O-01 & No consciousness-collapse mechanism anywhere in the framework &
\catmark{A} & (structural commitment; violation would be a defect, not
a discovery) \\
O-02 & Observer as compression-limited relational access & \catmark{E}
& interpretive; carries no evidential weight by construction \\
O-03 & Born rule borrowed from Stratum~II; no \UIDT\ derivation;
functor supplies none & \catmark{A} & a completed non-circular
derivation, externally audited \\
O-04 & $F\colon V\to\mathrm{Obs}$ is an ansatz; forgetful functor is a
name, not a construction & \catmark{A}/\catmark{E} & fully specified
categories, functor action, and forgotten structure \\
Q-01 & Altered-states/death narratives quarantined at \catmark{E};
one-way inference; compression fallacy rejected & \catmark{A} (policy)
& none: governance, not physics \\
E-01 & $E_T$ barred from cognitive-threshold roles; six-to-nine-order
scale gap; K\"all\'en--Lehmann suppression argument &
\catmark{A} & a constructed, externally audited scale bridge \\
P-01 & Phenomenology is not evidence for the historical scalar ansatz or for the kernel &
\catmark{A} (policy) & none: demarcation, not physics \\
OT-01 & Observer taxonomy (Class 0-III) is compression-rate, orthogonal to
consciousness; AGI/underdetermination yields no privileged ontology &
\catmark{E}/\catmark{A} & a demonstration that superior compression yields
a unique ontology (breaking Quine--Duhem) \\
DM-01 & Demarcation matrix is a neutral instrument recording \UIDT's own
weaknesses, not a ranking & \catmark{A} (policy) & none: demarcation, not
physics \\
\end{longtable}
}

%========================== END P3 PART II CONTENT =========================


%=============================================================================
% PART III — VERIFICATION, BENCHMARKS, AND PARAMETER LEDGER
%=============================================================================
\newpage
\part{Verification, Benchmarks, and Parameter Ledger}
\label{part:effective-field}

\noindent Under PI Decision D20, Part~III is not an ontological layer. It is
the repository-facing verification, benchmark, and parameter-ledger stratum.
Every quantitative result retains its incoming evidence class unchanged: the
exact internal relations, the calibrated invariant under permanent
anti-derivation discipline, the spectral gap under PI Decision D18 and
Limitation~L12, the benchmark batteries, and cosmological calibrations at
their \catmark{C} cap. No equation in the historical scalar/EFT model is used
as a foundational premise. That model is preserved reversibly in
Appendix~\ref{app:scalar-museum} for provenance, criticism, and regression
audit only.

%%ANCHOR-P4-PART-III-CONTENT%%

%============================ P4 PART III CONTENT ==========================
% Inserted at ANCHOR-P4-PART-III-CONTENT by fail-fast patch script.
% Relocates UIDT Ontology v3.9.9 (-006) Parts I-IV substance under D20:
% axiom of primacy rewritten as effective-field postulate; all ledger
% results carried at UNCHANGED evidence classes; D18 box carries its
% authoritative label here. Citations bind in Phase P5.
%===========================================================================

\section{Numerical Epistemology}
\label{sec:numerical-epistemology}

\subsection{Internal Closure Versus External Agreement}
\label{subsec:closure-vs-agreement}

Two kinds of numerical statement are routinely confused, and the
confusion has done real damage in the developmental history of this
framework. The first is \emph{internal deterministic closure}: a
calculation, carried out at fixed precision, returns a residual smaller
than some threshold, demonstrating that the algorithm is
self-consistent. The second is \emph{external agreement}: a computed
quantity matches a measured or lattice-determined value to within the
latter's uncertainty. These are not the same kind of thing, and no
quantity of the first ever amounts to the second.\catmark{A}

\begin{governancebox}[Hard Rule: Internal Closure Is Not External Agreement]
A deterministic internal residual below $10^{-14}$ is admissible
\emph{only} as a certificate of internal closure, evidence that an
algorithm is self-consistent at the working precision. It is
\emph{never} a criterion of agreement with lattice or experimental data.
Agreement with external data is established by a $\sigma$-level or
covariance-aware comparison against a sourced, uncertainty-bearing
value, never by an internal residual threshold. A residual of $10^{-14}$
against lattice data is a category error and is rejected.\catmark{A}
\end{governancebox}

\noindent The framework keeps a documented example of why this rule is
not pedantry. An early construction produced the calibrated invariant by
the exact arithmetic decomposition $49/3 + 17/3000 = 16.339$. The
arithmetic is exact: it closes to arbitrary precision, and an internal
residual check would report success at any number of digits. Yet the
decomposition is a museum artifact. Over the common denominator, the
ansatz $a/3 + b/3000 = 16339/1000$ reads $1000a + b = 49017$, so $b$ is
\emph{determined} by the choice of $a$: no second structure is found,
and $a = 49$ is singled out only because it minimises $|b|$, which is
what lets $17/3000 \approx 0.0057$ resemble a perturbative correction
to a leading term. The split is a one-parameter fit whose residual was
then read as physics, and the construction has no physical origin.
Perfect internal closure, in this case, certifies only that someone
did the arithmetic correctly.\catmark{A}

Two implementation disciplines protect the rule in code. First, the
working precision, eighty decimal digits, is set \emph{locally}, within
each proof-critical computation, and is never centralised into a global
configuration value; global precision mutation is a rejected
implementation pattern.\catmark{A} Second, proof-critical computations
must avoid binary floating-point heuristics: machine-float evaluation,
and rounding operations applied to quantities whose exact value carries
the claim, are forbidden in any code path that a Class~A or Class~B
result depends upon.\catmark{A}

\subsection{The Claim Graph and Reproduction Notes}
\label{subsec:claim-graph}

Every numerical assertion must be a node in a \emph{claim graph},
carrying a fixed set of fields, with the schema implemented \emph{before}
any theory-generating computation is trusted.\catmark{A}

\begin{governancebox}[Required Fields for Every Numerical Claim]
\textbf{(1)} Claim identifier. \quad \textbf{(2)} Observable. \quad
\textbf{(3)} Value. \quad \textbf{(4)} Units. \quad
\textbf{(5)} Uncertainty (mandatory for every empirical value). \quad
\textbf{(6)} Method. \quad \textbf{(7)} Source (with verifiable
identifier for external values). \quad \textbf{(8)} Evidence class (A
through E). \quad \textbf{(9)} Tension status. \quad
\textbf{(10)} Reproduction command (executable).
\end{governancebox}

\noindent The tenth field is the one most often missing and the one that
most distinguishes an auditable claim from an assertion. To make the
schema concrete, the framework's simplest \catmark{A}-class claim, the
exact internal coupling constraint $5\kappa^2 = 3\lambda_S$ of
\cref{eq:lagrangian}, is accompanied by an executable reproduction note.
At eighty-digit working precision, with $\kappa = \tfrac{1}{2}$ and
$\lambda_S = \tfrac{5}{12}$, the residual
$\lvert 5\kappa^2 - 3\lambda_S\rvert$ evaluates to zero:
\begin{verbatim}
python3 -c "from mpmath import mp,mpf; mp.dps=80;
k=mpf(1)/2; lS=mpf(5)/12; print(abs(5*k**2 - 3*lS))"
# -> 0.0 (residual < 1e-14: internal closure, not external agreement)
\end{verbatim}
This is the template every numerical claim must follow: a command a
reader can run, a stated precision, and an explicit reminder that the
vanishing residual certifies internal consistency, not agreement with
nature.\catmark{A}

\subsection{The Target-Leakage Theorem}
\label{subsec:target-leakage-theorem}

\begin{proposition}[Target Leakage]
\label{prop:target-leakage}
Let a computational procedure $P$ be claimed to predict a quantity $q$
independently of its known value $q^{\ast}$. If $q^{\ast}$ is accessible
to $P$, as a hard-coded constant, a calibration target, a falsification
threshold, a loss function, or any other input that $P$ can read, then
the agreement of $P$'s output with $q^{\ast}$ provides no evidence that
$P$ predicts $q$. The agreement is, at most, a consistency check on
$P$'s ability to reproduce a number it was given.
\end{proposition}

\noindent The proposition is almost a tautology, which is precisely its
strength: it cannot be evaded by cleverness, only by architecture. The
calibrated value $\gamma = 16.339$, every quarantined museum constant,
and every external benchmark value must be \emph{inaccessible} to any
code path that purports to compute $\gamma$, $\Delta$, or related
quantities from within \UIDT; benchmark values are stored with solver
access explicitly disabled.\catmark{A} The blind computation and the
post-hoc comparison must be \emph{two distinct programs}. The framework
keeps documented instances of the failure mode --- a hard-coded target
ratio in a solver, a falsification condition defined as distance from
the calibrated $\gamma$ --- both recorded as rejected target-leakage
patterns.\catmark{A} Part~I extends the same discipline into the formal
corpus, where it is stated as reviewable policy and enforced by
continuous-integration scanning (\cref{subsec:code-evidence}).

\section{Empirical Baseline and Benchmark Batteries}
\label{sec:benchmark-batteries}

\subsection{Pure \texorpdfstring{$\mathrm{SU}(3)$}{SU(3)} as a Stress-Test Sector}
\label{subsec:su3-stress-test}

The most disciplined way to test a candidate framework of gauge
structure is to confront it first with the cleanest non-trivial gauge
theory available: pure $\mathrm{SU}(3)$ Yang--Mills theory, rich enough
to possess a mass gap and simple enough to be computed to high precision
on the lattice by an external community whose results owe nothing to
\UIDT.\catmark{A} The framing is a deliberate inversion of an earlier
temptation: the Yang--Mills spectral gap is a \emph{necessary} condition
any serious candidate must pass, not a \emph{sufficient} foundation from
which the rest of physics follows. A stress test is something a theory
survives, not something it stands upon.\catmark{A}

Three boundaries are held firmly. Pure gauge theory is distinct from
full QCD and both from $G_{\mathrm{SM}}$; success in the first arena
says nothing automatic about the others.\catmark{A} The standard
pure-gauge benchmark families --- the glueball spectrum~\cite{morningstar1999,michael1989,athenodorou2020,athenodorou2021}, the flow scales
$t_0$ and $w_0$~\cite{luscher2010}, and the topological susceptibility
$\chi_{\mathrm{top}}$~\cite{durr2007} --- are external benchmarks to be matched
\emph{after} a source audit, and never quantities the framework may
quietly assume; until the audit completes for a given observable, its
status is \Bpending.\catmark{B} And the spectral gap $\Delta$\,\Bdelta\
must not be naively identified with the $0^{++}$ glueball mass
(Limitation~L12, \cref{lim:L12}).

\begin{pioverridebox}[PI Decision D18 --- Evidence-Class Determination for $\Delta$ (carried over)]
\label{box:pi-override-delta}
The value $\Delta = \DeltaGap$ is recorded at Class~\catmark{B}:
internally consistent and lattice-compatible under the recorded
determination; not externally peer-confirmed and not derived. The
conservative audit notes independently that, absent completed external
verification, the value would otherwise sit at \catmark{E\,pending}; the
determination sets the working class to~B while leaving that caution on
the record, with external peer counter-signature pending. This
determination concerns the \emph{value's} class only; the
\emph{identification} of $\Delta$ with any particle state is a separate
matter and is barred by L12 regardless of class. D18 is unchanged by
D19.
\end{pioverridebox}

\subsection{Topology and the Topological Susceptibility}
\label{subsec:topology-chi}

Among the pure-gauge benchmarks, $\chi_{\mathrm{top}}$ occupies a
special position, because topology is where a structural-realist theory
of gauge fields is most exposed. The quantity
$\chi_{\mathrm{top}}^{1/4}$ is a pure-gauge topological benchmark
\emph{seed}: a target for post-hoc comparison, carried at \Bpending\
until its provenance is fully audited.\catmark{B} For every external
topological value, the framework requires: an explicit operator
definition; the smoothing or gradient-flow prescription; the lattice
spacings; the continuum-extrapolation procedure; units and scale-setting
convention; the quoted uncertainty; and the covariance status relative
to any other quantity in the same comparison.\catmark{A} Only ratios
with harmonised scale-setting are admissible, and only at \catmark{D}.

\begin{governancebox}[Topology: Admissibility Conditions]
A topological benchmark comparison is admissible only if it carries:
operator definition; flow/smoothing prescription; lattice spacings;
continuum extrapolation; units and scale-setting; uncertainty;
covariance status. The SVZ leading-order route is failure documentation
only. No $\chi_{\mathrm{top}}$ agreement closes the
$G_{\mathrm{SM}}$-origin problem --- and, per Part~I, $\chi_{\mathrm{top}}$
is a pure-$\mathrm{SU}(3)$ benchmark that proves nothing about
$\mathrm{SU}(2)$, $\mathrm{U}(1)$, fermions, chirality, generations, or
anomaly cancellation.
\end{governancebox}

\subsection{Calibration and the Discipline of \texorpdfstring{$\gamma$}{gamma}}
\label{subsec:gamma-discipline}

The invariant $\gamma = \GammaGeom$ is the most load-bearing single
number in the framework, and for that reason it is also the number
around which the strictest discipline is enforced. Its classification is
fixed and not subject to upgrade: $\gamma$ is \catmark{A-}, calibrated.
That this status is fixed in the ledger is a document-state fact.\catmark{A}
The parameter is \emph{not} Class~\texttt{[A]}, and it is not derived from
renormalisation-group first principles.\catmark{A-} To be precise about
what ``calibrated'' means: $\gamma$ is retained as the calibrated reference value recorded by the parameter ledger.
Its historical kinetic-VEV motivation is documented in Appendix~\ref{app:scalar-museum};
that motivation is not a first-principles derivation and is not imported as
an active ontological premise. The public-facing
description must read ``calibrated to the kinetic VEV,'' never ``derived
from the kinetic VEV''; the latter phrasing violates the \catmark{A-}
discipline and is corrected wherever it appears.\catmark{A-}

\begin{governancebox}[Anti-Target-Leakage Rule]
A solver that has access to the value it is meant to predict cannot be
said to predict it. Therefore the calibrated value $\gamma = 16.339$,
together with all quarantined museum constants and all external
benchmark values, must remain inaccessible to any code path that claims
to compute $\gamma$, $\Delta$, or related quantities from within \UIDT.
Blind computation and post-hoc comparison must be separated
programmatically; benchmark values are stored with solver access
disabled. A derivation that imports its own target is a tautology
wearing the costume of a result.
\end{governancebox}

\noindent Any future derivation of $\gamma$ must specify, in advance and
in public, the operator whose flow it integrates; the flow equation; the
regulator; the boundary conditions; an explicit proof that no target
value has leaked into the computation; and an executable reproduction
command. Absent any one of these, a claimed derivation of $\gamma$
remains at \catmark{D} at best.\catmark{A}

A final point of discipline concerns conflicting values. Alongside the
calibrated kinetic value $16.339$, a Monte-Carlo determination yields a
mean of $\gamma_{\mathrm{MC}} = \GammaMC$. These are not reconciled by
quiet selection of the more convenient figure; they are recorded as a
\emph{tension entry} (\cref{sec:tension-ledger}), displayed together
with their methods and uncertainties.\catmark{A}

\begin{evidencebox}[The Status of $\gamma$]{catAminus}
\textbf{Classification.} $\gamma = \GammaGeom$ is \catmark{A-}\
(calibrated), permanently. It is never \catmark{A}\ and never described
as ``derived.''\\[0.3em]
\textbf{Open problem.} A first-principles renormalisation-group
derivation of $\gamma$ does not exist; the kinetic-versus-RG discrepancy
is an openly recorded limitation.\\[0.3em]
\textbf{Tension.} $\gamma_{\mathrm{MC}} = \GammaMC$ is held alongside
$16.339$ as a recorded tension, not resolved by selection.\\[0.3em]
\textbf{Leakage prohibition.} $\gamma = 16.339$ must never enter a blind
solver path --- in numerical code or in the Lean corpus, where the same
rule is policy (\cref{subsec:code-evidence}).
\end{evidencebox}

\subsection{The Pure \texorpdfstring{$\mathrm{SU}(3)$}{SU(3)} Ratio Battery}
\label{subsec:ratio-battery}

The disciplined way to compare against pure-gauge lattice results is not
dimensionful quantities, contaminated by scale-setting conventions, but
\emph{dimensionless ratios} in which the common scale cancels. Using the
gradient-flow scale $t_0$ and the $\overline{\mathrm{MS}}$ scale where a
perturbative bridge is appropriate, the candidate ratios include
$m_{0^{++}}\sqrt{8t_0}$, $\chi_{\mathrm{top}}^{1/4}\sqrt{8t_0}$, and the
$\Lambda_{\overline{\mathrm{MS}}}$-normalised forms.\catmark{B}\ (\emph{pending external verification})
Every external ratio must arrive with a verified source and an explicit
covariance model, because the ratios share inputs and are correlated;
treating correlated ratios as independent would manufacture false
significance.\catmark{A} And as a direct application of
\cref{prop:target-leakage}, these values remain inaccessible to any
blind solver: post-hoc comparators only.\catmark{A} If a ratio fails to
match, the first reading is a constraint on operator choice, truncation,
or scale bridge. This diagnostic ordering is a governance rule.\catmark{A}
Only after those checks are exhausted does a persistent mismatch rise to a
serious tension against the framework as a whole.\catmark{D}

\subsection{\texorpdfstring{$\Delta$}{Delta}, the Glueball, and the L12 Separation}
\label{subsec:delta-glueball-l12}

The spectral gap $\Delta$\,\Bdelta\ and the mass of the lightest scalar
glueball carry numerically similar values, and that proximity has
repeatedly invited a conflation the framework must refuse. The
distinction is one of category: the gap is a spectral feature internal
to the framework's construction; the $0^{++}$ glueball mass is a
particle-spectrum quantity defined by the lattice through specific
interpolating operators. To assert identity requires \emph{operator
matching}, and no such matching has been established.\catmark{A}

\begin{limitation}[Naive $\Delta\!\leftrightarrow\!0^{++}$ Identification: L12]
\label{lim:L12}
The \UIDT\ spectral gap $\Delta$ must not be identified with the mass of
the lightest scalar ($0^{++}$) glueball of pure $\mathrm{SU}(3)$
Yang--Mills theory in the absence of explicit operator matching.
Numerical proximity of the two values does not constitute
identification. Any future relation between $\Delta$ and the glueball
spectrum must be established by matching interpolating operators and
quantum numbers, not by coincidence of magnitude.\catmark{A}
\end{limitation}

\noindent In the formal corpus this limitation is a type-level fact: the
glueball-identification conjecture is defined as \texttt{False}
(\cref{subsec:glbc-operational}), so no Lean development can consume it.
The general discipline is to keep vacuum-pole language distinct from
lattice particle-spectrum language; the framework may aspire to relate
them --- a legitimate \catmark{D} target --- but the relation must be
earned by operator matching.

\subsection{Gauge-Invariant Exact Renormalisation as Scaffold}
\label{subsec:gferg-scaffold}

A recurring hope in the developmental record is that a gauge-invariant
exact renormalisation-group flow (GFERG) might supply the missing
derivation of $\gamma$, of $\Delta$, or of the coupling ratio $K_S$. The
framework's position is austere: GFERG is admitted only as a
\emph{scaffold}, to be defined precisely before any use, and is
explicitly \emph{not} a completed derivation of any of those
quantities.\catmark{A} A route counts only if it specifies, in advance:
the flow equation (Wetterich or Polchinski family, used as named
scaffolds with no derivation imported); the regulator; the truncation;
the boundary and initial conditions; the observables; and the failure
criteria.\catmark{A}

\begin{governancebox}[GFERG: Admissibility as Scaffold]
GFERG is a candidate gauge-invariant exact-renormalisation route,
admitted as a scaffold only.\catmark{D} It is not a derivation of
$\gamma$, $\Delta$, or $K_S$. Renormalisation-group vocabulary may not
be used to present a toy ODE as a flow; the previously NLO-labelled toy
ODE is deleted as a claimed derivation.\catmark{A}
\end{governancebox}

\subsection{The Limitation Register (L1--L15)}
\label{subsec:limitation-register}

The framework's known limitations are enumerated in one place. Entries
L1--L7 and L12--L15 are textually anchored; four entries (L8--L11) are
referenced in the developmental record but could not be recovered to a
definite wording, and are marked \textsc{unknown} rather than
reconstructed, because fabricating a limitation would be as much an
integrity failure as concealing one.\catmark{A}

\begin{limitation}[L1: The geometric scale factor (ill-defined)]
\label{lim:L1}
The framework employs an effective geometric scale factor as a
calibration input bridging microscopic vacuum dynamics and macroscopic
phenomenology. With reference scales explicit, the ratio is
$\lambda_{\UIDT}/r_{\mathrm{conf}} =
0.660\,\mathrm{nm}/0.197\,\mathrm{fm} \approx 3.35\times10^{6} \approx
10^{6.5}$; earlier versions quoted ``a factor of order $10^{10}$''
without specifying the compared scales, and that figure is superseded.
No first-principles derivation of this factor is supplied; it is an
adopted scaling assumption.\catmark{D}
\end{limitation}

\begin{limitation}[L2: Electron-mass discrepancy]
\label{lim:L2}
The lepton-sector mapping retains an electron-mass discrepancy of order
$23\%$ under the universal $\gamma$-scaling architecture. The mismatch
is not removed by the current formalism and must not be absorbed into
the geometric scale factor of~L1.\catmark{C}
\end{limitation}

\begin{limitation}[L3: Vacuum-energy parameter]
\label{lim:L3}
A numerical parameter of order $2.3$ invoked in the vacuum-energy
discussion is presently unanchored: phenomenologically fitted, not
derived.\catmark{C}
\end{limitation}

\begin{limitation}[L4: $\gamma$ not RG-derived]
\label{lim:L4}
$\gamma = \GammaGeom$ is permanently \catmark{A-}\ (calibrated). A
first-principles derivation does not exist; any future derivation must
satisfy \cref{prop:target-leakage} in full.\catmark{A-}
\end{limitation}

\begin{limitation}[L5: The $N=99$ RG-step count]
\label{lim:L5}
The choice of $N=99$ discrete cascade steps is empirical: selected to
reach the target magnitude, with no geometric or mathematical principle
dictating exactly $99$. Downstream predictions (such as a Casimir scale
near $\lambda \approx 0.66\,\mathrm{nm}$) presently lack direct
experimental corroboration.\catmark{D}
\end{limitation}

\begin{limitation}[L6: Glueball identification withdrawn]
\label{lim:L6}
The assertion that an exact renormalisation-group flow uniquely
identifies the spectral gap with a specific physical glueball state is
withdrawn from evidentiary use. The toy ODE once dressed in
renormalisation-group language is deleted as a claimed result; residual
mentions are historical only.\catmark{E}
\end{limitation}

\begin{limitation}[L7: EFT validity and the torsion scale]
\label{lim:L7}
The dimension-5 non-minimal coupling makes the Lagrangian an effective
field theory valid only below a cutoff $\Lambda_{\UIDT}$, which must lie
above the gap scale and is \emph{not} identified with the torsion scale
$E_T = \ETorsion$; the latter is a separate, much lower observer-context
quantity (\cref{subsec:et-scale-bridge}). The value of $\Lambda_{\UIDT}$
is not fixed and a non-perturbative UV completion remains
open.\catmark{D}
\end{limitation}

\begin{limitation}[L8--L11: Unrecovered register entries]
\label{lim:L8toL11}
Four further limitations are referenced in the developmental record but
could not be recovered to a definite wording. They are marked
\textsc{unknown} and remain so until the underlying source notes are
restored; they are not reconstructed.\catmark{E}
\end{limitation}

\begin{limitation}[L13: Background independence only partial]
\label{lim:L13}
Although the framework belongs to the pregeometric tradition,
\cref{ax:emergence} is written as a perturbation about a fixed Minkowski
metric, so $\eta_{\mu\nu}$ is inserted by hand rather than derived from
a relational phase. Full background independence has not been
established. The Weinberg--Witten theorem~\cite{weinbergwitten1980} bars any interpretation of the
emergent metric as a composite spin-2 graviton in the naive flat-space
sense --- consistent with, not a refutation of, the thermodynamic
reading borrowed from Jacobson~\cite{jenkins2009,carlip2014} --- and the framework has not yet
specified which loophole its construction exploits. Under D20 this
limitation acquires a sharper address: the dynamical origin of
$\eta_{\mu\nu}$ is exactly the Part-II passage
(\cref{subsec:no-go-points}), where it remains open.\catmark{D}
\end{limitation}

\begin{limitation}[L14: Scalar-glueball spectroscopy and the status of $m_S$]
\label{lim:L14}
Beyond the operator-matching discipline of L12, recent exploratory
lattice work using glueball, meson, and meson--meson operators has been
summarised as suggesting that no scalar state below about
$2\,\mathrm{GeV}$ can be regarded as predominantly a glueball state~\cite{morningstar2024} ---
an explicitly hedged single study that does not by itself overturn the
pure-gauge consensus near $1.5$--$1.7\,\mathrm{GeV}$~\cite{athenodorou2020,morningstar1999}, but sharpens the
concern. The candidate resonance $m_S = 1.705 \pm 0.015\GeV$ is
therefore carried at \catmark{D}\ and its falsification threshold is
read against the evolving lattice spectrum rather than a fixed
benchmark.\catmark{D}
\end{limitation}

\begin{limitation}[L15: Non-Abelian Casimir geometry versus the scalar prediction]
\label{lim:L15}
The Casimir-scale anomaly near $\lambda \approx 0.66\,\mathrm{nm}$ is
motivated by a scalar-field treatment. First-principles $\mathrm{SU}(3)$
lattice studies of the non-Abelian Casimir effect complicate any naive
scalar reading: between chromometallic plates the gluonic Casimir energy
is well described by a massive effective scalar~\cite{chernodub2023}, but in tube and box
geometries the non-Abelian force is found to be \emph{attractive} where
a massless scalar would predict repulsion~\cite{ngwenya2025}. The $0.66\,\mathrm{nm}$
figure is therefore a scalar-sector prediction awaiting both a
non-Abelian-consistent derivation and direct sub-nanometre metrology.\catmark{D}
\end{limitation}

\noindent The register is thus complete in enumeration: anchored where
the corpus permits (L1--L7, L12--L15), and honestly blank where it does
not (L8--L11).\catmark{A}

\section{The \texorpdfstring{$d^2 = 0$}{d2=0} Obstruction}
\label{sec:d2-obstruction}

\subsection{The Failed Scalar-Gradient Route to Curvature}
\label{subsec:failed-scalar-gradient}

The simplest conceivable way to extract a gauge field from a scalar is
to identify the gauge potential with the gradient of the scalar. Its
failure is fundamental, and it is stated with full precision because the
framework treats it as a decisive result --- and because, under D20, it
is the first of the negative results on which the re-architecture
rests (\cref{subsec:negative-results}).

\begin{proposition}[Vanishing Curvature of an Exact Connection]
\label{prop:d2-obstruction}
Let $S$ be a smooth real scalar field and let the connection one-form be
defined as its exterior derivative, $A = \dd S$. Then the curvature
two-form vanishes identically:
\[
 F \;=\; \dd A \;=\; \dd(\dd S) \;=\; \dd^2 S \;=\; 0,
\]
since the exterior derivative satisfies $\dd^2 = 0$ on any smooth
manifold. Consequently the scalar-gradient construction yields only a
flat, pure-gauge connection and cannot produce a non-trivial Yang--Mills
field strength.\catmark{A}
\end{proposition}

\noindent The result is elementary, and that is exactly its force. A
one-form that is exact is automatically closed, and its curvature is
zero by identity, not by approximation; the route fails identically,
everywhere, for any smooth scalar.\catmark{A}\footnote{\textbf{Audit
note (Clebsch loophole).} A known mathematical loophole exists via
Clebsch parameterisations built from higher jets of $S$: with
functionally independent scalars $u,v$ constructed from $S$ and its
derivatives, the one-form $A = u\,\dd v$ yields a non-vanishing
composite \emph{abelian} curvature $F = \dd u \wedge \dd v$.\catmark{A}
This does not rescue the gauge-origin programme: the construction is
tensorial, providing no gauge redundancy and no gauge orbit; it offers
no non-abelian structure; and it supplies no chiral matter. The
obstruction remains exactly as stated: construction-specific, with no
universal impossibility theorem claimed.\catmark{A}} The discrete
counterpart fails for the same reason: a discrete exterior calculus
inherits $\dd^2 = 0$ in its coboundary operator.\catmark{A} With
$A=\dd S$, both contributions to $F = \dd A + A\wedge A$ vanish --- the
first by nilpotency, the second by $\omega\wedge\omega=0$ for any
$1$-form:
\begin{verbatim}
# sympy differential-forms check (schematic):
# d(dS) = 0      (Poincare, d^2 = 0)
# dS ^ dS = 0    (omega ^ omega = 0 for a 1-form)
# => F = d(dS) + dS^dS = 0 identically, for any scalar S
\end{verbatim}
The vanishing is an identity; no precision setting and no input data can
change it.\catmark{A}

\subsection{Topological Consequences and the Interface to Part I}
\label{subsec:topological-consequences}

If $F$ vanishes identically, every topological invariant built from the
scalar-gradient connection vanishes with it: no instanton number, no
non-trivial winding, nothing that could seed a topological
susceptibility.\catmark{A} This is why the topological observables of
\cref{subsec:topology-chi} cannot be sourced from $S$ alone, and why the
effective sector couples $S$ to a Yang--Mills sector that is already
present.

What this does and does not establish must be stated precisely. It does
\emph{not} establish that non-abelian gauge symmetry can never arise
from a scalar primitive; no such impossibility theorem exists, and the
rejected ``algebraic emptiness'' impossibility claim stays
rejected.\catmark{A} Nor may numerical proximity be presented as a
``seal'' that closes the gauge-origin question; seal language is
rejected wherever it appears.\catmark{A} What the obstruction
establishes is narrower and certain: the specific route $A = \dd S$
fails identically --- and the space of non-naive routes is exactly what
Part~I now maps, with its localization of the obstruction to the
multiplicity insertion (G1--G4), its corrected admissibility lattice,
its lower bounds on any intermediate algebra, and its candidate
programmes, all at their stated classes. The v3.9 fork --- strict scalar
versus enriched primitive --- has been taken: D20 commits the research
programme to the enriched arm, at \catmark{D}, with the strict-scalar
arm's honest boundary preserved here as the permanent record of why.

\section{Cosmological Calibration}
\label{sec:cosmological-calibration}

\noindent Every cosmological quantity in this section is
Category~\catmark{C}: calibrated to external data, not derived from the
pregeometric kernel or from the historical scalar ansatz, and never presented
as resolving an observational tension. The
framework's cosmological role is that of a consumer of empirical inputs,
not a producer of predictions.\catmark{C}

\subsection{Empirical Inputs and the Calibration Boundary}
\label{subsec:cosmology-inputs}

\paragraph{Hubble rate $H_0$.}
The defining feature of the current $H_0$ situation is a persistent
$\sim 5\sigma$ discrepancy between early- and late-universe
determinations. In units of km\,s$^{-1}$\,Mpc$^{-1}$, the
\textit{Planck} 2018 acoustic-scale analysis in base $\Lambda$CDM gives
$H_0 = 67.4 \pm 0.5$~\cite{planck2018}, while the SH0ES
distance-ladder programme reports $H_0 = 73.04 \pm
1.04$~\cite{riess2022}; combined in quadrature the two sit roughly
$4.9\sigma$ apart. The framework
calibrates its macroscopic metric parameters to this empirical situation
and records the disagreement in the tension ledger; it does \emph{not}
claim to resolve the Hubble tension, and any such claim is a forbidden
overclaim.\catmark{C}

\paragraph{Dark-energy equation of state.}
In the Chevallier--Polarski--Linder parametrisation
$w(a)=w_0+w_a(1-a)$, a pure cosmological constant requires
$w_0=-1,\ w_a=0$. The DESI Data Release~2 baryon-acoustic-oscillation
measurements~\cite{desidr2} report a preference for the quadrant $w_0>-1,\ w_a<0$
(evolving dark energy), with significance depending on the data
combination --- a mild $\sim 2.3\sigma$ BAO--CMB tension in
$\Lambda$CDM, rising to $\sim 2.8$--$4.2\sigma$ for $w_0$--$w_a$ fits
with supernova samples. Combination-dependent central values (e.g.\
$w_0 \approx -0.42 \pm 0.21$, $w_a \approx -1.75 \pm 0.58$ for
DESI+CMB) are carried as Stratum-I inputs at \catmark{C}, flagged for
final confirmation against the published DESI DR2 table. The
supernova-inclusive combinations trace an equation of state that crosses
$w=-1$ at intermediate redshift --- a phantom crossing that a single
canonical real scalar cannot smoothly realise without a ghost
instability, which is precisely why these are imported as external
constraints rather than derived from the kernel or the historical scalar ansatz.\catmark{C}

\paragraph{Matter fluctuation amplitude $S_8$.}
The parameter $S_8\equiv\sigma_8\sqrt{\Omega_m/0.3}$ is currently
survey-dependent: KiDS-Legacy cosmic shear reports
$S_8 = 0.815^{+0.016}_{-0.021}$ (agreeing with \textit{Planck} at
$\sim 0.73\sigma$), while the DES Year-6 $3\times2$pt analysis reports
$S_8 = 0.789 \pm 0.012$~\cite{kids2025,des2026}. The framework treats $S_8$ as a
survey-and-pipeline-dependent input its coarse-graining parameters are
bounded by; it calibrates to the spread and does not relabel it as a
resolved tension.\catmark{C}

\begin{table}[H]
\footnotesize
\caption{Stratum-I empirical inputs to which \UIDT\ calibrates. All
entries are Category~[C]; none is a prediction or a tension resolution.}
\label{tab:cosmology-inputs}
\renewcommand{\arraystretch}{1.3}
\begin{tabular*}{\textwidth}{@{}p{0.18\textwidth}@{\extracolsep{\fill}}p{0.35\textwidth}p{0.24\textwidth}c@{}}
\toprule
\textbf{Parameter} & \textbf{Value $\pm$ uncertainty} & \textbf{Source} & \textbf{Class} \\
\midrule
$H_0$ (early) & $67.4 \pm 0.5\ \mathrm{km\,s^{-1}\,Mpc^{-1}}$ & Planck 2018 & \catmark{C} \\
$H_0$ (late) & $73.04 \pm 1.04\ \mathrm{km\,s^{-1}\,Mpc^{-1}}$ & SH0ES (Riess 2022) & \catmark{C} \\
$w_0, w_a$ & $w_0\!\approx\!-0.42\pm0.21,\ w_a\!\approx\!-1.75\pm0.58$ (DESI+CMB; combination-dependent) & DESI DR2 & \catmark{C} \\
$S_8$ & $0.815^{+0.016}_{-0.021}$ / $0.789\pm0.012$ & KiDS-Legacy / DES~Y6 & \catmark{C} \\
\bottomrule
\end{tabular*}
\end{table}

\subsection{The Relational-Proliferation Scenario: Interpretive Status}
\label{subsec:information-big-bang}

The developmental corpus contains a cosmogonic narrative, retained
strictly as Category~\catmark{E}\ interpretation within Stratum~III ---
a way of thinking, fed into a Category~\catmark{C}\ calibration, not a
derivation of cosmological data.\catmark{E} The ``beginning'' is read
not as a physical instant $t=0$ but as the first logical distinction in
the sense of \cref{subsec:sx-zero-not-unmarked}; cosmic inflation is
reinterpreted, again at \catmark{E}, as a rapid proliferation of
relational structure that, viewed under coarse-graining, appears in the
effective regime as an exponentially expanding metric. The picture
computes no inflationary observable, predicts no primordial spectrum,
and does not touch the \catmark{C}\ cap.\catmark{E}

\section{The Verification Return Path and Final Claims}
\label{sec:closed-circle}

\subsection{The Empirical-to-Ledger Return Path}
\label{subsec:return-path}

The verification architecture traces a controlled return path from external
comparators to the parameter ledger without promoting that path to ontology. One begins from the empirical and lattice stress
tests of pure $\mathrm{SU}(3)$ (\catmark{B}\,pending their source and
covariance audit); computes the framework-side quantities \emph{blindly}
per \cref{prop:target-leakage} (\catmark{D}); compares \emph{post hoc}
through an explicit uncertainty and covariance model; identifies structural audit facts \catmark{A} and open programme failures
\catmark{D}; and interprets last, at \catmark{E}. The closure is governed throughout by the two structural
limits this draft has doubled down on: the $d^2=0$ obstruction
(\cref{prop:d2-obstruction}) and the localized
$G_{\mathrm{SM}}$-origin gap (\cref{sec:g1g4}). The return path terminates in a disciplined accounting rather than an
ontological closure. Its auditability is a software and governance fact
\catmark{A}; its physical interpretation remains bounded by the individual
claim classes.

\subsection{Final Allowed Claims}
\label{sec:final-claims}

{\small
\renewcommand{\arraystretch}{1.3}
\begin{longtable}{p{0.10\textwidth} p{0.56\textwidth} p{0.10\textwidth} p{0.12\textwidth}}
\caption{Final allowed claims, with conservative evidence classes and
strata. Negative and open entries are included deliberately; classes are
unchanged from the incoming ledger; D20 changes architectural role only (ONT-11).}
\label{tab:claims}\\
\toprule
\textbf{ID} & \textbf{Claim} & \textbf{Class} & \textbf{Stratum} \\
\midrule
\endfirsthead
\toprule
\textbf{ID} & \textbf{Claim} & \textbf{Class} & \textbf{Stratum} \\
\midrule
\endhead
\bottomrule
\endfoot
ONT-01 & In the historical scalar grammar, a vanishing field value is a marked state, not the unmarked state; this is retained only in Appendix~\ref{app:scalar-museum}. & \catmark{A} & II/III \\
ONT-02 & A morphism $f\colon A\to B$ presupposes distinguishable $A,B$; morphism is a formal tool, not an ontological generator. & \catmark{A} & II/III \\
ONT-03 & With $A=\dd S$, the curvature $F=\dd^2 S=0$ identically; the naive scalar-gradient route fails. & \catmark{A} & II \\
ONT-04 & $\gamma = \GammaGeom$ is calibrated, never derived from RG first principles. & \catmark{A-} & III \\
ONT-05 & $\Delta = \DeltaGap$: internally consistent, lattice-compatible; determination D18 (\cref{box:pi-override-delta}); not externally peer-confirmed, not derived. & \catmark{B} & I/III \\
ONT-06 & $\kappa = \KappaRG$ exact framework input (fixed-point reading a candidate condition, untested by any defined flow); $5\kappa^2 = 3\lambda_S$ closes internally to $<10^{-14}$. The empirically calibrated reading $\kappa = 0.500 \pm 0.008$ is a separate, lower-class statement. & \catmark{A} & II \\
ONT-07 & Pure $\mathrm{SU}(3)$ ratio comparisons are post-hoc stress tests with solver access disabled. & \catmark{B}\ (\emph{pending external verification}) & I/III \\
ONT-08 & Primitive $\to\mathcal{A}\to G_{\mathrm{SM}}$ is the research programme; localization, lower bounds, and candidates are recorded in Part~I. & \catmark{D} & III \\
ONT-09 & Cosmological parameters are calibrated, capped at~C; no tension is declared resolved. & \catmark{C} & I/III \\
ONT-10 & The observer functor $F\colon V\to\mathrm{Obs}$ and the compression taxonomy are bounded interpretive frameworks. & \catmark{E} & III \\
ONT-11 & Under PI Decision D20, Part~III is a verification/benchmark/ledger stratum rather than active ontology; the kernel candidate of Part~I remains a research programme at \catmark{D}. & \catmark{D} & III \\
\end{longtable}
}

\noindent The conservatism of \cref{tab:claims} is its principal
feature: its firmest ground is its boundary results, and its ambitions
are labelled as ambitions.\catmark{A}

\subsection{Final Forbidden Claims}
\label{subsec:forbidden-claims}

\begin{governancebox}[Final Forbidden Claims]
\textbf{(1)} No derivation of the full Standard-Model gauge group
$G_{\mathrm{SM}}$ --- from the historical scalar ansatz, and equally none from the
kernel candidate: Part~I records selection, not derivation.\\[0.3em]
\textbf{(2)} No derivation of the Born rule.\\[0.3em]
\textbf{(3)} No solution of consciousness or qualia.\\[0.3em]
\textbf{(4)} No validation of the historical scalar ansatz or of the kernel by psychedelic,
near-death, or mystical experience.\\[0.3em]
\textbf{(5)} No use of $E_T$ as a threshold for life, death, ego,
psychedelic states, or artificial general intelligence.\\[0.3em]
\textbf{(6)} No identification of $\Delta$ with the $0^{++}$ glueball
without operator matching (L12).\\[0.3em]
\textbf{(7)} No use of quarantined museum artifacts as derivation or
factual evidence.\\[0.3em]
\textbf{(8)} No presentation of $5\kappa^2 = 3\lambda_S$ as an
independent empirical discovery where $\lambda_S$ is fixed
definitionally by that same constraint.\\[0.3em]
\textbf{(9)} No claim that DESI, $H_0$, $S_8$, or $w_0$--$w_a$ data
validate \UIDT\ beyond calibrated Category~\catmark{C}\
status.\\[0.3em]
\textbf{(10)} No claim that a lattice-compatible number establishes the
Standard-Model gauge group, a confinement mechanism, or a Clay-grade
Yang--Mills construction.\\[0.3em]
\textbf{(11)} No presentation of the conditional $N=6$ selection
proposition (\cref{prop:unique321}) as an unconditional derivation of
the Standard-Model algebra; its grade is bounded by H1 and H2.\\[0.3em]
\emph{All eleven are Class~\catmark{A}\ standing constraints.}
\end{governancebox}

\noindent The seventh prohibition guards the integrity of all the
others. The museum of refuted constructions --- the numerological
decomposition of $\gamma$, the failed scalar-gradient operator, the toy
ODE dressed in renormalisation-group language, and now the quarantined
block-repetition audit of \cref{subsec:multiplicity-precision} --- is
retained for its documentary value, as evidence of method; its contents
may never be cited as factual evidence for any claim.\catmark{A}

\section{Falsification Gates and the Tension Ledger}
\label{sec:falsification-tension}

\subsection{The Falsification Gates}
\label{subsec:falsification-gates}

The framework's discipline is enforced by gates through which every
proposed claim, computation, or contribution must pass. They are
\emph{advisory and non-approving}: a gate can block, but a gate that is
passed authorises nothing. Passing all eight is necessary hygiene, never
sufficient warrant.\catmark{A}

\begin{governancebox}[The Eight Falsification Gates (Advisory, Non-Approving)]
\textbf{F1: No target leakage.} No solver may access the value it claims
to predict (\cref{prop:target-leakage}); in the Lean corpus, no
definition or hypothesis may contain the candidate partition
(\cref{subsec:code-evidence}).\\[0.3em]
\textbf{F2: Local precision discipline.} Working precision is set
locally per proof-critical block; no global precision mutation; no
proof-critical floating-point or rounding heuristics.\\[0.3em]
\textbf{F3: Complete benchmark accounting.} The requirement that every
external benchmark carry a verified source, uncertainty budget, and
covariance model is an internal audit rule. \catmark{A} Satisfaction of
that external benchmark condition remains pending at \catmark{B}.\\[0.3em]
\textbf{F4: L12 separation.} $\Delta$ is not identified with the
$0^{++}$ glueball without operator matching (\cref{lim:L12}).\\[0.3em]
\textbf{F5: $G_{\mathrm{SM}}$-origin honesty.} The gauge-origin gap is
stated as open (\cref{sec:g1g4}); no derivation is claimed; selection
propositions carry their conditional grades.\\[0.3em]
\textbf{F6: Observer evidence boundary.} No phenomenological report
validates the ontology (\cref{sec:phenomenology-boundary}).\\[0.3em]
\textbf{F7: Born-rule status.} The Born rule is recorded as borrowed and
open, never as a derivation (\cref{subsec:born-rule}).\\[0.3em]
\textbf{F8: Governance boundary.} No machine-generated verdict
authorises a merge, freeze, or evidence-class upgrade --- and no formal
proof, once machine-checked, authorises a physics upgrade without
governance.\\[0.3em]
\emph{Each gate is Class~\catmark{A}. Passing is advisory; it approves
nothing.}
\end{governancebox}

\subsection{The Tension Ledger}
\label{sec:tension-ledger}

{\small
\renewcommand{\arraystretch}{1.3}
\begin{longtable}{p{0.08\textwidth} p{0.40\textwidth} p{0.22\textwidth} p{0.18\textwidth}}
\caption{Tension ledger. Unresolved disagreements and gaps, displayed
rather than smoothed.}
\label{tab:tensions}\\
\toprule
\textbf{ID} & \textbf{Tension / gap} & \textbf{Status} & \textbf{Resolution path} \\
\midrule
\endfirsthead
\toprule
\textbf{ID} & \textbf{Tension / gap} & \textbf{Status} & \textbf{Resolution path} \\
\midrule
\endhead
\bottomrule
\endfoot
T1 & $\gamma_{\mathrm{MC}} = \GammaMC$ versus calibrated $\gamma = \GammaGeom$. & \textbf{[consistent within MC variance]}; displayed separately, not merged by selection. & Reduce MC variance; non-leaking derivation of $\gamma$. \\
T2 & $\Delta = \DeltaGap$ evidence class. & Would sit \catmark{E}\ (\emph{pending bridge}) absent external work; determination D18 sets \catmark{B}\ (\cref{box:pi-override-delta}). & External peer counter-signature. \\
T3 & $\Delta$ versus $0^{++}$ glueball identity. & \catmark{A}\ separation enforced (L12); identity unproven. & Operator matching. \\
T4 & SVZ leading-order topology route versus lattice topology. & Separation of the two routes is an audit fact \catmark{A}; external lattice agreement remains pending at \catmark{B}. The SVZ route is museum-only. & Audited lattice $\chi_{\mathrm{top}}$ ratios. \\
T5 & $E_T$ observer/neural scale bridge. & Its absence and the stated scale mismatch are audit facts \catmark{A}; any proposed bridge remains interpretive at \catmark{E}. & Derive or refute a principled scale bridge. \\
T6 & Born-rule derivation. & Absent; borrowed from Stratum~II. \catmark{A} & Non-circular coarse-graining derivation. \\
T7 & Cosmological tensions ($H_0$, $S_8$). & Capped at \catmark{C}; calibrated, \emph{not} closed. & Independent prediction, not fit. \\
T8 & $[2,2,2]$ rejection (H2) rests on an unverified mass-degeneracy argument (\cref{tab:lattice}). & \catmark{D}; the weakest verdict in the corrected lattice; the conditional selection proposition inherits it. & Primary-source theorem or free-energy stability analysis settling the case either way. \\
T9 & Scalar-glueball spectroscopy versus $m_S = 1.705 \pm 0.015\GeV$ and the empirical grounding of $\Delta$. & \catmark{D}; exploratory lattice work~\cite{morningstar2024} suggests no state below $2\,\mathrm{GeV}$ is predominantly a glueball; the $1.5$--$1.7\,\mathrm{GeV}$ pure-gauge consensus is sharpened, not overturned, and the empirical grounding of $\Delta$ is correspondingly under-determined (L14). & Read $m_S$ against the evolving lattice spectrum, not a fixed benchmark; monitor the sub-$2\,\mathrm{GeV}$ picture. \\
\end{longtable}
}

\noindent Two entries warrant closing remarks. T1: the calibrated
kinetic value and the Monte-Carlo mean differ by $0.035\,\mathrm{GeV}$,
i.e.\ $0.035\sigma$ given the $\pm1.005$ spread, and are displayed together
rather than reconciled by selection.\catmark{A} T8 is new in this draft
and is deliberate: the formal selection layer of Part~I made the
weakest link of the admissibility lattice \emph{sharper}, not weaker,
and the ledger records it where it can be seen. The ledger remains the
most honest page in the manuscript: a framework that hid this page would
be easier to admire and impossible to trust.\catmark{A}

%========================== END P4 PART III CONTENT ========================


%=============================================================================
% BIBLIOGRAPHY
%=============================================================================
\newpage

%============================== R1 CODA CONTENT ============================
% Restores UIDT Ontology v3.9.9 (-006) Chapters 20-21 and the Adversary Test,
% adapted to the D20 architecture. Dropped in error during P2-P4; recovered in
% the R0 reading stage. Evidence classes unchanged from -006 throughout.
%===========================================================================

\newpage
\part*{Coda: Epistemological Position and Self-Audit}
\addcontentsline{toc}{part}{Coda: Epistemological Position and Self-Audit}
\label{part:coda}

\noindent The three parts of this draft have stated what the framework
holds, at what evidence class, and under what falsification conditions.
This coda states something different and equally necessary: the
philosophical position from which those statements are made, the company
the framework keeps, and the outcome of turning the framework's own
adversarial machinery against its strongest-looking claims. Every
assertion in this coda is Category~\catmark{E}, Stratum~III, unless
explicitly tagged otherwise, and none of it alters any evidence class
established earlier.

\section{Epistemological Position and Comparisons}
\label{sec:epistemology}

\subsection{Structural Realism, Held as a Position}
\label{subsec:structural-realism}

The ontological proposal that physical reality is constituted by
mathematical structure raises an epistemological question: does \UIDT\
claim the universe \emph{is} a mathematical object, or merely that it is
\emph{best described by} one? The framework's working answer leans toward
the former, an ontic-structural-realist position, but it is held as a
philosophical \emph{position}, not asserted as a result, and nothing in
the physics strata depends on it.\catmark{E} The proposal is narrow: not
that all mathematical structures are physically real, but that one
specific structure is offered as the candidate constituting structure.

\begin{governancebox}[What D20 Does to the Structural-Realist Position]
Under PI Decision D20 the identity of the candidate constituting
structure changes, and the position is thereby \emph{weakened}, not
strengthened. In v3.9.x the candidate was the self-interacting vacuum
information density of \cref{eq:lagrangian}: a concrete Lagrangian, with
a stated minimality argument. In this draft that Lagrangian describes the
\emph{effective} stratum, and the candidate constituting structure is the
finite algebraic pregeometric kernel of Part~I --- which is held at
\catmark{D}, is not uniquely specified, and whose Wedderburn invariants,
real structure, and KO-dimension are free moduli. A structural realism
whose structure is a research programme is a position about a
\emph{placeholder}. The framework states this rather than concealing it:
the ontological ambition of Part~I exceeds its formal content by exactly
the distance between \catmark{D} and \catmark{A}, and that distance is
the honest measure of the position's current standing.\catmark{E}
\end{governancebox}

\noindent The selection is meant to be structural rather than anthropic.
For the effective sector, the Lagrangian of \cref{eq:lagrangian} is the
minimal gauge-invariant, vacuum-stable, EFT-admissible self-interaction
of a real scalar non-minimally coupled to $\mathrm{SU}(3)$ in $3+1$
dimensions; for the kernel, the analogous minimality claim does not yet
exist, and \cref{oq:staircase} together with the conditional selection
proposition (\cref{prop:unique321}) marks the place where it would have
to be constructed. That ``minimality'' is in both cases a claim about the
framework's construction, not a proven uniqueness theorem, and is carried
as such.\catmark{E}

The position is explicitly bounded by its own incompleteness. The
framework acknowledges that any sufficiently expressive formal system is
subject to G\"odel-type incompleteness, and it reads its limitation
register (\cref{subsec:limitation-register}) as a concrete, honest
catalogue of what it cannot presently derive rather than as a temporary
embarrassment. The bound is sharper than it was: the undecidability of
the spectral-gap problem for general local translation-invariant
Hamiltonians (Cubitt--P\'erez-Garc\'ia--Wolf) is an external theorem
showing that some questions of exactly this shape admit no algorithm and
are independent of the axioms.\catmark{A} That result is not a statement
about the Yang--Mills gap and not a statement about $\Delta$; it is a
statement about the class of questions to which such questions
belong.\catmark{A}

The framework also distinguishes itself from a simulation hypothesis:
\UIDT\ posits no external simulator and no computational substrate.
Whatever force the structural-realist picture has, it is ontological
speculation, not a claim that the universe is being computed by
something.\catmark{E} The computational vocabulary that appears in the
developmental record --- interfaces, rendering, frontend and backend ---
is metaphor, and is barred from any argument in this manuscript.

\subsection{Comparison with Alternative Frameworks}
\label{subsec:alt-frameworks}

It is useful, and disciplined, to state where \UIDT\ sits relative to
established programmes, always as comparison, never as a claim of
superiority.

Against string theory, \UIDT\ is far more ontologically frugal but
correspondingly far less developed and far less mathematically mature;
frugality is not a virtue if it buys its simplicity by leaving the hard
problems unsolved, and the $G_{\mathrm{SM}}$-origin gap is exactly such
an unsolved problem.\catmark{E} Against entropic gravity
(Verlinde~\cite{verlinde2011}, Padmanabhan~\cite{padmanabhan2010}),
\UIDT\ shares the intuition that geometry is emergent and thermodynamic
in spirit but does not claim to have improved on it; the relationship is
one of shared intuition, not demonstrated advance.\catmark{E} Against
loop quantum gravity, \UIDT\ shares the \emph{aspiration} to background
independence but does not yet achieve it in full: the metric-emergence
map (\cref{ax:emergence}) is written as a perturbation about a fixed
Minkowski background, so background independence is at present only
partial (\cref{lim:L13}); it takes the opposite route to geometry (a
continuum coarse-graining limit rather than quantised geometric
structures); and like LQG it has not recovered the full Standard-Model
particle content --- indeed \UIDT\ has not recovered it at all, which is
the content of the open gap.\catmark{E}

Two comparisons are new to this draft, and both are unflattering. Against
the noncommutative-geometry programme of Connes, Chamseddine, and
collaborators, Part~I is a strict consumer: every structural constraint
it uses --- Krajewski combinatorics, multiplicity matrices, the
KO-dimension-6 real structure, the classification of irreducible
geometries --- is imported, and \UIDT\ has added no theorem to that
literature. Against the matrix-model programme (IKKT, BFSS, and their
condensation studies), the framework's dynamical-selection search
returned an honest null. The relationship in both cases is that of a
framework standing on external results, and the correct posture is
acknowledgement rather than positioning.\catmark{E}

The framework's one genuine methodological commitment here is an
epistemic razor: it admits no unobservable extra dimensions, no
fundamental supersymmetry, and no infinite landscape, and it ties its
claims to a small set of parameters with explicit falsification
thresholds (\cref{sec:falsification-tension}). This makes the framework
vulnerable, and that vulnerability, not any claimed explanatory triumph,
is what the framework offers as its scientific virtue.\catmark{E} The
earlier developmental claim that \UIDT\ ``eliminates'' dark matter or
dark energy is withdrawn: cosmology is capped at \catmark{C}
(\cref{sec:cosmological-calibration}), and no such elimination is
asserted.

\subsection{Relation to the Information-Theoretic and Structural-Realist Tradition}
\label{subsec:prior-art}

The general move that \UIDT\ makes --- taking an information-theoretic
primitive as ontologically basic and recovering spacetime, and
aspirationally gauge structure, as emergent --- is not itself novel, and
intellectual honesty requires locating it explicitly within an
established tradition rather than presenting it as unprecedented.

The lineage begins with Wheeler's ``it from bit'' programme, in which
information is treated as the substrate from which physical things
derive~\cite{wheeler1990}. It continues through the thermodynamic and
entropic-gravity derivations of geometry
(Jacobson~\cite{jacobson1995}; Padmanabhan~\cite{padmanabhan2010};
Verlinde~\cite{verlinde2011}), the matrix-model emergent-geometry
programme~\cite{ikkt1997,steinacker2010}, and group field theory, in
which spacetime emerges from a condensate. Within holography, Van
Raamsdonk argues that the emergence of classically connected spacetimes
is intimately related to the entanglement of degrees of freedom in a
non-perturbative description of quantum gravity, and that disentangling
the degrees of freedom associated with two regions causes those regions
to pull apart and pinch off, quantifiable by standard entanglement
measures~\cite{vanraamsdonk2010}. That statement is a
Stratum-II\,\catmark{B} result \emph{within} AdS/CFT; \UIDT\ has no
holographic dual and no boundary conformal field theory, so the
generalisation ``spatial connectivity \emph{is} entanglement'' is a
Stratum-III reading at \catmark{E} and is not inherited by this
framework.\catmark{E} A particularly close ontological cousin is the
Carroll--Singh proposal~\cite{carrollsingh2018}, in which the only
fundamental objects are a vector in Hilbert space and a Hamiltonian;
\UIDT\ under D20 has moved \emph{toward} that austerity, replacing a
concrete real-scalar field with a finite algebraic candidate, while
sharing the aim of a maximally frugal primitive.

\begin{evidencebox}[Prior-Art Register (Thermodynamic-Gravity Bridge)]{catE}
\textbf{RMC-01} In AdS/CFT, disentangling boundary degrees of freedom
disconnects the bulk geometry~\cite{vanraamsdonk2010}. \catmark{B},
Stratum~II.\\[0.25em]
\textbf{RMC-01b} The generalisation ``spatial connectivity \emph{is}
entanglement'' beyond holography. \catmark{E}, Stratum~III. Not inherited
by \UIDT.\\[0.25em]
\textbf{RMC-02} Time as a local informational calibration measure.
\catmark{E}, Stratum~III.\\[0.25em]
\textbf{RMC-03} Continuous gravity as an effective field theory; the
thermodynamic-gravity results of Jacobson, Padmanabhan, and Verlinde are
\emph{precedent} that emergent geometry is a respectable position, not
derivations imported into \UIDT. \catmark{A} re: corpus (as a statement
about what \UIDT\ does and does not derive); \catmark{E} for the mapping.
\end{evidencebox}

\noindent Against this background, what is specific to \UIDT\ --- and
what is offered as its actual contribution, as distinct from the shared
tradition above --- is narrow and concrete: the particular Lagrangian of
\cref{eq:lagrangian} with its non-minimal $S\,F^2$ coupling at a stated
cutoff; the identification of the spectral gap as a contraction fixed
point; the calibrated dimensionless invariant $\gamma =
\GammaGeom$\,\catmark{A-}; the localization of the
$G_{\mathrm{SM}}$-origin obstruction to the four conditions G1--G4 with
quantitative lower bounds on any admissible intermediate algebra
(\cref{subsec:lower-bounds}); and the explicit evidence-grading
discipline under which every claim in this manuscript is carried. None of
these is a claim to have solved the emergence problem that the tradition
leaves open; they are the specific, falsifiable form in which \UIDT\
states its version of it.\catmark{E}

\section{Scope, Contingency, and Interpretive Minimality}
\label{sec:inevitability-sociology}

\subsection{From ``Unreasonable Effectiveness'' to a Modest Minimality}
\label{subsec:inevitability}

Wigner's remark on the unreasonable effectiveness of mathematics invites
a strong reading: that if the universe \emph{is} mathematical structure,
the effectiveness is no longer surprising. The framework records this as
a tempting interpretive move and explicitly declines its strong form. To
call the effectiveness a ``tautology'' or the structure ``inevitable''
would be exactly the kind of rhetorical overreach the manuscript's
discipline forbids; the honest statement is far more modest. Within the
framework's own axioms and within $3+1$ dimensions, the $\mathrm{SU}(3)$
Lagrangian of \cref{eq:lagrangian} is presented as the \emph{minimal}
self-interacting structure compatible with gauge invariance, vacuum
stability, and EFT admissibility. That is a statement about minimality
under stated constraints, not a proof of cosmic necessity, and the
framework does not claim the universe \emph{had} to be this
way.\catmark{E}

The question ``why this universe and not another?'' is therefore
answered, if at all, by structural minimality rather than by modal
necessity or anthropic selection. Under D20 even that answer contracts:
the minimality argument applies to the effective sector, and no
minimality argument exists for the kernel. The spectral gap
$\Delta = \DeltaGap$ remains, in this reading, a structurally determined
feature of the construction's fixed point, but its evidence class is
unchanged by any of this philosophy: it is \Bdelta, and the philosophical
gloss adds no evidential weight.\catmark{E}

\begin{governancebox}[A Framework Cannot Be Made Irrefutable]
\label{gov:irrefutable}
A recurring temptation, in a research programme that has survived several
audit cycles, is to seek a formulation so guarded that it can no longer
be wrong. That ambition is incoherent, and \cref{ax:falsifiability}
forbids it. A framework that cannot fail is not a physics framework; it
is a vocabulary. What can be sought, and what this manuscript seeks, is a
formulation whose claims are so precisely bounded that a hostile reader
finds no gap between what is asserted and what is shown. That is not
irrefutability. It is the opposite: it is the condition under which
refutation becomes possible and informative.\catmark{A}
\end{governancebox}

\subsection{Reception Risk and Non-Evidential Context}
\label{subsec:sociology}

Finally, the framework reflects, again purely as \catmark{E}\ context, on
how a radically simplifying ontology is liable to be received. A
framework whose evidentiary classes are not kept explicit can be misread
in both directions: dismissed wholesale on account of its speculative
stratum, or embraced wholesale on the strength of its calibrated
agreements.\catmark{E} The framework's response is not to claim
vindication but to insist on the opposite safeguard: the
evidence-classification system and the verification discipline exist as
bulwarks against \emph{both} premature rejection \emph{and} uncritical
acceptance. The right attitude toward a radical simplification is neither
enthusiasm nor dismissal but exactly the falsification-first scrutiny the
manuscript has tried to embody; by that standard, the framework's
simplifications must earn their place by surviving tests, not by being
elegant.\catmark{E}

\begin{evidencebox}[Philosophical Context: Standing Status]{catE}
\textbf{This entire coda is \catmark{E}, Stratum~III,} except where a
tag states otherwise. Structural realism is a position held, not a
result; under D20 its candidate structure is a \catmark{D} research
programme, which weakens the position rather than strengthening it.
Minimality is a construction claim about the effective sector, not a
uniqueness proof, and no minimality argument exists for the kernel.
``Inevitability'' and ``tautology'' language is withdrawn. Comparisons
with string theory, entropic gravity, LQG, noncommutative geometry, and
matrix models are stated without claim of superiority. The ``eliminates
dark matter/energy'' claim is withdrawn. No philosophical statement here
alters any evidence class.\catmark{E}
\end{evidencebox}

\section{The Adversary Test: A Self-Audit}
\label{sec:adversary}

\noindent The framework maintains an internal adversary test: a
formalised attempt to break its own strongest-looking claim. The
discipline of this edition requires that the test be run honestly, which
means reporting where the adversary \emph{wins}, not staging a defence
that quietly declares victory. This draft runs the test twice: once against the v3.9 scalar claim it
inherits as history, and once against the D20 Phase-6 rebase.

\subsection{First Round: The Origin of \texorpdfstring{$\gamma$}{gamma} (Limitation L4)}
\label{subsec:adversary-gamma}

\paragraph{Adversarial argument.}
If $\gamma$ is a phenomenological fit parameter rather than a derived
constant, then any vacuum-energy or cosmological calculation scaled by a
power of $\gamma$ risks tautology: it would scale by precisely the amount
needed to match observation, draining the result of independent
explanatory power. The adversary presses further: the developmental
record once decomposed $\gamma$ numerologically ($49/3 + 17/3000$, with
the integer $16339$ prime), which is positive evidence that the value was
\emph{constructed} to a target rather than found.

\paragraph{Honest adjudication.}
The framework does not claim the adversary is refuted. On the contrary:
the numerological decomposition is quarantined in the museum precisely
because the adversary is right about it --- that construction is not
evidence and is barred from use. The rejection is \emph{methodological}
before it is arithmetical: the term $17/3000$ was obtained by working
backwards from the calibration target, and the primality of $16339$ is
the evidence that no mechanism could have produced the split. A solver
whose stopping criterion reads $\gamma$ is corrupt regardless of the
elegance of its output.\catmark{A} As for the broader charge, the defence
is partial and is recorded as partial. $\gamma$ is fixed by matching to
a historical internal kinetic reference documented in
Appendix~\ref{app:scalar-museum}, independently of cosmology, so a
vacuum-energy calculation using $\gamma$ is not circular \emph{with
respect to the cosmological data}; but this does not make $\gamma$
derived. The honest verdict is that the adversarial argument
\emph{stands} on its central point: $\gamma$ is \catmark{A-}, calibrated,
not derived, and the framework's protection against tautology is not a
first-principles derivation but the structural separation of strata
together with the anti-target-leakage rule
(\cref{prop:target-leakage}).\catmark{A-}

\subsection{Second Round: The D20 Phase-6 Rebase}
\label{subsec:adversary-d19}

\paragraph{Adversarial argument.}
The Phase-6 rebase, says the adversary, is an evasion dressed as
rigour. Faced with a primitive that cannot generate gauge structure, the
framework has not solved the problem; it has renamed the primitive.
Three charges follow. \emph{First}, the kernel is unestablished: Part~I
offers a \catmark{D} candidate whose block partition, real structure, and
KO-dimension are free moduli, so the ontological office has been
transferred from one unproven occupant to another. \emph{Second}, the
technical move that appears to rescue the programme --- replacing the
exterior derivative by an inner derivation $\delta(\cdot) = [\mathbf
S,\cdot]$, so that the nilpotency $d^2=0$ no longer applies --- buys
nothing: the algebra generated by a single self-adjoint operator is
commutative by the spectral theorem, so $\delta$ vanishes identically on
it, and a non-vanishing commutator requires an ambient noncommutative
algebra that has been \emph{inserted}. \emph{Third}, the one crisp
positive result of Part~I, the uniqueness of the $[3,2,1]$ partition at
$N=6$, is conditional on a design-level filter and a heuristic filter,
and its Lean proof obligation is open at the audited head.

\paragraph{Honest adjudication.}
All three charges are conceded.\catmark{A} On the first: the transfer is
real, and \cref{subsec:structural-realism} states it in those terms; what
the re-architecture buys is not a foundation but a sharper obstruction
--- G1--G4 with quantitative lower bounds, where v3.9 had only a
declaration that the gap was open. On the second: the adversary's algebra
is correct, and the framework adopts the point as its own result rather
than resisting it; the passage to inner derivations evades the $d^2=0$
identity but relocates the obstruction to the insertion of the ambient
algebra, which is condition G1. Any claim that this passage ``lays the
foundation for gauge groups'' is withdrawn. On the third: the
conditionality is stated in \cref{prop:unique321} and its grade is
\catmark{D}; discharging the Lean obligation would upgrade the formal
layer and nothing else. The framework's answer to all three is the same:
these are the honest boundaries, and the manuscript exists to make them
explicit rather than to argue past them.

\paragraph{What the self-audit establishes.}
The adversary test is not passed by winning the argument; it is passed by
surviving the argument with the evidence classes intact. $\gamma$ remains
\catmark{A-}; the museum construction stays refuted; the calibration is
non-circular with respect to cosmology but is not a derivation; the
kernel remains \catmark{D}; the inner-derivation escape is a relocation,
not a closure; and the selection proposition remains conditional. This is
the appropriate outcome of an honest self-audit: the framework's weakest
load-bearing points are identified, the adversary's valid points are
conceded, and nothing is relabelled to make the defence look stronger
than it is.\catmark{A}

\begin{governancebox}[Adversary Test: Outcome]
\textbf{Round 1 (v3.9, $\gamma$).} The strongest attack targets the
origin of $\gamma$ (L4). The numerological construction is conceded to
the adversary and stays quarantined; the rejection is methodological
(backward construction), with primality as the evidence. The calibration
of $\gamma$ is non-circular with respect to cosmology but is \emph{not} a
derivation; $\gamma$ remains \catmark{A-}.\\[0.35em]
\textbf{Round 2 (D20).} All three charges conceded: the kernel is an
unestablished \catmark{D} candidate; the inner-derivation route relocates
the $d^2=0$ obstruction rather than closing it, and any ``foundation for
gauge groups'' reading is withdrawn; the $N=6$ selection proposition is
conditional on a design-level and a heuristic filter, with its formal
obligation open.\\[0.35em]
\textbf{Verdict.} The test is ``passed'' only in the sense that the
evidence classes survive the attack unchanged, not in the sense that
either adversary is refuted.\catmark{A-}
\end{governancebox}

%============================ END R1 CODA CONTENT ==========================

%%ANCHOR-P5-BIBLIOGRAPHY%%

%============================ P5 BACK MATTER ===============================
% Inserted at ANCHOR-P5-BIBLIOGRAPHY. Bibliography carried over verbatim from
% UIDT Ontology v3.9.9 (-006), 94 verified entries, with two audit corrections
% applied and marked in place. No entry added without a resolved DOI/arXiv id.
%===========================================================================

\section{Reproduction Note}
\label{sec:reproduction-note}

\noindent Every quantitative and formal statement in this draft is
reproducible from the following commands and artefacts. Nothing in this
section constitutes an authorization.

\paragraph{Numerical claims.}
The exact coupling identity retained from the historical scalar/EFT record
is reproduced from Appendix~\ref{app:scalar-museum}; it is an internal
ledger relation, not a premise of the active ontology. Working precision
is set locally, never globally:
\begin{verbatim}
python3 -c "from mpmath import mp,mpf; mp.dps=80; assert mp.dps==80
k=mpf(1)/2; lS=mpf(5)/12; print(mp.nstr(abs(5*k**2 - 3*lS), 80))"
# -> 0.0   internal closure < 1e-14; NOT external agreement
\end{verbatim}
No other numerical value in this draft is computed here; ledger values
are reproduced at their recorded classes, and no new eighty-digit value
is introduced, estimated, or interpolated.

\paragraph{Symbolic claims.}
\cref{prop:d2-obstruction} is an identity of differential forms and
requires no numerics; the schematic check is given at
\cref{subsec:failed-scalar-gradient}. The rigidity results underlying
G1--G4 are representation-theoretic identities: the maximal compact
subgroup of $GL(1,\mathbb{R})$ is $\{\pm 1\}$; loop phases factor
through $H_1(\Gamma) = \pi_1(\Gamma)^{\mathrm{ab}}$; Skolem--Noether
gives $\mathrm{Aut}(M_N(\mathbb{C})) = PU(N)$; and the rank budget for
level-one simply-laced affine algebras yields $\mathrm{rank} = 1$.

\paragraph{Formal corpus.}
The Lean~4 statements quoted in Part~I remain a historical snapshot of the
repository sources at commit \texttt{61bc24da}; MSS-000 does not claim that
this manuscript excerpt is the live repository bridge. The audit reported in
\cref{subsec:formalization-status} is a \emph{static structural audit}
of that snapshot: proof-obligation counts are from a whole-word scan for
\texttt{sorry}; no \texttt{lake build} was executed, and no build result
is asserted anywhere in this draft. Reproduction of that historical audit requires fetching the stated commit
and repeating the scan; live bridge reconciliation is a separate
post-application executor obligation; Findings F1--F3 are stated so that
they can be checked, and refuted, at that commit.

\paragraph{Empirical values.}
All Stratum-I values (\cref{tab:cosmology-inputs} and the benchmark
discussion of \cref{sec:benchmark-batteries}) are quoted from the cited
external sources with their published uncertainties. They are post-hoc
comparators, stored with solver access disabled
(\cref{prop:target-leakage}); none enters any computation whose output
it is meant to test.

\section{Reference and DOI Check}
\label{sec:doi-check}

\noindent The bibliography of this draft is the verified bibliography of
UIDT Ontology v3.9.9, carried over without addition. Two corrections
were made in the present audit pass and are marked in place:

\begin{enumerate}[leftmargin=1.6em,itemsep=0.25em]
\item \textbf{Attribution correction.} The survey arXiv:1904.12392,
cited in the v3.9 bibliography as by van Suijlekom alone, is by
\emph{Chamseddine and van Suijlekom}. The first author has been
restored. The claim this source supports in
\cref{rem:multiplicity} --- that the Standard-Model algebra with each
summand appearing once carries a KO-dimension-6 real structure ---
rests, in this draft, on Chamseddine--Connes--Marcolli~\cite{ccm2007},
with the survey cited as context. \emph{No theorem number is asserted
for a source whose theorem statement was not read in the present pass.}
\item \textbf{Bibliographic completion.} The page range and place of
publication for \'Ca\'ci\'c~\cite{cacic2011} (Wiesbaden, pp.~9--68) were
verified and added. The source's own abstract confirms the framing used
in Part~I: the structure theory of finite real spectral triples due to
Krajewski and to Paschke--Sitarz is there generalised to arbitrary
KO-dimension and to the failure of orientability and Poincar\'e duality,
and that theory reduces the study of finite spectral triples largely to
\emph{multiplicity matrices}. This is the primary-source basis of the
bimodule-multiplicity reading of G1.
\end{enumerate}

\noindent No source is cited in this draft for more than it claims.
Preprints are carried at the class of their strongest independent
confirmation, not at the class of their abstract. Any bibliography entry
required by future additions to this draft must pass a DOI or arXiv
resolution check before insertion; entries that cannot be resolved are
marked \textsc{unverified} and may not carry evidential weight.

\section{Consolidated Claims Register}
\label{sec:consolidated-claims}

\noindent The three per-part registers (\cref{tab:part1-claims},
\cref{tab:part2-claims}, \cref{tab:claims}) are authoritative in their
respective domains. The table below is a navigational summary of the
evidence distribution and is not a substitute for them.

\begin{table}[htbp]
\centering\scriptsize
\begin{tabular}{@{}p{2.5cm}p{3.4cm}p{3.2cm}p{3.4cm}@{}}
\toprule
Class & Part I & Part II & Part III ledger \\
\midrule
\catmark{A} (proof/identity/audit fact) & K-02, G-01--G-04, L-01, L-02 & O-01, O-03, Q-01, E-01, P-01 & ONT-01--ONT-03, ONT-06 \\
\catmark{A-} (calibrated) & --- & --- & ONT-04 ($\gamma$) \\
\catmark{B} (benchmarked) & lattice verdicts, \cref{tab:lattice} & --- & ONT-05 ($\Delta$, D18), ONT-07 \\
\catmark{C} (cosmology cap) & --- & --- & ONT-09 \\
\catmark{D} (programme) & K-01, G-05, S-01, S-02, C-01, C-02 & T-01, T-02 & ONT-08, ONT-11 \\
\catmark{E} (interpretive/museum) & quarantined audit doc & O-02, O-04 & ONT-10, L6, L8--L11 \\
\bottomrule
\end{tabular}
\caption{Evidence distribution across the three parts. The firmest
entries are boundary and negative results; every positive structural
ambition sits at \catmark{D}.}
\label{tab:consolidated}
\end{table}

\noindent Read as a whole, the distribution is the point. Not one
positive claim about the origin of gauge structure rises above
\catmark{D} in any part of this draft. The classes above \catmark{D} are
occupied by identities, obstructions, calibrations, audit facts, and
governance commitments. A framework whose strongest statements are the
statements of its own limits has not thereby proved anything about
nature; it has, at least, made itself checkable.

\section{Ledger Draft: PI Decision D19}
\label{sec:d19-ledger-draft}

\begin{pioverridebox}[D19 --- Ledger Entry, DRAFT for PI Signature]
\textbf{Decision.} The continuous real scalar field $\Sx$ is demoted
from ontological primitive (Variant~B, canonical in v3.9.x) to a
macroscopic effective-field-theory limit. A finite algebraic
pregeometric kernel is recorded as the candidate primitive under
investigation, at evidence class \catmark{D}.

\textbf{Binding conditions.}
(i)~Part~I is a formal research programme\catmark{D}, not a theorem
corpus; no Lean result, once machine-checked, upgrades a physics claim
without governance.
(ii)~No Standard-Model-like block partition is hard-coded into any core
formal structure; partitions enter as inserted candidate instances only,
under GLBC.
(iii)~All v3.9 ledger evidence classes are carried unchanged:
$\kappa = \KappaRG$\,\catmark{A}, $\lambda_S = \lambdaS$\,\catmark{A},
$v = \vVEV$\,\catmark{A}, $\gamma = \GammaGeom$\,\catmark{A-},
$\Delta = \DeltaGap$\,\catmark{B} (D18 unchanged), cosmology
$\leq$\,\catmark{C}. D19 reclassifies no constant.

\textbf{New ledger rows.} ONT-11 (\catmark{D}, Stratum~III);
tension entry T8 ($[2,2,2]$ rejection unverified, \catmark{D}).

\textbf{Scope.} D19 is an ontological reassignment, not a physics
result. It does not assert that the kernel is established. It removes a
stipulation; it adds no evidence.

\textbf{Status.} DRAFT. Takes effect only upon PI signature in
\texttt{LEDGER/CLAIMS.json}. Until signature, UIDT Ontology v3.9.9
(\texttt{-006}) remains the canonical manuscript of record and Variant~B
remains the canonical grammar. This document authorizes nothing; AI
audit output is advisory only (AI\_AUDIT\_POLICY Sec.~1, Gate~F8).

\medskip
\noindent\textit{Principal Investigator signature:}\\[1.4em]
\rule{6cm}{0.4pt}\qquad\textit{Date:}\ \rule{2.5cm}{0.4pt}
\end{pioverridebox}

\section{Ledger Record: PI Decision D20}
\label{sec:d20-ledger-record}

\begin{pioverridebox}[D20 --- Ledger Record: Architectural Reassignment]
\textbf{Decision.} The quantitative content of the former scalar/EFT Part~III
loses its role as active ontology and is reorganized as a verification,
benchmark, and parameter-ledger stratum. The historical scalar model is
retained in Appendix~\ref{app:scalar-museum} and is not imported into the
active architecture. The finite algebraic pregeometric kernel remains a
candidate under investigation at evidence class D.

\textbf{Binding conditions.}
(i)~Part~I remains a formal research programme at evidence class D, not a
theorem corpus; no machine-checked statement upgrades a physics claim without
separate governance.
(ii)~No Standard-Model-like block partition is hard-coded into a core formal
structure; partitions enter as inserted candidate instances under GLBC.
(iii)~All incoming evidence classes remain unchanged exactly as recorded in
the preceding D19 ledger entry; D20 changes architectural role only.
(iv)~PI-DELTA-P6, PI-ET-P6, PI-MUSEUM, and P2 remain separate decisions.

\textbf{Application boundary.} This executor candidate is subject to a
pre-application audit against blob \texttt{601cc57d}. Claude Desktop may
recommend GO or HOLD. Explicit PI GO is required before the single-writer
application. Final integration and merge remain PI-only.
\end{pioverridebox}

\section*{Open Obligations Carried Forward}
\addcontentsline{toc}{section}{Open Obligations Carried Forward}

\noindent The following are recorded so that no reader of this draft
mistakes its completeness for closure.

\begin{enumerate}[leftmargin=1.6em,itemsep=0.25em]
\item Lean findings F1 (register--code contradiction on the selection
proposition), F2 (target-shaped admissibility wrapper), and F3
(under-registered proof debt: 51 open obligations at head, 35
unregistered) are open. Discharging F1 upgrades the \emph{formal} layer
only; the proposition's physical grade remains bounded by H1 and H2.
\item Hypothesis H1 (intersection-form filter) is design-level: its
derivation from spectral-triple axioms is the stated upgrade path.
Hypothesis H2 (mass non-degeneracy) is a heuristic proxy for a claim
about the Yukawa sector. The $[2,2,2]$ rejection rests on H2 alone
(tension T8).
\item The Relocation Necessity Conjecture (\cref{oq:rnc}) and the
staircase conjecture (\cref{oq:staircase}) are open and independently
publishable.
\item The thermodynamic passage of Part~II has no construction. The
three no-go points --- absence of a dimensional-emergence theorem,
UV/IR mixing, and thermal-time circularity --- stand.
\item The Born rule, the $E_T$ scale bridge, the $H_0$ and $S_8$
tensions, and the origin of $G_{\mathrm{SM}}$ remain exactly where the
tension ledger (\cref{tab:tensions}) places them: unresolved, displayed,
and available for falsification.
\end{enumerate}


\clearpage
\appendix
\part*{Historical Appendix: Non-Imported Scalar/EFT Museum Record}
\addcontentsline{toc}{part}{Historical Appendix: Non-Imported Scalar/EFT Museum Record}

\begin{governancebox}[Museum Route B --- Reversible Retention under D20]
The material in this appendix is preserved to prevent silent loss, support
regression audit, and document the superseded v3.9/D19 scalar and EFT grammar.
It is not imported into Parts~I--III. Occurrences of
$\LegacySx$, $\partial_\mu S$, $\eta_{\mu\nu}$, the effective metric ansatz,
the scalar potential, or scalar correlators below are historical model content,
not active primitives, operators, acceptance conditions, or derivations. D20
does not authorize deletion; a future PI-MUSEUM decision may retain, relocate,
or remove this record through a separately audited patch.
\end{governancebox}

\section{Historical Effective-Scalar Axiomatics}
\label{sec:axiomatics}
\label{app:scalar-museum}

\noindent The v3.9 effective-sector record that associated the spectral gap $\Delta$, the
calibrated invariant $\gamma$, the coupling $\kappa$, and the
self-coupling $\lambda_S$ owes the reader the Lagrangian in which those
quantities actually live. This section preserves the four historical axioms, the effective
Lagrangian, and the dimensional bookkeeping for audit. None of these equations
is imported as a premise of the active Phase-6 architecture.
One axiom --- and only one --- changes under D19; the change is stated
first and in full.

\subsection{The Four Axioms in Demoted Form}
\label{subsec:four-axioms}

\begin{axiom}[Historical Effective-Field Postulate (D19 record)]
\label{ax:primacy}
There exists a real scalar effective field $\LegacySx$, the vacuum information
density, of canonical mass dimension $[\LegacySx]=1$. Under PI Decision D19
(draft), $\LegacySx$ is \emph{not} an ontological primitive: it is postulated
as the macroscopic, order-parameter-level description in which the
kernel programme's continuum limit (\cref{sec:thermo-passage}) is
expressed, and within the grammar of this Part it is the sole dynamical
scalar degree of freedom of the effective sector. The v3.9 reading of
this axiom --- $\LegacySx$ as sole fundamental degree of freedom, with all
structure emergent from its dynamics (Variant~B) --- is recorded in
\cref{subsec:variant-b} as the superseded grammar. Its preservation in this
appendix is a document-state fact.\catmark{A} The former ontological status
remains a historical interpretive commitment at \catmark{E}.
\end{axiom}

\begin{axiom}[Gauge-Invariant Self-Interaction]
\label{ax:gauge}
The dynamics of $\LegacySx$ are governed by a Lagrangian density
$\mathcal{L}_{\UIDT}$ invariant under $\mathrm{SU}(3)$ gauge
transformations and respecting vacuum stability ($V''(v)>0$). It is an
effective field theory valid below a cutoff $\Lambda_{\UIDT}$. The
non-minimal coupling
$(\bar\kappa/\Lambda_{\UIDT})\,S\,F_{\mu\nu}^a F^{a\mu\nu}$ has mass
dimension~5 and is therefore not renormalisable by strict power
counting; its validity is confined to the EFT domain (Limitation~L7). A
full non-perturbative UV completion is an open research
question.\catmark{A}
\end{axiom}

\noindent The minimal realisation satisfying \cref{ax:gauge} is the
effective Lagrangian
\begin{equation}
 \mathcal{L}_{\UIDT}
 = -\tfrac{1}{4}F_{\mu\nu}^a F^{a\mu\nu}
 + \tfrac{1}{2}(\partial_\mu S)(\partial^\mu S)
 - \tfrac{1}{2}\mu^2 S^2
 - \frac{\lambda_S}{4!}S^4
 + \frac{\bar\kappa}{\Lambda_{\UIDT}}\,S\,F_{\mu\nu}^a F^{a\mu\nu},
 \label{eq:lagrangian}
\end{equation}
where $\bar\kappa$ is the dimensionless EFT-normalised non-minimal
gauge--scalar coupling and $\lambda_S$ the scalar self-coupling; the
dimension-5 operator carries the dimensionful prefactor
$\bar\kappa/\Lambda_{\UIDT}$, in agreement with
\cref{tab:eft-dimensions}. Because $S$ is a real gauge singlet, its
kinetic term uses the ordinary derivative $\partial_\mu$. In the
framework's ledger and renormalisation-group statements this
dimensionless coefficient is written simply $\kappa$ (so
$\kappa\equiv\bar\kappa$, with ledger value $\kappa=\kappaVal$). The
exact internal coupling constraint of the framework is
$5\kappa^2 = 3\lambda_S$ (a candidate fixed-point condition, untested by
any defined functional renormalisation-group flow), which for
$\kappa = \tfrac{1}{2}$ gives the exact value $\lambda_S = \lambdaS$ and
the vacuum expectation value $v = \vVEV$; like $\lambda_S$, $v$ is fixed
by this exact constraint rather than independently measured, which is
why it carries no uncertainty. This constraint closes to a residual
below $10^{-14}$ at eighty-digit working precision
(\cref{prop:target-leakage} discipline applies: the closure is an
internal-consistency certificate, not an external agreement).

\begin{table}[H]
\scriptsize
\caption{Dimensional and EFT status of the terms entering
\cref{ax:gauge}. The dimension-5 operator marks the sector as an
effective field theory.}
\label{tab:eft-dimensions}
\renewcommand{\arraystretch}{1.25}
\begin{tabular*}{\textwidth}{@{}p{0.31\textwidth}@{\extracolsep{\fill}}ccp{0.34\textwidth}@{}}
\toprule
\textbf{Term} & \textbf{Op.\ dim.} & \textbf{Coupling dim.} & \textbf{Renormalisability} \\
\midrule
$F_{\mu\nu}^a F^{a\mu\nu}$ & $4$ & dimensionless & Strictly renormalisable Yang--Mills kinetic term \catmark{A} \\
$(\partial_\mu S)(\partial^\mu S)$ & $4$ & dimensionless & Canonical scalar kinetic term \catmark{A} \\
$\mu^2 S^2$ & $2$ & $+2$ & Relevant scalar mass operator \catmark{A} \\
$\lambda_S S^4$ & $4$ & dimensionless & Power-counting renormalisable self-interaction \catmark{A} \\
$(\bar\kappa/\Lambda_{\UIDT})\,S F_{\mu\nu}^a F^{a\mu\nu}$ & $5$ & $-1$ & Non-renormalisable; EFT-only below $\Lambda_{\UIDT}$ \catmark{A} \\
\bottomrule
\end{tabular*}
\end{table}

\begin{axiom}[Geometric Emergence]
\label{ax:emergence}
The effective spacetime metric $g_{\mu\nu}^{\mathrm{eff}}$ is not
fundamental but a derived quantity determined by the vacuum
configuration of $\LegacySx$. A candidate parametrisation is
\begin{equation}
 g_{\mu\nu}^{\mathrm{eff}} = \eta_{\mu\nu}
 + \alpha\,\partial_\mu\LegacySx\,\partial_\nu\LegacySx
 + \beta\,\LegacySx^2\,\eta_{\mu\nu},
 \label{eq:metric-emergence}
\end{equation}
with $\eta_{\mu\nu}$ the Minkowski metric and $\alpha,\beta$
dimensionful constants ($[\alpha]=-4$, $[\beta]=-2$) fixed by the
requirement that the emergent geometry reproduce general relativity in
the low-energy limit. The specific form of \cref{eq:metric-emergence}
and the values of $\alpha,\beta$ are model-dependent and carried at
predictive status; under D19 this axiom describes emergence \emph{within
the effective stratum}, with the deeper origin of $\eta_{\mu\nu}$
deferred to the Part-II passage and Limitation~L13.\catmark{D}
\end{axiom}

\begin{axiom}[Falsifiability]
\label{ax:falsifiability}
The theory defines a finite set of experimentally testable predictions
(\cref{sec:falsification-tension}), each with an explicit falsification
threshold. If any critical prediction is excluded at the specified
confidence level, the corresponding \UIDT\ mapping, calibration, or
prediction must be revised or abandoned.\catmark{A}
\end{axiom}

\begin{remark}[Falsifiability and critical rationalism]
\label{rem:popper}
Including falsifiability as a formal axiom is a methodological
commitment to the critical-rationalist tradition: a theory that cannot
in principle be wrong has no claim to being right.
\cref{ax:falsifiability} operationalises this by assigning every central
prediction a quantitative falsification threshold. The axiom is not a
physical law but a meta-theoretical commitment to epistemic honesty,
continuous with the principle that internal consistency is not external
confirmation (\cref{subsec:closure-vs-agreement}).
\end{remark}

\begin{evidencebox}[Historical Effective-Scalar Record: Preserved Status]{catA}
\textbf{Axioms.} Preservation of the Effective-Field Postulate is a
document-state fact \catmark{A}; its historical ontological content is
\catmark{E} (D19 form). Gauge-Invariant Self-Interaction \catmark{A} (EFT, dimension-5
operator); Geometric Emergence \catmark{D} (model-dependent metric);
Falsifiability \catmark{A} (meta-theoretical).\\[0.3em]
\textbf{Lagrangian.} \cref{eq:lagrangian} with $\kappa=\kappaVal$,
$\lambda_S=\lambdaS$, $v=\vVEV$; internal coupling constraint
$5\kappa^2=3\lambda_S$ closes internally to $<10^{-14}$ at eighty
digits.\\[0.3em]
\textbf{Status.} The Lagrangian is an effective field theory below
$\Lambda_{\UIDT}$; renormalisability is explicit per
\cref{tab:eft-dimensions}; the gauge sector is coupled, not generated
(\cref{subsec:coarse-geometry-eft}).
\end{evidencebox}

\section{Unmarked State, First Distinction, and Superseded Variant B (Historical Record)}
\label{sec:unmarked-distinction-sx}

\subsection{Why \texorpdfstring{$S(x)=0$}{S(x)=0} Is Not the Unmarked State}
\label{subsec:sx-zero-not-unmarked}

A foundational theory must be exact about its own beginning, and the
most common error at the beginning is to confuse \emph{absence of value}
with \emph{absence of structure}. \UIDT\ takes the unmarked state,
``nothing'' in the precise sense of Spencer-Brown's calculus of
distinctions, as the condition in which no field-value assignment exists
at all. It is essential to see that this is \emph{not} the same as the
field configuration $S(x) = 0$.\catmark{A}

The point is structural, not rhetorical. To write $S(x) = 0$ is already
to presuppose a great deal: that there is a field $S$; that there is a
domain indexed by $x$ over which it is defined; that evaluation of the
field at a point is a meaningful operation; and, most importantly, that
the value $0$ is \emph{distinguishable} from other possible values. A
vanishing field value is a marked state: it is the assertion ``the field
here is zero,'' and an assertion of that form lives entirely within an
already-constituted system of field, domain, and distinguishability. The
unmarked state precedes all of this. It is not the number zero; it is
the absence of the apparatus within which numbers could be
assigned.\catmark{A}

Because the unmarked state is not a physical configuration, it cannot be
assigned a physical history, a temperature, a fluctuation spectrum, or a
cosmological epoch. One cannot say that ``in the beginning $S(x)=0$ and
then it became non-zero,'' because that sentence smuggles a temporal
ordering and a pre-existing field into a condition that, by
construction, contains neither. The framework explicitly refuses
cosmological, theological, and metaphysical elaborations of the unmarked
state. It is a logical boundary, not a primordial substance, and the
prohibition on treating it as a substance is as binding as any numerical
constraint in the theory.\catmark{A} The result is unaffected by D19:
it concerned the grammar of beginnings, not the identity of the
primitive, and it survives the primitive's replacement intact.

\subsection{Variant~B, Recorded and Superseded}
\label{subsec:variant-b}

The relationship between the first distinction and the scalar field
$S(x)$ admits two readings, and the developmental record was right to
treat the choice between them as a potential source of circularity. Under
Variant~A, the distinction is strictly prior: first-distinction, then
$S(x)$, then relations, then geometry. Under Variant~B, the field
$S(x)$ is itself the primitive, and the first distinction is identified
with its first non-trivial realisation, $S(x)\neq 0$. The v3.9
manuscript adopted Variant~B as a matter of \emph{internal ontological
grammar} --- explicitly not as a physical proof that the field is the
unique primitive of nature. That adoption is a document-state fact.\catmark{A}
The choice itself was an interpretive commitment at \catmark{E}. Any inference
of unique physical primitiveness from that grammar remains at \catmark{E}.

Under D19 that grammar is superseded. The office of candidate primitive
transfers to the pregeometric kernel of Part~I, held at \catmark{D}; the
first distinction is no longer identified with $S(x)\neq 0$, and the
question of what, in a finite algebraic primitive, plays the role of
first non-trivial structure is recorded as open rather than answered by
relabelling. Two elements of the Variant-B analysis are retained at full
force because they were never specific to the scalar. The first is the
anti-circularity discipline itself: naming the primitive choice, and
labelling it as a grammar rather than a proof, is what keeps the
definitional circle from closing unnoticed, and D19 is written in
exactly that register. The second is the terminology guard.

\begin{governancebox}[Terminology Guard: $S(x)$ and ``Information'']
The \UIDT\ information density $S(x)$ is not Shannon, semantic, Bateson,
algorithmic, thermodynamic, or epistemic information. It is a real
scalar effective field (D19). ``Information'' in \UIDT\ is the emergent
counting-measure of difference, attaching at the first distinction.
Under no reading is $S(x)$ identified with consciousness; that
identification is a category error and is prohibited --- and the same
prohibition applies, unchanged, to the kernel candidate of
Part~I.\catmark{A}
\end{governancebox}

\subsection{The Corrected Ontological Sequence: Morphism as Formal Tool, Not Generator}
\label{subsec:corrected-sequence}

The developmental dialogue at one stage placed ``morphism'' within the
vertical sequence of ontological stages, as though a morphism were a
step in the genesis of being. The honesty audit corrected this, and the
correction survives D19 unchanged because it is a statement about
mathematics, not about the primitive. The corrected architecture rests
on a separation between two axes: a vertical \emph{ontological} axis,
the genesis of reality, and a horizontal \emph{formal} axis, the
mathematical apparatus used to describe that genesis.\catmark{A}

The decisive observation is elementary and therefore robust. A morphism
$f\colon A\to B$ presupposes a distinguishable domain $A$ and a
distinguishable codomain $B$; the very act of writing down a morphism
requires that two distinct objects already exist to be related. A
morphism therefore cannot be the generator of distinction, because it
presupposes distinction.\catmark{A} It belongs to the formal
representation layer, not to the primitive genesis layer. Information
docks at the first distinction, as the emergent counting-measure of
difference; morphisms dock at relations, as the computational interface.
Their placement outside the formal representation layer is a structural
bookkeeping statement.\catmark{A} The claim that neither is a constituent
of genesis remains interpretive at \catmark{E}. Under D19 the
vertical line is re-read with the kernel candidate at its base and the
effective field far downstream; the placement of the formal apparatus on
the horizontal axis is unchanged, and the major categorical and
relational programmes --- topos theory, causal sets, constructor theory,
QBism --- are catalogued here as external lenses.\catmark{A} Their exclusion
from the definition of the primitive remains an interpretive boundary at
\catmark{E}.

\subsection{Relations, Graph Discipline, and the Effective-Field Boundary}
\label{subsec:coarse-geometry-eft}

The graph stage of the vertical line is a \emph{theorem target}, not a
completed construction: the honest form of the proposal is not ``the
world is a graph'' but ``reconstructing an effective spacetime from a
correlation network is a legitimate theorem target.''\catmark{D} No
quantitative graph claim --- no count of edges, no degree distribution,
no spectral property --- is admissible until nodes, edges, weights, and
dynamics are explicitly defined; graph notation represents relational
structure and is not that structure; and the schematic $I \to R(I)$ is
toy notation, a recorded direction of construction, not a defined
map.\catmark{E} The candidate primary observable
$G(x,y) = \langle S(x)\,S(y)\rangle$ is recorded as a candidate whose
canonical status must be confirmed against the ledger before any
load-bearing use.\catmark{D} One prohibition carries Class~A force and
is inherited by Part~II's boundary: there is no direct inference from
phenomenology to graph structure.\catmark{A}

Geometry, in this architecture, is a coarse-grained representation of
underlying structure --- the pregeometric tradition's position, with the
thermodynamic derivation of the Einstein equations from horizon entropy
(Jacobson)~\cite{jacobson1995} and the later entropic-gravity programme~\cite{padmanabhan2010,verlinde2011,chirco2014} invoked as
Stratum-II precedent that emergent geometry is a respectable position,
not as a derivation imported into \UIDT.\catmark{E} Downstream of
geometry lies the effective field regime, whose most important sentence
must be stated plainly, because it is the honest boundary this Part
inherits from v3.9 and now shares with Part~I: the framework couples the
scalar $S(x)$ to a Yang--Mills sector that is \emph{already present}.
The gauge fields are part of the effective description, not something
the scalar has been shown to generate.\catmark{A}

\begin{evidencebox}[Historical Effective-Field Boundary]{catA}
\textbf{Regime.} The gauge sector is an effective-field-theory regime,
in which $S(x)$ is coupled to an already-present Yang--Mills sector. It
is not a generative regime.\catmark{A}\\[0.3em]
\textbf{Negative result.} \UIDT\ does \emph{not} derive the
Standard-Model gauge group $G_{\mathrm{SM}}$ from a structureless real
scalar $S(x)\in\mathbb{R}$ --- and, per Part~I, no known route derives
it from the kernel candidate either; on every known route it is
selected, not forced.\catmark{A}\\[0.3em]
\textbf{Consequence.} The transition from primitive to
$G_{\mathrm{SM}}$ is a research programme, not a result; its
localization (G1--G4), lower bounds, and candidate algebras are the
subject of Part~I.\catmark{D}
\end{evidencebox}

\newpage

\begin{thebibliography}{9}
\footnotesize

\bibitem{uidt-ontology}
P.~Rietz, \textit{Vacuum Information Density as the Fundamental Geometric Scalar: UIDT Ontology v3.9} (this work). Zenodo. DOI: \href{https://doi.org/10.5281/zenodo.20319634}{10.5281/zenodo.20319634}.

\bibitem{uidt-framework}
P.~Rietz, \textit{UIDT Framework v3.9 (parent record)}. Zenodo. DOI: \href{https://doi.org/10.5281/zenodo.17835200}{10.5281/zenodo.17835200}.

\bibitem{uidt-clay}
P.~Rietz, \textit{UIDT Clay-Submission Record}. Zenodo. DOI: \href{https://doi.org/10.5281/zenodo.18003018}{10.5281/zenodo.18003018}. Historical provenance record; cited as documentation only, not as an active claim.

\bibitem{bib-audit-note}
\textbf{Bibliography audit note.} The external references below were resolved to verifiable primary sources (DOI or arXiv) in a dedicated literature audit. In keeping with the framework's evidence discipline, citation of these sources changes no \UIDT\ evidence class: external literature supports the quantitative pure-gauge benchmarks, is contested on the interpretive physics, explicitly identifies open chiral-fermion obstructions, is interpretive for the consciousness frameworks, and refutes any MeV-to-neural coupling. A small number of secondary venue/page details remain flagged for final primary confirmation before submission.

% Cluster 1: pure SU(3) lattice benchmarks
\bibitem{athenodorou2020}
A.~Athenodorou, M.~Teper, ``The glueball spectrum of SU(3) gauge theory in 3+1 dimensions,'' JHEP \textbf{11} (2020) 172. arXiv:2007.06422.
\bibitem{athenodorou2021}
A.~Athenodorou, M.~Teper, ``SU(N) gauge theories in 3+1 dimensions: glueball spectrum, string tensions and topology,'' JHEP \textbf{12} (2021) 082. arXiv:2106.00364.
\bibitem{morningstar1999}
C.~J.~Morningstar, M.~Peardon, ``The glueball spectrum from an anisotropic lattice study,'' Phys.\ Rev.\ D \textbf{60} (1999) 034509. arXiv:hep-lat/9901004.
\bibitem{morningstar2024}
C.~Morningstar, ``Update on Glueballs,'' PoS \textbf{LATTICE2024} (2024) 004. DOI 10.22323/1.466.0004. arXiv:2502.02547.
\bibitem{michael1989}
C.~Michael, M.~Teper, ``The glueball spectrum in SU(3),'' Nucl.\ Phys.\ B \textbf{314} (1989) 347.
\bibitem{luscher2010}
M.~L\"uscher, ``Properties and uses of the Wilson flow in lattice QCD,'' JHEP \textbf{08} (2010) 071. arXiv:1006.4518.
\bibitem{borsanyi2012}
S.~Bors\'anyi et al.\ (BMW), ``High-precision scale setting in lattice QCD,'' JHEP \textbf{09} (2012) 010. arXiv:1203.4469.
\bibitem{chitop2025}
S.~D\"urr, G.~Fuwa, ``Topological susceptibility and excess kurtosis in SU(3) Yang--Mills theory,'' Phys.\ Rev.\ D \textbf{113} (2026) 054508. arXiv:2501.08217.
\bibitem{durr2007}
S.~D\"urr, Z.~Fodor, C.~Hoelbling, T.~Kurth, ``Precision study of the SU(3) topological susceptibility in the continuum,'' arXiv:hep-lat/0612021.
\bibitem{dallabrida2019}
M.~Dalla Brida, A.~Ramos, ``The gradient flow coupling at high-energy and the scale of SU(3) Yang--Mills theory,'' Eur.\ Phys.\ J.\ C \textbf{79} (2019) 720. arXiv:1905.05147.

% Cluster 1b: non-Abelian Casimir effect on the lattice
\bibitem{chernodub2023}
M.~N.~Chernodub, V.~A.~Goy, A.~V.~Molochkov, A.~S.~Tanashkin, ``Boundary states and non-Abelian Casimir effect in lattice Yang--Mills theory,'' Phys.\ Rev.\ D \textbf{108} (2023) 014515. arXiv:2302.00376.
\bibitem{ngwenya2025}
B.~A.~Ngwenya, A.~K.~Rothkopf, W.~A.~Horowitz, ``The Non-Abelian Casimir Effect for Plates, Symmetrical Tube and Box on the Lattice,'' arXiv:2507.21333 (2025); B.~A.~Ngwenya, ``The Casimir effect in non-Abelian gauge theories on the lattice,'' Ph.D.\ thesis, University of Cape Town (2025), arXiv:2505.02875.

% Cluster 2: noncommutative geometry SM
\bibitem{chamseddine1997}
A.~H.~Chamseddine, A.~Connes, ``The Spectral Action Principle,'' Commun.\ Math.\ Phys.\ \textbf{186} (1997) 731. arXiv:hep-th/9606001.
\bibitem{chamseddine1996}
A.~H.~Chamseddine, A.~Connes, ``Universal Formula for Noncommutative Geometry Actions,'' Phys.\ Rev.\ Lett.\ \textbf{77} (1996) 4868.
\bibitem{ccm2007}
A.~H.~Chamseddine, A.~Connes, M.~Marcolli, ``Gravity and the standard model with neutrino mixing,'' Adv.\ Theor.\ Math.\ Phys.\ \textbf{11} (2007) 991. arXiv:hep-th/0610241.
\bibitem{cc2012}
A.~H.~Chamseddine, A.~Connes, ``Resilience of the Spectral Standard Model,'' JHEP \textbf{09} (2012) 104.

\bibitem{cc2006}
A.~H.~Chamseddine, A.~Connes, ``Scale invariance in the spectral action,'' J.\ Math.\ Phys.\ \textbf{47} (2006) 063504. arXiv:hep-th/0512169.

\bibitem{ccvs2013}
A.~H.~Chamseddine, A.~Connes, W.~D.~van Suijlekom, ``Beyond the Spectral Standard Model: Emergence of Pati--Salam Unification,'' JHEP \textbf{11} (2013) 132. arXiv:1304.8050.

\bibitem{connes2013}
A.~Connes, ``On the spectral characterization of manifolds,'' J.\ Noncommut.\ Geom.\ \textbf{7} (2013) 1. arXiv:0810.2088.
\bibitem{fuzzy1012}
M.~Ihl, C.~Sachse, C.~S\"amann, ``Fuzzy Scalar Field Theory as Matrix Quantum Mechanics,'' JHEP \textbf{03} (2011) 091. arXiv:1012.3568. DOI 10.1007/JHEP03(2011)091. (Candidate source for the induced-matrix-structure mechanism; \catmark{D}.)

% Cluster 3: string-net / topological order
\bibitem{levinwen2005}
M.~A.~Levin, X.-G.~Wen, ``String-net condensation: A physical mechanism for topological phases,'' Phys.\ Rev.\ B \textbf{71} (2005) 045110. arXiv:cond-mat/0404617.

\bibitem{kitaev2003}
A.~Yu.~Kitaev, ``Fault-tolerant quantum computation by anyons,'' Ann.\ Phys.\ \textbf{303} (2003) 2. arXiv:quant-ph/9707021.

\bibitem{mueger2003}
M.~M\"uger, ``From subfactors to categories and topology II: The quantum double of tensor categories and subfactors,'' J.\ Pure Appl.\ Algebra \textbf{180} (2003) 159. arXiv:math/0111205.

\bibitem{eno2005}
P.~Etingof, D.~Nikshych, V.~Ostrik, ``On fusion categories,'' Ann.\ Math.\ \textbf{162} (2005) 581. arXiv:math/0203060.

\bibitem{lankongwen2018}
T.~Lan, L.~Kong, X.-G.~Wen, ``Classification of (3+1)D bosonic topological orders: The case when pointlike excitations are all bosons,'' Phys.\ Rev.\ X \textbf{8} (2018) 021074. arXiv:1704.04221.

\bibitem{lanwen2019}
T.~Lan, X.-G.~Wen, ``Classification of (3+1)D bosonic topological orders (II): The case when some pointlike excitations are fermions,'' Phys.\ Rev.\ X \textbf{9} (2019) 021005. arXiv:1801.08530.
\bibitem{wen2003}
X.-G.~Wen, ``Quantum order from string-net condensations and the origin of light and massless fermions,'' Phys.\ Rev.\ D \textbf{68} (2003) 065003. arXiv:hep-th/0302201. DOI 10.1103/PhysRevD.68.065003. [Article number recertified to 065003 against APS primary metadata; the previously-circulating 024501 is a cross-contamination artifact and is incorrect.]
\bibitem{walkerwang2012}
K.~Walker, Z.~Wang, ``(3+1)-TQFTs and topological insulators,'' Front.\ Phys.\ \textbf{7} (2012) 150. arXiv:1104.2632.
\bibitem{vonkeyserlingk2013}
C.~W.~von Keyserlingk, F.~J.~Burnell, S.~H.~Simon, ``Three-dimensional topological lattice models with surface anyons,'' Phys.\ Rev.\ B \textbf{87} (2013) 045107. arXiv:1208.5128.

% Cluster 4: quantum link / D-theory
\bibitem{brower1999}
R.~Brower, S.~Chandrasekharan, U.-J.~Wiese, ``QCD as a quantum link model,'' Phys.\ Rev.\ D \textbf{60} (1999) 094502. arXiv:hep-lat/9704106.
\bibitem{brower2004}
R.~Brower, S.~Chandrasekharan, S.~Riederer, U.-J.~Wiese, ``D-theory: field quantization by dimensional reduction of discrete variables,'' Nucl.\ Phys.\ B \textbf{693} (2004) 149. arXiv:hep-lat/0309182. DOI 10.1016/j.nuclphysb.2004.06.007.
\bibitem{beard1998}
B.~B.~Beard, R.~Brower, S.~Chandrasekharan, D.~Chen, A.~Tsapalis, U.-J.~Wiese, ``D-theory: field theory via dimensional reduction of discrete variables'' (proceedings), Nucl.\ Phys.\ B Proc.\ Suppl.\ \textbf{63} (1998) 775. arXiv:hep-lat/9709120.
\bibitem{wiese2022}
U.-J.~Wiese, ``From quantum link models to D-theory: a resource efficient framework,'' Phil.\ Trans.\ R.\ Soc.\ A \textbf{380} (2022) 20210068. arXiv:2107.09335.

% Cluster 5: gauged tensor networks / PEPS
\bibitem{tagliacozzo2014}
L.~Tagliacozzo, A.~Celi, M.~Lewenstein, ``Tensor Networks for Lattice Gauge Theories with Continuous Groups,'' Phys.\ Rev.\ X \textbf{4} (2014) 041024.
\bibitem{haegeman2015}
J.~Haegeman, K.~Van Acoleyen, N.~Schuch, J.~I.~Cirac, F.~Verstraete, ``Gauging quantum states,'' Phys.\ Rev.\ X \textbf{5} (2015) 011024.
\bibitem{cirac2021}
J.~I.~Cirac, D.~P\'erez-Garc\'ia, N.~Schuch, F.~Verstraete, ``Matrix product states and projected entangled pair states,'' Rev.\ Mod.\ Phys.\ \textbf{93} (2021) 045003.
\bibitem{nielsen1981}
H.~B.~Nielsen, M.~Ninomiya, ``Absence of neutrinos on a lattice,'' Nucl.\ Phys.\ B \textbf{185} (1981) 20; \textbf{193} (1981) 173.

\bibitem{weinbergwitten1980}
S.~Weinberg, E.~Witten, ``Limits on Massless Particles,'' Phys.\ Lett.\ B \textbf{96} (1980) 59. DOI 10.1016/0370-2693(80)90212-9.

% Cluster 6: thermodynamic gravity
\bibitem{jacobson1995}
T.~Jacobson, ``Thermodynamics of Spacetime: The Einstein Equation of State,'' Phys.\ Rev.\ Lett.\ \textbf{75} (1995) 1260. arXiv:gr-qc/9504004.
\bibitem{padmanabhan2010}
T.~Padmanabhan, ``Thermodynamical Aspects of Gravity: New insights,'' Rep.\ Prog.\ Phys.\ \textbf{73} (2010) 046901. arXiv:0911.5004.
\bibitem{verlinde2011}
E.~Verlinde, ``On the Origin of Gravity and the Laws of Newton,'' JHEP \textbf{04} (2011) 029. arXiv:1001.0785.
\bibitem{chirco2014}
G.~Chirco, H.~M.~Haggard, A.~Riello, C.~Rovelli, ``Spacetime thermodynamics without hidden degrees of freedom,'' Phys.\ Rev.\ D \textbf{90} (2014) 044044. arXiv:1401.5262.

% Cluster 6b: information-theoretic / structural-realist prior art and emergent-gravity constraints
\bibitem{wheeler1990}
J.~A.~Wheeler, ``Information, Physics, Quantum: The Search for Links,'' in W.~H.~Zurek (ed.), \emph{Complexity, Entropy, and the Physics of Information}, Addison-Wesley (1990), pp.~3--28.
\bibitem{carrollsingh2018}
S.~M.~Carroll, A.~Singh, ``Mad-Dog Everettianism: Quantum Mechanics at Its Most Minimal,'' arXiv:1801.08132 (2018); in \emph{What is Fundamental?}, The Frontiers Collection, Springer (2019), DOI 10.1007/978-3-030-11301-8\_10.
\bibitem{jenkins2009}
A.~Jenkins, ``Constraints on emergent gravity,'' Int.\ J.\ Mod.\ Phys.\ D \textbf{18} (2009) 2249. arXiv:0904.0453.
\bibitem{carlip2014}
S.~Carlip, ``Challenges for Emergent Gravity,'' Stud.\ Hist.\ Phil.\ Sci.\ B \textbf{46} (2014) 200. arXiv:1207.2504.

% Cluster 7: observer / neuroscience comparison (interpretive)
\bibitem{friston2010}
K.~Friston, ``The free-energy principle: a unified brain theory?,'' Nat.\ Rev.\ Neurosci.\ \textbf{11} (2010) 127.
\bibitem{metzinger2003}
T.~Metzinger, \textit{Being No One: The Self-Model Theory of Subjectivity}, MIT Press (2003).
\bibitem{hoffman2014}
D.~Hoffman, C.~Prakash, ``Objects of Consciousness,'' Front.\ Psychol.\ \textbf{5} (2014) 577.
\bibitem{hameroff2014}
S.~Hameroff, R.~Penrose, ``Consciousness in the universe: A review of the `Orch OR' theory,'' Phys.\ Life Rev.\ \textbf{11} (2014) 39.
\bibitem{tononi2004}
G.~Tononi, ``An information integration theory of consciousness,'' BMC Neurosci.\ \textbf{5} (2004) 42.

% Cluster 8: scale-bridge / decoupling
\bibitem{attwell2001}
D.~Attwell, S.~B.~Laughlin, ``An energy budget for signaling in the grey matter of the brain,'' J.\ Cereb.\ Blood Flow Metab.\ \textbf{21} (2001) 1133.
\bibitem{harris2012}
J.~J.~Harris, R.~Jolivet, D.~Attwell, ``Synaptic energy use and supply,'' Neuron \textbf{75} (2012) 762.
\bibitem{levy2021}
W.~B.~Levy, V.~G.~Calvert, ``Communication consumes 35 times more energy than computation in the human cortex,'' PNAS \textbf{118} (2021) e2008173118.
\bibitem{appelquist1975}
T.~Appelquist, J.~Carazzone, ``Infrared singularities and massive fields,'' Phys.\ Rev.\ D \textbf{11} (1975) 2856.

% v4.x G_SM-gap research frontiers (literature-status, [D]/[E])
\bibitem{cc2008}
A.~H.~Chamseddine, A.~Connes, ``Why the Standard Model,'' J.\ Geom.\ Phys.\ \textbf{58} (2008) 38. arXiv:0706.3688. DOI 10.1016/j.geomphys.2007.09.011.
\bibitem{ccm2015}
A.~H.~Chamseddine, A.~Connes, V.~Mukhanov, ``Quanta of Geometry: Noncommutative Aspects,'' Phys.\ Rev.\ Lett.\ \textbf{114} (2015) 091302. arXiv:1409.2471.
\bibitem{grosse2005}
H.~Grosse, H.~Steinacker, ``Finite Gauge Theory on Fuzzy $CP^2$,'' Nucl.\ Phys.\ B \textbf{707} (2005) 145. arXiv:hep-th/0407089. DOI 10.1016/j.nuclphysb.2004.11.058.
\bibitem{azuma2004}
T.~Azuma, S.~Bal, K.~Nagao, J.~Nishimura, ``Nonperturbative studies of fuzzy spheres in a matrix model with the Chern--Simons term,'' JHEP \textbf{05} (2004) 005. arXiv:hep-th/0401038. DOI 10.1088/1126-6708/2004/05/005.
\bibitem{farnsworth2015}
S.~Farnsworth, L.~Boyle, ``Rethinking Connes' approach to the standard model of particle physics via non-commutative geometry,'' New J.\ Phys.\ \textbf{17} (2015) 023021. arXiv:1408.5367. DOI 10.1088/1367-2630/17/2/023021.
\bibitem{aastrup2025}
J.~Aastrup, J.~M.~Grimstrup, ``On the emergence of an almost-commutative spectral triple from a geometric construction on a configuration space,'' (2025). arXiv:2504.03391.
\bibitem{luscher1999}
M.~L\"uscher, ``Abelian chiral gauge theories on the lattice with exact gauge invariance,'' Nucl.\ Phys.\ B \textbf{549} (1999) 295. arXiv:hep-lat/9811032. DOI 10.1016/S0550-3213(99)00115-7.
\bibitem{golterman2023}
M.~Golterman, Y.~Shamir, ``Propagator zeros and lattice chiral gauge theories,'' (2023). arXiv:2311.12790; and ``Constraints on the symmetric mass generation paradigm,'' (2025). arXiv:2505.20436.
\bibitem{dupuis2021}
N.~Dupuis et al., ``The nonperturbative functional renormalization group and its applications,'' Phys.\ Rept.\ \textbf{910} (2021) 1. arXiv:2006.04853. DOI 10.1016/j.physrep.2021.01.001.
\bibitem{schnoerr2013}
D.~Schnoerr, I.~Boettcher, J.~M.~Pawlowski, C.~Wetterich, ``Error estimates and specification parameters for functional renormalization,'' (2013). arXiv:1301.4169.

% Appendix A endogenous-enhancement cluster (added in 004-appA-gsm; entries
% resolved to primary sources; flagged for the regular DOI sweep)
\bibitem{wilson1974}
K.~G.~Wilson, ``Confinement of quarks,'' Phys.\ Rev.\ D \textbf{10} (1974) 2445. DOI 10.1103/PhysRevD.10.2445.
\bibitem{frenkelkac1980}
I.~B.~Frenkel, V.~G.~Kac, ``Basic representations of affine Lie algebras and dual resonance models,'' Invent.\ Math.\ \textbf{62} (1980) 23.
\bibitem{goddardolive1986}
P.~Goddard, D.~Olive, ``Kac--Moody and Virasoro algebras in relation to quantum physics,'' Int.\ J.\ Mod.\ Phys.\ A \textbf{1} (1986) 303.
\bibitem{kuranishi1951}
M.~Kuranishi, ``On everywhere dense imbedding of free groups in Lie groups,'' Nagoya Math.\ J.\ \textbf{2} (1951) 63.
\bibitem{madore1992}
J.~Madore, ``The fuzzy sphere,'' Class.\ Quantum Grav.\ \textbf{9} (1992) 69.
\bibitem{ikkt1997}
N.~Ishibashi, H.~Kawai, Y.~Kitazawa, A.~Tsuchiya, ``A large-$N$ reduced model as superstring,'' Nucl.\ Phys.\ B \textbf{498} (1997) 467. arXiv:hep-th/9612115.
\bibitem{steinacker2010}
H.~Steinacker, ``Emergent geometry and gravity from matrix models: an introduction,'' Class.\ Quantum Grav.\ \textbf{27} (2010) 133001. arXiv:1003.4134.
\bibitem{krajewski1998}
T.~Krajewski, ``Classification of finite spectral triples,'' J.\ Geom.\ Phys.\ \textbf{28} (1998) 1. arXiv:hep-th/9701081.
\bibitem{paschkesitarz1998}
M.~Paschke, A.~Sitarz, ``Discrete spectral triples and their symmetries,'' J.\ Math.\ Phys.\ \textbf{39} (1998) 6191. arXiv:q-alg/9612029.

% Appendix B Connes-vacuum attractor cluster (added in 005-appB-connes;
% entries resolved to primary sources; flagged for the regular DOI sweep)
\bibitem{agbm1995}
E.~Alvarez, J.~M.~Gracia-Bond\'ia, C.~P.~Mart\'in, ``Anomaly cancellation and the gauge group of the standard model in non-commutative geometry,'' Phys.\ Lett.\ B \textbf{364} (1995) 33. arXiv:hep-th/9506115.
\bibitem{connes2006}
A.~Connes, ``Noncommutative geometry and the standard model with neutrino mixing,'' JHEP \textbf{11} (2006) 081. arXiv:hep-th/0608226.
\bibitem{barrett2007}
J.~W.~Barrett, ``A Lorentzian version of the non-commutative geometry of the standard model of particle physics,'' J.\ Math.\ Phys.\ \textbf{48} (2007) 012303. arXiv:hep-th/0608221.
\bibitem{barrettglaser2016}
J.~W.~Barrett, L.~Glaser, ``Monte Carlo simulations of random non-commutative geometries,'' J.\ Phys.\ A \textbf{49} (2016) 245001. arXiv:1510.01377.
\bibitem{cotler2017}
J.~S.~Cotler et al., ``Black holes and random matrices,'' JHEP \textbf{05} (2017) 118. arXiv:1611.04650.

% Cosmology inputs (Stratum-I, calibration only)
\bibitem{planck2018}
Planck Collaboration, ``Planck 2018 results. VI. Cosmological parameters,'' Astron.\ Astrophys.\ \textbf{641} (2020) A6. arXiv:1807.06209.
\bibitem{riess2022}
A.~G.~Riess et al., ``A Comprehensive Measurement of the Local Value of the Hubble Constant\ldots (SH0ES),'' Astrophys.\ J.\ Lett.\ \textbf{934} (2022) L7. arXiv:2112.04510.
\bibitem{desidr2}
DESI Collaboration, ``DESI DR2 Results II: Measurements of Baryon Acoustic Oscillations and Cosmological Constraints,'' (2025). arXiv:2503.14738.
\bibitem{kids2025}
KiDS Collaboration (A.~H.~Wright et al.), ``KiDS-Legacy: Cosmological constraints from cosmic shear with the complete Kilo-Degree Survey,'' arXiv:2503.19441 (2025).
\bibitem{des2026}
DES Collaboration (T.~M.~C.~Abbott et al.), ``Dark Energy Survey Year 6 Results: Cosmological Constraints from Galaxy Clustering and Weak Lensing,'' arXiv:2601.14559 (2026).

\bibitem{vansuijlekom2019}
A.~H.~Chamseddine, W.~D.~van Suijlekom, ``A survey of spectral models of gravity coupled to matter,'' arXiv:1904.12392 (2019); in \emph{Advances in Noncommutative Geometry}, pp.~1--51, Springer (2019). \textsc{[v4.0 audit: first author restored; the v3.9 entry attributed this survey to van Suijlekom alone.]}
\bibitem{cacic2011}
B.~\'{C}a\'{c}i\'{c}, ``Moduli spaces of Dirac operators for finite spectral triples,'' in \emph{Quantum Groups and Noncommutative Spaces}, Aspects of Mathematics \textbf{41}, Vieweg+Teubner, Wiesbaden (2011), pp.~9--68, DOI 10.1007/978-3-8348-9831-9\_2; arXiv:0902.2068.
\bibitem{iss2004}
B.~Iochum, T.~Sch\"{u}cker, C.~Stephan, ``On a classification of irreducible almost commutative geometries,'' J.\ Math.\ Phys.\ \textbf{45} (2004) 5003. arXiv:hep-th/0312276.
\bibitem{furey2018}
C.~Furey, ``$\mathrm{SU}(3)_C\times\mathrm{SU}(2)_L\times\mathrm{U}(1)_Y$ as a symmetry of division algebraic ladder operators,'' Eur.\ Phys.\ J.\ C \textbf{78} (2018) 375. arXiv:1806.05209.
\bibitem{furey2025}
C.~Furey, ``A superalgebra within,'' Ann.\ Phys.\ (Berlin) (2025), DOI 10.1002/andp.202500229; arXiv:2505.07923.
\bibitem{gourlay2024}
B.~Gourlay, N.~Gresnigt, ``Three generations of standard model fermions from octonionic Clifford algebras,'' arXiv:2407.01580 (2024).
\bibitem{vanraamsdonk2010}
M.~Van Raamsdonk, ``Building up spacetime with quantum entanglement,'' Gen.\ Rel.\ Grav.\ \textbf{42} (2010) 2323--2329, DOI 10.1007/s10714-010-1034-0; also Int.\ J.\ Mod.\ Phys.\ \textbf{D19} (2010) 2429--2435. arXiv:1005.3035. \textsc{[Essay; scope is AdS/CFT. Cited for the in-holography statement only.]}
\bibitem{cubitt2015}
T.~S.~Cubitt, D.~P\'erez-Garc\'ia, M.~M.~Wolf, ``Undecidability of the spectral gap,'' Nature \textbf{528} (2015) 207--211, DOI 10.1038/nature16059; full version arXiv:1502.04573; letter arXiv:1502.04135. \textsc{[Scope: general local translation-invariant Hamiltonians on a 2D lattice. Asserts nothing about the Yang--Mills gap or about $\Delta$.]}
\bibitem{connesrovelli1994}
A.~Connes, C.~Rovelli, ``Von Neumann algebra automorphisms and time--thermodynamics relation in generally covariant quantum theories,'' Class.\ Quantum Grav.\ \textbf{11} (1994) 2899--2917, DOI 10.1088/0264-9381/11/12/007; arXiv:gr-qc/9406019.
\bibitem{parker2019}
D.~E.~Parker, X.~Cao, A.~Avdoshkin, T.~Scaffidi, E.~Altman, ``A Universal Operator Growth Hypothesis,'' Phys.\ Rev.\ X \textbf{9} (2019) 041017, DOI 10.1103/PhysRevX.9.041017; arXiv:1812.08657. \textsc{[Hypothesis, verified in specific models; thermodynamic-limit setting. Bounds $\lambda_L \le 2\alpha$; asserts nothing about state-space differentiability.]}
\bibitem{rabinovici2021}
E.~Rabinovici, A.~S\'anchez-Garrido, R.~Shir, J.~Sonner, ``Operator complexity: a journey to the edge of Krylov space,'' JHEP \textbf{06} (2021) 062, DOI 10.1007/JHEP06(2021)062; arXiv:2009.01862. \textsc{[Finite-entropy/finite-dimensional setting; rigorous bounds on K-complexity; SYK4 saturates, SYK2 stays exponentially below.]}
\end{thebibliography}


\end{document}
